\documentclass[11pt,letterpaper]{article}
\usepackage[T1]{fontenc}
\usepackage[utf8]{inputenc}
\usepackage{lmodern}
\usepackage[margin=1in,headheight=15pt]{geometry}
\usepackage{amsmath,amssymb,amsthm,mathtools}
\usepackage{microtype,booktabs,longtable,array,enumitem}
\usepackage{xcolor,fancyhdr,xurl,aliascnt}
\usepackage[colorlinks=true,linkcolor=blue!45!black,citecolor=green!35!black,urlcolor=blue!45!black]{hyperref}
\usepackage[nameinlink,noabbrev]{cleveref}
\hypersetup{pdftitle={Birthday Cutoffs and Hereditary Sets},pdfauthor={Merged research report for the Surreal project},pdfsubject={Birthday cutoffs, hereditary set universes, bi-interpretation, elementary inclusions, axiom spectra, outer models and surcomplex orientation}}
\pagestyle{fancy}\fancyhf{}
\fancyhead[L]{\small Birthday cutoffs and hereditary sets}
\fancyhead[R]{\small\thepage}
\renewcommand{\headrulewidth}{0.3pt}
\setlength{\parskip}{0.25em}
\setlength{\emergencystretch}{3em}
\setlist{nosep,leftmargin=2em}
\clubpenalty=10000\widowpenalty=10000\displaywidowpenalty=10000
\allowdisplaybreaks[2]\numberwithin{equation}{section}
\setcounter{tocdepth}{1}
\newtheorem{theorem}{Theorem}[section]
\newaliascnt{lemma}{theorem}
\newtheorem{lemma}[lemma]{Lemma}\aliascntresetthe{lemma}
\newaliascnt{proposition}{theorem}
\newtheorem{proposition}[proposition]{Proposition}\aliascntresetthe{proposition}
\newaliascnt{corollary}{theorem}
\newtheorem{corollary}[corollary]{Corollary}\aliascntresetthe{corollary}
\theoremstyle{definition}
\newaliascnt{definition}{theorem}
\newtheorem{definition}[definition]{Definition}\aliascntresetthe{definition}
\newaliascnt{example}{theorem}
\newtheorem{example}[example]{Example}\aliascntresetthe{example}
\newaliascnt{convention}{theorem}
\newtheorem{convention}[convention]{Convention}\aliascntresetthe{convention}
\newaliascnt{question}{theorem}
\newtheorem{question}[question]{Question}\aliascntresetthe{question}
\theoremstyle{remark}
\newaliascnt{remark}{theorem}
\newtheorem{remark}[remark]{Remark}\aliascntresetthe{remark}
\crefname{lemma}{lemma}{lemmas}
\crefname{proposition}{proposition}{propositions}
\crefname{corollary}{corollary}{corollaries}
\crefname{definition}{definition}{definitions}
\crefname{example}{example}{examples}
\crefname{remark}{remark}{remarks}
\crefname{convention}{convention}{conventions}
\crefname{question}{question}{questions}
\crefname{part}{Part}{Parts}
\Crefname{part}{Part}{Parts}
\newcommand{\No}{\mathbf{No}}
\newcommand{\Ord}{\operatorname{Ord}}
\newcommand{\bd}{\operatorname{b}}
\newcommand{\TC}{\operatorname{TC}}
\newcommand{\B}{\mathcal B}
\newcommand{\C}{\mathcal C}
\newcommand{\I}{\mathcal I}
\newcommand{\J}{\mathcal J}
\newcommand{\Code}{\operatorname{Code}}
\newcommand{\Bit}{\operatorname{Bit}}
\newcommand{\Pre}{\operatorname{Pre}}
\newcommand{\Vis}{\operatorname{VisOrd}}
\newcommand{\Coll}{\operatorname{Coll}}
\newcommand{\SG}{\operatorname{SignGraph}}
\newcommand{\OC}{\operatorname{OC}}
\newcommand{\Et}{\mathsf{Et}}
\newcommand{\Pack}{\mathsf{Pack}}
\newcommand{\PowI}{\mathsf{Pow}^{\I}}
\newcommand{\RepB}{\mathsf{Rep}_{\beth}}
\newcommand{\New}{\operatorname{New}}
\newcommand{\Pow}{\mathcal P}
\newcommand{\om}{\omega}
\newcommand{\nsum}{\mathbin{\oplus}}
\newcommand{\nprod}{\mathbin{\otimes}}
\newcommand{\card}[1]{\lvert#1\rvert}
\newcommand{\dom}{\operatorname{dom}}
\newcommand{\ran}{\operatorname{ran}}
\newcommand{\rank}{\operatorname{rank}}
\newcommand{\cf}{\operatorname{cf}}
\newcommand{\id}{\operatorname{id}}
\newcommand{\Aut}{\operatorname{Aut}}
\newcommand{\Emb}{\operatorname{Emb}_{\mathrm{el}}}
\newcommand{\cb}{\operatorname{cb}}
\newcommand{\name}{\operatorname{name}}
\newcommand{\HWO}{\mathrm{HWO}}
\newcommand{\ZFC}{\mathrm{ZFC}}
\newcommand{\ceilcard}[1]{\kappa(#1)}
\newcommand{\src}[1]{\textsf{[#1]}}
\newcommand{\lbl}[1]{\nolinkurl{#1}}
\newcommand{\repo}[1]{\nolinkurl{#1}}
\newcommand{\repoSHA}{\texttt{4f2645599121fa872c7104995e47f86f0382351f}}
\newcolumntype{P}[1]{>{\raggedright\arraybackslash}p{#1}}
\begin{document}
\hypersetup{pageanchor=false}
\begin{titlepage}
\centering
{\large\scshape A merged research report\par}
\vspace{0.2in}
{\Huge\bfseries Birthday Cutoffs\par and Hereditary Sets\par}
\vspace{0.25in}
{\Large The set-theoretic content of bounded surreal arithmetic:\par
bi-interpretations, elementary inclusions, axiom spectra,\par
outer models, and surcomplex orientation\par}
\vspace{0.15in}
{\large September 22, 2026\par}
\vspace{0.1in}
\begin{minipage}{0.96\textwidth}\small
\begin{abstract}
\noindent
Expand the ordered field of surreal numbers of birthday below an ordinal
$\lambda$ by its birthday function. The global connection of this kind of
structure with set theory is the subject of a bi-interpretation theorem
announced by Chen, Hamkins and Yang. This report, merged from five
manuscripts that prove overlapping versions of one local theorem, asks
\emph{which set universe is recovered from a particular birthday cutoff, and
what that recovery implies}.

For an epsilon number $\lambda$ let $\kappa(\lambda)$ be the least cardinal
not below $\lambda$. One parameter-free interpretation recovers exactly
$H_{\kappa(\lambda)}$. At cardinal cutoffs, singular ones included, this is a
uniform bi-interpretation with $H_\lambda$; at noncardinal cutoffs it is a
bi-interpretation with $(H_{\kappa(\lambda)},\in,\lambda)$, and the cutoff is
recovered by a definable collector. Thus the field of birthdays below
$\varepsilon_0$ already recovers all hereditarily countable sets. Small
arithmetic certificates justify the reverse interpretation without Replacement
in $H_\kappa$, and a second route codes graphs modulo bisimulation. The
forward interpretation extends to the ring cutoffs, including $\omega$.

Consequences within one universe: for epsilon numbers $\lambda<\eta$ the
inclusion of birthday structures is elementary exactly when both are cardinals
and $H_\lambda\prec H_\eta$; elementary embeddings correspond to those of the
hereditary universes; one sentence detects noncardinal cutoffs and another
the strong-limit cardinals; birthday eternity is Replacement at cardinal
cutoffs, and with fiber packing it selects the worldly cardinals; in ZF the
graph coding reaches exactly the sets with well-orderable transitive closure.
Across universes: old arithmetic is absolute, a bounded structure is unchanged
exactly when no bounded subset is added, the first new birthday is the least
ordinal with a new subset and is a ground-model cardinal, explicit formulas
detect new subsets, and, assuming a measurable cardinal, Prikry forcing leaves
the enriched field unchanged while singularizing its cutoff. Under GBC the full
structure has no nontrivial elementary self-embedding. For surcomplex numbers
with conjugation and coordinate birthday, each real elementary embedding has
exactly two lifts, and one nonreal parameter, but no set of real parameters,
defines an orientation.
\end{abstract}
\end{minipage}
\vfill
{\footnotesize Built from five manuscripts written independently in relation to the
\repo{VladimirReshetnikov/Surreal} project,\par
all pinned to commit \texttt{4f26455}, and merged after placement commit
\texttt{d4e71b7}.\par
An AI-assisted research draft with written proofs and finite regression checks;
not refereed;\par not a claim of any Lean formalization. The global
bi-interpretation is prior announced work of Chen--Hamkins--Yang.\par}
\end{titlepage}
\hypersetup{pageanchor=true}
\pagenumbering{roman}
\begingroup\small\setlength{\parskip}{0pt}\tableofcontents\endgroup
\clearpage\pagenumbering{arabic}

\part{The cutoff theorem}\label{hset:part:spine}

\section{The question, the results, and their provenance}
\label{hset:sec:intro}

The field structure of the surreal numbers and their construction history
carry very different information. Real-closed-field quantifier elimination
makes all inclusions between real closed subfields elementary in the ordered
field language. A birthday function remembers much more: together with the
order it reads individual sign coordinates, and thereby encodes subsets,
relations and well-founded membership graphs. The issue addressed here is not
merely whether this coding is possible globally, but how its strength changes
when only birthdays below one ordinal are allowed, and what the resulting
dictionary says about elementary inclusions, set-existence axioms, outer
models of set theory, and the surcomplex numbers.

\subsection{What is already known, and credited}
\label{hset:sub:prior}

\paragraph{The global theorem is prior work.}
Chen, Hamkins, and Yang announced that the surreal ordered field with birthday
order is bi-interpretable with the set-theoretic universe, and described a
first-order theory of surreal arithmetic intended to be bi-interpretable with
ZFC. Hamkins's September and November 2025 slides give the global coding by
pointed well-founded extensional relations, together with birthday-growth
principles; the March 2026 workshop announcement and slides repeat the global
theorem and label the axiomatic work as still being checked
\cite{CHYKobe,CHYNotreDame,CHYCUNY}. This report does \emph{not} claim that
theorem, its pointed-graph coding idea, or the association of birthday growth
with set-existence axioms as new. The collection's
\repo{foundations-and-computation/foundations} report already records the
March 2026 announcement, at slide level and with its status boundary, in its
subsection \lbl{found:sub:announcement}, and lists it among its non-claims as
work in progress that is not a proved dependency of that report. Of the five
manuscripts merged here, sources 02, 06 and 09 credit the announcement
prominently. Sources 05 and 07 do not cite it; source 07 presented the
full-class bi-interpretation as a proposed new contribution. That presentation
is corrected here: \cref{hset:thm:chy} states the announced theorem with its
credit, and the argument printed there, taken from source 07, is a written
proof of the announced theorem, not a new result.

\paragraph{Other prior inputs.}
If $\lambda$ is an epsilon number, $\No_{<\lambda}$ is a real closed subfield.
The original van den Dries--Ehrlich paper must be used together with its
erratum \cite{vdDE,vdDEerr}; Bournez--Guilmant state the corrected ring, field
and real-closedness results in their Theorems 1.1--1.2 \cite{BG}. The
collection's foundations report states the same classical fact in its
subsection \lbl{found:sub:cutoffs}, including that $\varepsilon_0$ is already an
epsilon number, so that no large cardinal is needed. The sign construction and
recursive arithmetic are classical \cite{Conway,Gonshor}. Mostowski collapse,
well-founded recursion, hereditary-size universes and forcing are standard
\cite{Jech}; quantifier elimination for real closed and algebraically closed
fields is standard \cite{Marker,Deloro}. The general obstruction to
parameter-free bi-interpretation caused by complex conjugation, and the role of
a named imaginary unit, were explained by Je\v{r}\'abek in 2021 \cite{Jerabek};
every use of that obstruction below (\cref{hset:sec:surcomplex}) is credited to
him. Bi-interpretation methodology follows Freire--Hamkins \cite{FreireHamkins}.
The Prikry forcing facts are imported from Poveda--Rinot--Sinapova \cite{PRS},
and the class form of Kunen's inconsistency from
Hamkins--Kirmayer--Perlmutter \cite{HKP}. Rangel--Mariano's algebraic set
theory of surreal numbers and Bertram's realization of numbers as sets of
ordinals are separate programs with different languages and aims; nothing here
settles or refutes them \cite{RM,Bertram}.

\subsection{The principal statements}
\label{hset:sub:results}

We use $\TC(\{x\})$ for the transitive closure containing $x$ and all its
iterated members, and write
\[
 H_\kappa=\{x:\card{\TC(\{x\})}<\kappa\}.
\]
For an epsilon number $\lambda$ (an ordinal with $\omega^\lambda=\lambda$ in
ordinary ordinal exponentiation) put
\[
 \B_\lambda=(\No_{<\lambda};0,1,+,-,\cdot,<,\bd),
\]
where $\bd(x)$ is the birthday of $x$, regarded as the all-plus surreal
ordinal of that length. ``Uniform'' means that fixed finite formulas work at
every cutoff in the stated class, not that a separate language is chosen for
each cardinal. The source numbers in brackets are explained in
\cref{hset:sub:merge}.

\begin{theorem}[Cardinal rounding and cutoff recovery; \cref{hset:thm:biinterpretation}]
\label{hset:thm:main-rounding}
Let $\lambda$ be an epsilon number and let
\[
 \ceilcard{\lambda}=\min\{\kappa:\kappa\text{ is a cardinal and }\lambda\le\kappa\}.
\]
There is one parameter-free interpretation $\I$ such that
\[
 \I(\B_\lambda)\cong(H_{\ceilcard{\lambda}},\in).
\]
If $\lambda$ is a cardinal, $\B_\lambda$ and $(H_\lambda,\in)$ are uniformly
bi-interpretable. If $\lambda$ is not a cardinal, $\B_\lambda$ and
$(H_{\ceilcard{\lambda}},\in,\lambda)$ are uniformly bi-interpretable, and the
named ordinal $\lambda$ is definable, as an interpreted set, in $\B_\lambda$
itself.
\end{theorem}

This is the most general form in the cluster, and it is source 09's. The other
four sources prove cases of it with their own codes:
\begin{itemize}
\item source 02: the forward interpretation onto $H_{\kappa(\lambda)}$ at every
epsilon number, one way only (it explicitly claims no reverse interpretation);
\item source 05: the bi-interpretation at uncountable \emph{regular} cardinals,
with graphs modulo bisimulation (\cref{hset:sec:bisim});
\item sources 06 and 07: the bi-interpretation at every uncountable cardinal,
singular ones included; source 07 also proves the forward interpretation onto
$H_{\kappa(\lambda)}$ at every epsilon number, which it calls a
\emph{cardinality plateau} and states forward only.
\end{itemize}
This is a cardinal-rounding statement, not a rank-cutoff statement: for
example $\kappa(\varepsilon_0)=\omega_1$, whereas $V_{\varepsilon_0}$ is not
the recovered universe. At a singular cardinal $\kappa$ the recovered structure
is still $H_\kappa$; no regularity assumption is hidden in the theorem.

\begin{theorem}[Elementary-inclusion spectrum; \cref{hset:thm:elementary}]
\label{hset:thm:main-elementary}
For epsilon numbers $\lambda<\eta$,
\[
 \B_\lambda\prec\B_\eta
 \quad\Longleftrightarrow\quad
 \lambda,\eta\text{ are cardinals and }H_\lambda\prec H_\eta.
\]
The ordered-field reducts, in contrast, always form an elementary inclusion.
\end{theorem}

Source 09 proves the equivalence at all epsilon numbers. Source 02 proves the
forward implication and left the converse as its first open question; that
question is answered by source 09 (\cref{hset:rem:answered-02}). Sources 05
(regular cardinals), 06 and 07 (uncountable cardinals) prove the equivalence
between cardinal cutoffs, and 05 and 06 add the correspondence of elementary
embeddings (\cref{hset:thm:embeddings}).

\begin{theorem}[Growth spectrum; \cref{hset:thm:worldly}]
\label{hset:thm:main-growth}
For an epsilon-number cutoff, birthday eternity $\Et$ fails at every
noncardinal cutoff. At an uncountable cardinal $\kappa$,
\[
 \B_\kappa\models\Et
 \quad\Longleftrightarrow\quad
 H_\kappa\models\text{Replacement}.
\]
Fiber packing $\Pack$, and also the translated Power Set sentence $\PowI$,
hold exactly at the uncountable strong-limit cardinal cutoffs. Consequently
\[
 \B_\lambda\models\Et+\Pack
 \quad\Longleftrightarrow\quad
 \lambda\text{ is a cardinal and }V_\lambda\models\mathrm{ZFC}.
\]
\end{theorem}

The eternity, packing and worldly statements are source 09's. The translated
Power Set sentence is source 02's, at every epsilon number, and source 06's at
uncountable cardinals (\cref{hset:thm:power}).

\begin{theorem}[Extension rigidity at a strong-limit cutoff; \cref{hset:thm:extension-rigidity}]
\label{hset:thm:main-extension}
Let $M\subseteq N$ be transitive models of ZFC with the same ordinals, and let
$\kappa$ be an uncountable cardinal in both, strong limit in $M$. Then
\[
 \B_\kappa^M\prec\B_\kappa^N
 \quad\Longleftrightarrow\quad
 \B_\kappa^M=\B_\kappa^N
 \quad\Longleftrightarrow\quad
 \Pow^M(\alpha)=\Pow^N(\alpha)\quad(\alpha<\kappa).
\]
Without the strong-limit hypothesis, a new subset of $\alpha<\kappa$ still
prevents elementarity whenever $(2^{\card\alpha})^M<\kappa$.
\end{theorem}

The detector is sharp for the old-power-set \emph{parameter}
(\cref{hset:prop:detector-cost}), not a claim that every other possible
detector needs that much space. The report also proves:
\begin{itemize}
\item the forward interpretation at ring cutoffs, including $\No_{<\omega}$
onto the hereditarily finite sets (\cref{hset:prop:ring}, source 02);
\item one sentence true exactly at noncardinal epsilon cutoffs
(\cref{hset:thm:collector}, sources 02 and 09);
\item a separation of $\B_{\beth_\omega}$ from $\B_{\beth_\omega^+}$ by one
Replacement instance (\cref{hset:thm:beth}, source 06);
\item the exact range of the graph coding over ZF without Choice
(\cref{hset:thm:choice}, source 06);
\item the transfer of elementary equivalence, the interpretation of full
second-order arithmetic, the independence property of the bit formula, an
omitted parameter-free type, and nondefinability of birthday in the field
reduct (\cref{hset:sec:modeltheory}, sources 05, 06, 07);
\item the exact preservation criterion across outer models, the first new
birthday, three explicit new-subset detectors, and the full-class statement
$\B^M\prec\B^N\iff M=N$ (\cref{hset:part:outer}, sources 02, 05, 06, 07, 09);
\item conditional on a measurable cardinal, an unchanged enriched field whose
cutoff changes cofinality, the preserved $H_\kappa=V_\kappa^M$, and the exact
onset of new elements above the cutoff (\cref{hset:sec:prikry}, sources 02 and
06);
\item under GBC, no nontrivial class elementary self-embedding of the full
birthday structure, by transfer of Kunen's inconsistency
(\cref{hset:thm:kunen}, source 07);
\item for surcomplex numbers with conjugation and coordinate birthday: exactly
two elementary lifts of each real elementary embedding (source 06), the exact
parameter threshold for a definable orientation (source 06), the
mutual-interpretation versus bi-interpretation distinction (sources 05, 06, 09,
after Je\v{r}\'abek), and, under GBC, the elementary self-embeddings of the full
structure (source 07) (\cref{hset:sec:surcomplex}).
\end{itemize}

\subsection{Five manuscripts, one report}
\label{hset:sub:merge}

This report is built from five manuscripts, written independently on
September 22, 2026, all pinned to repository commit \repoSHA{} and placed in
the collection by commit \texttt{d4e71b7}. Their numbers are those of the
placement batch and are local to this directory.

\begin{longtable}{@{}P{0.07\textwidth}P{0.33\textwidth}P{0.52\textwidth}@{}}
\toprule
Source & Manuscript & Contribution to this report\\
\midrule\endhead
09 & \emph{The Set-Theoretic Content of Bounded Surreal Arithmetic: hereditary-size universes, coherent bi-interpretations, axiom spectra, and extension rigidity} &
The base. Its text is the spine of Part~I
(\cref{hset:sec:foundations,hset:sec:certificates,hset:sec:bits,hset:sec:codes,hset:sec:biinterpretation}),
the elementary-inclusion classification (\cref{hset:sec:elementary}), the
eternity, packing and worldly calibration (\cref{hset:sec:growth}), the
cross-model coherence, detector and extension rigidity
(\cref{hset:sec:detectors}), the orientation obstruction
(\cref{hset:thm:orientation}), and most of \cref{hset:sec:scope}. Files prefixed \texttt{09-bounded-arithmetic-}.\\[0.4em]
02 & \emph{What a Birthday Cutoff Remembers: hereditary set universes, elementary obstructions, and forcing invisibility} &
The one-way forward theorem at all epsilon numbers (a special case of the
base), ring cutoffs (\cref{hset:sec:rings}), the collector sentence
(printed once with 09's), the translated Power Set sentence at every epsilon
number, the preservation criterion at epsilon cutoffs and its countable
corollary, the Prikry invisibility theorem with its preserved $H_\kappa$ and the
onset above the cutoff (\cref{hset:sec:prikry}), the concrete example in
Appendix~\ref{hset:app:example}. Files prefixed \texttt{02-birthday-cutoffs-}, and its
source audit \path{02-birthday-cutoffs-SOURCE_AUDIT.md}.\\[0.4em]
05 & \emph{Birthdays Recover Sets: bi-interpretability, elementary embeddings, and forcing in surreal and surcomplex fields} &
The second route by codes modulo bisimulation (\cref{hset:sec:bisim}), the
explicit embedding formula and critical-point corollary, the separation of
$\B_{\omega_1}$ and $\B_{\omega_2}$, nondefinability of birthday, the explicit
missing-row detector $\New$, the full-class statement $\B^M\prec\B^N\iff M=N$,
the first new birthday for arbitrary same-ordinal outer models, the countably
closed example, and the surcomplex forcing contrast. Files prefixed
\texttt{05-birthdays-recover-sets-}.\\[0.4em]
06 & \emph{What Birthday-Enriched Surreal Fields Remember: hereditary-size interpretations, choice, forcing, and surcomplex orientation} &
The bi-interpretation at singular cardinals (printed once with the base),
the second route by prefix-pair arithmetic tables (\cref{hset:sub:tables}),
definitional equivalence with $(\No_{<\kappa};<,\bd,G_\varpi)$, transfer of
elementary equivalence, the $\beth_\omega$ Replacement instance, the ZF range
theorem, the four-way preservation criterion, the Prikry corollaries, the two
complex lifts, and the orientation threshold. Files prefixed
\texttt{06-hereditary-size-orientation-}.\\[0.4em]
07 & \emph{The Birthday Expansion of the Surreal Field: recovering set theory, detecting forcing, and rigidity of surcomplex embeddings} &
The written proof of the announced full-class theorem (\cref{hset:thm:chy}),
the radix pairing, the $\varepsilon_0$ localization (printed once as a case of
the base), recovery of models of ZFC, stage agreement at every ordinal, the
two-parameter outer-model witness, the independence property and omitted type,
the o-minimal nondefinability, Kunen transfer (\cref{hset:thm:kunen}) and the
surcomplex elementary self-embeddings. Files prefixed
\texttt{07-birthday-expansion-}.\\
\bottomrule
\end{longtable}

\paragraph{Why one report.} All five prove the same spine: the ordered field
of surreals below a cutoff, expanded by birthday, interprets the hereditary
sets $H_\kappa$ by sign bits read from order and birthday, a bounded ordinal
pairing polynomial, and pointed well-founded relation codes, and in most
versions is bi-interpretable with them. They differ in the cutoff class and in
whether the result is one-way or two-way. Source 09 has the widest cutoff class
(all epsilon numbers, which include all uncountable cardinals by
\cref{hset:lem:card-estimates}) and the strongest conclusion (two-way, with the
correct pointed target), so it is the base.

\paragraph{Printed once.} The prefix and bit formulas (all five), the bounded
pairing (all five, in four variants compared in \cref{hset:rem:pairings}), the
rooted collapse and exact range (02, 06, 07, 09), the arithmetic certificates
(05, 06, 07, 09), the bi-interpretation at cardinal cutoffs (05 regular; 06, 07,
09 all), the $\varepsilon_0$ case (02, 07, 09), the collector sentence (02, 09),
the elementary-inclusion classification (02 one direction; 05, 06, 07, 09),
the translated Power Set threshold (02, 06, with 09's version inside
$H_\kappa$), absoluteness of old arithmetic (02, 05, 07, 09), the preservation
criterion (02, 06, 07, 09), the first new birthday (02, 05, 07), the
outer-model non-elementarity (05, 07, 09), the Prikry invisibility (02, 06),
rigidity of the bounded structures (05, 06, 07, 09), the surcomplex real-axis
recovery (all five) and the two-element automorphism group (05, 06, 07, 09).
Each is printed once and credited to every source that proves it.

\paragraph{Kept twice, as different proofs.} The reverse arithmetic has two
routes: source 09's option-graph certificates built from finite signed words
(\cref{hset:thm:certificates}), and the prefix-pair tables of sources 05, 06 and
07 (\cref{hset:sub:tables}), whose single-$\mu$ size bound is source 06's
argument at singular cutoffs. The coding of sets has two routes: source 09's
extensional rooted codes with isomorphism as equality (\cref{hset:sec:codes}),
and source 05's codes modulo bisimulation (\cref{hset:sec:bisim}). The first new
birthday has two existence arguments: the new-subset lemma of 05 and 07 for any
outer model, and the generic filter of 02 for set forcing
(\cref{hset:rem:firstbirthday-02}). The new-subset detectors are three different
formulas (\cref{hset:sec:detectors}).

\paragraph{Where the merge had to choose.}
(i) Hypotheses: whenever a result is proved by several sources under different
hypotheses, the weakest proved hypothesis is printed, and the proof printed is
one that uses only it; where this attaches a source's argument to a weaker
hypothesis than the source stated, the fact is said at the point of use
(\cref{hset:rem:05-regular,hset:rem:bounded-new}). (ii) The first new birthday
is stated for arbitrary same-ordinal transitive outer models (05, 07), with
02's set-forcing theorem as a special case. (iii) The full-class
bi-interpretation is printed as the announced theorem of Chen--Hamkins--Yang,
not as a result of source 07. (iv) Symbols were renamed as listed in
\cref{hset:sub:conventions}.

\subsection{Conventions and renamed symbols}
\label{hset:sub:conventions}

\begin{convention}[One meaning per symbol]\label{hset:conv:notation}
The following conventions hold throughout, including in material taken from
sources that used other letters.
\begin{center}\small
\begin{longtable}{@{}P{0.27\textwidth}P{0.34\textwidth}P{0.31\textwidth}@{}}
\toprule
Notion & Here & In the sources\\
\midrule\endhead
Birthday of $x$, as a surreal ordinal & $\bd(x)$ & $\mathsf b$ (02); $b$ (05, 07, 09); $b$ for the ordinal and $\mathbf b$ for its surreal (06)\\
Cutoff carrier & $\No_{<\lambda}$ & $F_\lambda$ (02), $S_\kappa$ (05, 06), $\No_{<\kappa}$ (07), $K_\lambda$ (09)\\
Birthday structure & $\B_\lambda$; full class $\B$ & $\mathfrak B_\lambda$ (02), $\mathcal S_\kappa$ (05), $\mathcal A_\kappa$ (06), $\B_\kappa$ (07, 09)\\
Cardinal ceiling & $\kappa(\lambda)$ & $\theta(\lambda)$, the least \emph{uncountable} cardinal $\ge\lambda$ (07); equal for epsilon $\lambda$\\
Pairing length & $\Lambda(\theta)=\theta^2+\theta$ & $B_\theta$ (02); $L(a,d)$ (05); $\tau(\delta)$ (06); $q^2$ (07)\\
Code & $C=(\delta,r,t)$: domain, relation word, root & $(\theta,d,e,r)$ (02); $(\lambda,c,r)$ with $\lambda$ a node count (05); $(\delta,e,r)$ (06); $(q,d,e,r)$ (07)\\
Interpreted $=$, $\in$ & $\sim$, $\in_\I$ & $\mathsf{Eq},\mathsf{Mem}$ (02); $\equiv,\in^*$ (06); $\simeq,\in_*$ (07)\\
Ordinal code & $\OC(a)$ & $\mathsf{OC}$ (02), $\operatorname{Sync}$ (05), $\mathsf{OrdCode}$ (06), $\operatorname{OCode}$ (07), $O(a)$ (09)\\
Collector & $\Coll(C)$; sentence $\exists C\,\Coll(C)$ & $\mathsf{CollectOrd}$ (02), $\operatorname{Coll}$ (09)\\
Replacement instance at $\beth_\omega$ & $\RepB$ & $\mathsf{Collect}_\beth$ (06)\\
Translated Power Set & $\PowI$ & $\mathsf{Power}$ (02), $\mathsf{Pow}^\sharp$ (06)\\
Eternity, fiber packing & $\Et$, $\Pack$ & $\mathsf E$, $\mathsf P$ (09)\\
Conjugation & $c$ & $c$ (02, 06); overline (05, 07); $\sigma$ (09)\\
Coordinate birthday & $\cb(x+iy)=\max(\bd x,\bd y)$ & $\beta$ (02, 05, 09); $B$ (06, 07)\\
Surcomplex structures & $\C_\lambda$ (no $i$ named), $\C^i_\lambda$ ($i$ named); full $\C$, $\C^i$ & $\mathfrak K_\lambda$ with $i$ (02); $\C_\kappa,\C^-_\kappa$ (05); $\mathcal K_\kappa$ (06); $\C_b,\C_b^i$ (07); $\C^\pm_\lambda$ (09)\\
First new birthday & $\beta=\beta(M,N)$ & $\delta$ (02, 05, 07)\\
New-subset formula of 05 & $\New(m,\theta,\nu)$ & $\New(c,\theta,\lambda)$, $\lambda$ an enumeration length (05)\\
Eternity bound & $\varrho$ & $\beta$ (09)\\
\bottomrule
\end{longtable}
\end{center}
\end{convention}

Three collisions with the collection's shared notation are resolved as follows
(\repo{docs/NOTATION.md}). The collection reserves $F_\Gamma$ and $K_\Gamma$ for
the real and complex Hahn workspaces and $K$ for a set-sized \emph{complex}
field; sources 02 and 09 used $F_\lambda$ and $K_\lambda$ for birthday cutoffs
(09's $K_\lambda$ a \emph{real} field), so this report writes $\No_{<\lambda}$
and $\No_{<\lambda}[i]$ instead. The collection's $H_\alpha$ is a subgroup of a
value group (bounded multiples of $\alpha$); here $H_\kappa$ is always the
universe of sets of hereditary size below $\kappa$, and no value group
occurs. The collection's $\mathcal O$ is a valuation ring, hence the ordinal
code is written $\OC$. The letter $\beta$ is reserved for the first new
birthday of an outer model, as in the sibling report
\repo{foundations-and-computation/surreal-fields-across-universes}; that report
uses $\delta$ for the first new \emph{ordinal-sequence length}, a notion not
named here, while here $\delta$ is only the domain ordinal of a code. Sources
02, 05 and 07 wrote $\delta$ for the first new birthday.

Two further words have two meanings in the sources. ``Collect'' named a
sentence collecting the literal ordinals in source 02 and a Replacement instance
in source 06; the two are $\exists C\,\Coll(C)$ and $\RepB$ here. A
sans-serif $\mathsf P$ was fiber packing in source 09 and Power Set in sources
02 and 06; they are $\Pack$ and $\PowI$ here, and $\Pack$ is \emph{not} the
Power Set axiom. Birthday precedence and prefix simplicity are different
relations (\lbl{found:rem:simplicity} of the foundations report); prefix is
defined below from order and birthday, and is never identified with
precedence.

\subsection{Relation to the collection and corrected repository statements}
\label{hset:sub:repo}

\paragraph{Foundations.} The foundations report records the Chen--Hamkins--Yang
announcement (\lbl{found:sub:announcement}), the epsilon-number cutoffs
(\lbl{found:sub:cutoffs}) and the distinction between simplicity and birthday
precedence (\lbl{found:rem:simplicity}). Its inventory notes that the field
language has no birthday symbol; \cref{hset:prop:nondef} supplies the standard
nondefinability proof. None of the interpretation theorems of this report was
in the collection before it.

\paragraph{Surcomplex automorphisms.} The report
\repo{surcomplex/surcomplex-field-automorphisms} proves $\Aut(K/F)=\{\id,c\}$
for $K=F(i)$ with $F$ real closed (\lbl{saut:thm:axis}) and, in its section
\lbl{saut:sec:exponential}, subsection ``Simplicity restores the canonical
surreal numbers'', that a bijection preserving numerical order and prefix
simplicity is the identity, so that a real-axis-preserving automorphism whose
real restriction preserves simplicity is the identity or conjugation. Since
order and birthday define prefix (\cref{hset:lem:prefix}), the bounded
automorphism statements of sources 05, 06, 07 and 09 are consequences of those
results; they are printed once, as \cref{hset:lem:rigidity,hset:thm:orientation},
and credited there. What is new relative to that report is only source 07's
classification of the class elementary self-embeddings of the full structure,
which are not assumed surjective (\cref{hset:thm:complex-rigid}), and source 06's
two-lift classification of elementary embeddings between different cutoffs
(\cref{hset:thm:two-lifts}). No question named in that report is answered or
refuted; in particular its question whether $\{\id,c\}$ is a proper subgroup of
its group $\operatorname{Aut}_L$ of field automorphisms commuting with
$L(z)=\log\lvert z\rvert$ (equivalently, whether a nontrivial exponential
automorphism of $\No$ exists) is untouched.

\paragraph{Product birthdays.} No argument here uses a bound
$\bd(xy)\le\bd(x)\nprod\bd(y)$. That bound is Gonshor's product-birthday
conjecture; the collection's \repo{surreal/gonshor-product-birthdays} report
proves it on a specified ring of ordinal-support series and states that the
general conjecture is not settled. Sources 05, 06, 07 and 09 all say that they
avoid it, and the size estimates below use only field closure of the
epsilon-number cutoffs.

\paragraph{Sibling report.} The report
\repo{foundations-and-computation/surreal-fields-across-universes}, placed in the
same batch, studies the external saturation of the \emph{pure} ordered field
$\No^M$ inside an outer universe, without a birthday symbol. Its setting meets
this report's outer-model part (\cref{hset:part:outer}): the absoluteness of old
arithmetic, the first new birthday, and the forcing examples. The two languages
give opposite phenomena: the pure-field inclusion stays elementary but can lose
saturation, while the birthday expansion loses elementarity
(\cref{hset:rem:sibling}).

\paragraph{Other reports.} The \repo{foundations-and-computation/computable-surreals}
report concerns countable effective subfields and is a different program; the
undecidability remarks of \cref{hset:sec:modeltheory} are standard. No report of
the collection is contradicted by this one.

\paragraph{Repository statements of the sources, corrected.}
All five sources compared themselves with the collection at the pin \repoSHA,
which is 32 commits before the placement commit \texttt{d4e71b7}.
\begin{itemize}
\item Sources 05 and 07 report that the catalogue at the pin described forty
catalogued reports and four newly placed drafts, and that the formalization
ledger counted forty-four main texts. This was true at the pin. Before this
batch was placed the catalogue listed 44 research reports in five families.
\item Sources 05 and 07 report that no matching reconstruction theorem was
located in the material they inspected, and source 07 that the uniform
bi-interpretation had not been located in earlier published work or in the
repository. Source 05 read only the first 230 lines of the foundations report
and source 07 its introduction; the Chen--Hamkins--Yang announcement is recorded
in that report at the pin and now. The full-class statement is therefore
credited to the announcement (\cref{hset:thm:chy}), and the local refinements of
05 and 07 carry the same priority caveat as those of 02, 06 and 09.
\item Sources 05 and 06 cite line ranges and section structure of the
surcomplex automorphisms report at the pin. That report has since been
rewritten; its subsection ``Simplicity restores the canonical surreal numbers''
still exists, now in its section \lbl{saut:sec:exponential}, and the relative
result is its \lbl{saut:thm:axis}.
\item Source 06 reports that a GitHub search for ``Hamkins'' returned no
indexed matches. That was a search-index observation; a direct search of the
repository finds the announcement in the foundations report.
\item Sources 02 and 07 report that the incoming-report directory contained
only its README at the pin; that was true at the pin.
\end{itemize}
Source 02's own source audit, kept byte-identical as
\path{02-birthday-cutoffs-SOURCE_AUDIT.md}, contains no statement about the
repository that is false; its closing description of the elementary-inclusion
criterion as ``necessary, not claimed sufficient'' describes source 02 and is
superseded inside this report by \cref{hset:thm:elementary}.

\section{Foundations and the two ordinal languages}
\label{hset:sec:foundations}

\subsection{Ambient theory and class conventions}

The proofs take place in ZFC, except in \cref{hset:sec:choice} (ZF) and in
\cref{hset:sub:kunen} (GBC, stated there). All cutoff structures are sets. The
statements about the full surreal class are definable-class statements,
understood one first-order formula at a time. Neither a set of all surreals nor
an internally available truth predicate for the entire universe is assumed.
Quotient interpretations use definable equivalence relations; the Mostowski
collapse supplies set representatives when we compare them with $H_\kappa$.
Global Choice is not used outside \cref{hset:sub:kunen}.

When a transitive model $M$ is explicitly mentioned, its surreals are the sign
sequences belonging to $M$, and cardinalities with superscript $M$ are computed
there. A statement proved internally in an arbitrary model of ZFC does not
assert that the model's well-founded relations are externally well founded.
The outer-model comparison theorems explicitly use transitive models.

\subsection{The carrier and the birthday function}

A surreal is a function $s:\alpha\longrightarrow\{-,+\}$ for an ordinal
$\alpha$. Comparison is lexicographic, with termination between $-$ and $+$.
The all-plus sequence of length $\alpha$ is the embedded ordinal $\alpha$.
Throughout, $\bd(s)$ is the \emph{surreal ordinal} corresponding to the length
of $s$. Thus it is a function on the surreal carrier, not an extra ordinal
sort. Set-theoretic lengths and their surreal embeddings are distinguished
by context.

For an ordinal $\lambda$ put $\No_{<\lambda}=\{x\in\No:\bd(x)<\lambda\}$, a set.
For an epsilon number $\lambda$ define
\[
 \B_\lambda=(\No_{<\lambda};0,1,+,-,\cdot,<,\bd),
 \qquad
 \B=(\No;0,1,+,-,\cdot,<,\bd).
\]
The language includes no exponentiation, support predicate, membership
predicate, or second-order quantifier.

\begin{lemma}[Birthday order and birthday function]
\label{hset:lem:birthday-language}
On each $\B_\lambda$, and on $\B$, the birthday function and the
birthday-precedence relation $x\prec_b y\iff\bd(x)<\bd(y)$ are definitionally
interconvertible. So is the equal-birthday relation
$E_b(x,y)\iff\bd(x)=\bd(y)$. Moreover, the embedded ordinals are exactly the
fixed points of $\bd$.
\end{lemma}
\begin{proof}
Given the function, precedence and equal birthday are immediate. Conversely,
equality of birthdays is expressed by neither $x\prec_b y$ nor $y\prec_b x$. In
the class of sign sequences of length $\alpha$, the largest in numerical order
is the all-plus sequence: any other sequence of that length has a first minus
sign and is smaller. Therefore $\bd(x)$ is the unique largest member of $x$'s
equal-birthday class,
\[
 y=\bd(x)\iff E_b(x,y)\land\forall z\,(E_b(x,z)\to z\le y),
\]
and this uses only $E_b$ and the order. A sign sequence is fixed by this
function exactly when it is all plus. Each cutoff contains the all-plus
representative of each of its birthdays, so the definitions work in every
$\B_\lambda$.
\end{proof}

The interconvertibility with precedence is proved in sources 02, 06 and 09;
the sufficiency of the coarser equal-birthday relation is source 07's
observation. In particular the results below concern a definitional variant of
the birthday-precedence language of the Chen--Hamkins--Yang announcement, not
an essentially stronger coding apparatus. We abbreviate $\Ord(a)$ by
$\bd(a)=a$. Quantification over embedded ordinals in $\No_{<\lambda}$ is
precisely quantification over $\alpha<\lambda$.

\subsection{Natural arithmetic, not ordinary ordinal arithmetic}

For ordinal arguments, surreal addition and multiplication are the
Hessenberg operations
\[
 \alpha+\gamma=\alpha\nsum\gamma,\qquad
 \alpha\gamma=\alpha\nprod\gamma.
\]
They are strictly increasing in each positive argument where appropriate,
and $\alpha+1$ agrees with ordinary ordinal successor. We use ordinary
ordinal exponentiation only in metamathematical bounds and in explicitly
set-theoretic constructions. It is not silently identified with a field
operation.

\begin{lemma}[Closure and cardinal estimates]
\label{hset:lem:card-estimates}
Every uncountable cardinal is an epsilon number. For infinite ordinals
$\alpha,\gamma$, natural sums and products have cardinality at most
$\max(\aleph_0,\card\alpha,\card\gamma)$. In particular, if $\kappa$ is an
uncountable cardinal and $\alpha,\gamma<\kappa$, their natural sum and product
are below $\kappa$.
\end{lemma}
\begin{proof}
In Cantor normal form, a natural sum joins finitely many exponent terms; a
natural product uses finitely many natural sums of their exponents. The
cardinality estimate follows from the corresponding estimate for
$\omega^\xi$: for infinite $\xi$, its cardinality is at most
$\max(\aleph_0,\card\xi)$. One can see the latter by representing ordinals
below $\omega^\xi$ by finite Cantor-normal-form lists with exponents below
$\xi$ and natural-number coefficients. These observations also cover
finite arguments after adjoining $\aleph_0$ to the bound.

If $\kappa$ is uncountable and initial, $\omega^\alpha<\kappa$ for every
$\alpha<\kappa$. Continuity gives $\omega^\kappa\le\kappa$. The opposite
inequality follows from $\omega^\alpha\ge\alpha$ and taking suprema. Hence
$\omega^\kappa=\kappa$. Neither step uses regularity of $\kappa$.
\end{proof}

All five sources prove the first assertion; van den Dries and Ehrlich already
note that every uncountable cardinal is an epsilon number \cite{vdDE}, as
source 06 records. For a general epsilon number, closure of natural
arithmetic below $\lambda$ also follows from the classical field-cutoff
theorem. This simple closure is all the ordinal arithmetic needed by the
relation codes.

\subsection{Hereditary size rather than rank}

\begin{lemma}[Basic hereditary-size facts]
\label{hset:lem:hereditary}
Let $\kappa$ be an uncountable cardinal.
A sign sequence of length $\alpha$ belongs to $H_\kappa$ exactly when
$\alpha<\kappa$. Every subset of an ordinal $\alpha<\kappa$ belongs to
$H_\kappa$. Also $H_\kappa\subseteq V_\kappa$, and $\Ord\cap H_\kappa=\kappa$.
\end{lemma}
\begin{proof}
The transitive closure of the usual set-coded graph of a sign sequence has
cardinality at most $\max(\aleph_0,\card\alpha)$, and it contains enough of
the domain ordinal to have cardinality at least $\card\alpha$ when $\alpha$
is infinite. The same estimate applies to a subset of $\alpha$, and
$\TC(\{\alpha\})=\alpha+1$.

If $\card{\TC(\{x\})}\le\mu$ for an infinite $\mu<\kappa$, the membership
relation on this transitive closure is a well-founded relation of size at
most $\mu$. Its rank function takes values below $\mu^+$. This follows by
well-founded induction, using regularity of the successor cardinal $\mu^+$:
the supremum of at most $\mu$ ordinals below $\mu^+$ remains below $\mu^+$.
Thus $\rank(x)<\mu^+\le\kappa$, and $x\in V_\kappa$.
\end{proof}

The statement that every subset of $\alpha$ belongs to $H_\kappa$ does
\emph{not} say that the collection $\Pow(\alpha)$ is an element of
$H_\kappa$. This distinction drives the growth results. Three notions must
not be interchanged: a family whose individual members are small, a family
that is a small set, and a family that belongs to $H_\kappa$. For example,
all subsets of $\omega$ belong to $H_{\omega_1}$, but their collection does not
(source 06).

\section{Small certificates for surreal arithmetic}
\label{hset:sec:certificates}

A reverse interpretation cannot simply say ``perform the usual set-theoretic
recursion inside $H_\kappa$.'' At a singular cutoff, Replacement may fail in
that structure. We instead prove that each individual arithmetic computation
has a certificate of the same hereditary cardinal scale as its inputs. Two
routes are given: source 09's option-graph certificates
(\cref{hset:thm:certificates}) and the prefix-pair tables of sources 05, 06
and 07 (\cref{hset:sub:tables}).

\subsection{Legal cuts can be checked using prefixes}

For sets of sign sequences $L<R$, let $\{L\mid R\}$ denote their simplest
separator. The classical simplicity theorem supplies this separator and the
bound
\begin{equation}
 \bd(\{L\mid R\})\le
 \sup\{\bd(a)+1:a\in L\cup R\}.
 \label{hset:eq:cut-bound}
\end{equation}
The supremum of the empty set is zero. These statements can equivalently be
proved by the sign-by-sign cut construction \cite{Gonshor}.

\begin{lemma}[Prefix test for a cut value]
\label{hset:lem:cut-test}
A sign sequence $z$ equals $\{L\mid R\}$ if and only if it separates $L$ and
$R$ and no proper sign prefix of $z$ separates them.
\end{lemma}
\begin{proof}
The forward implication is immediate. For the converse, suppose $y$ is a
separator shorter than $z$. If $y$ is a prefix of $z$, the stated test fails.
If it is not, their longest common prefix lies numerically between them and
is a proper prefix of $z$. The set of separators is convex, so this common
prefix is another separator. Again the test fails. Thus $z$ has minimum
birthday among separators, and uniqueness of the simplest separator gives
the result.
\end{proof}

The check in this lemma quantifies only over the displayed option sets and
the prefixes of $z$. In set coding it is a bounded check: one does not
quantify over all sign sequences of earlier birthdays.

\subsection{Well-founded option graphs of small size}

A labeled option graph has a set $D$ of vertices, left and right successor
relations, and a rank label that strictly decreases along every option edge.
A valuation of such a graph assigns a sign sequence to each vertex so that
its value passes \cref{hset:lem:cut-test} for the values at its options. When
the graph is numeric, this valuation exists and is unique.

\begin{lemma}[Small numeric graphs have small certificates]
\label{hset:lem:graph-size}
Suppose $\mu$ is infinite and a numeric, well-founded option graph has at
most $\mu$ vertices, with its structural data of hereditary size at most
$\mu$. Then it has a rank-labeled valuation certificate of hereditary size
at most $\mu$.
\end{lemma}
\begin{proof}
The graph rank of every vertex is below $\mu^+$, by the successor-cardinal
argument in \cref{hset:lem:hereditary}. Applying \eqref{hset:eq:cut-bound} by
induction shows that a vertex value has birthday at most that vertex's graph
rank. Thus each rank label and each sign value has hereditary size at most
$\mu$. There are at most $\mu$ vertices and at most $\mu^2=\mu$ edges, so the
union of all these data still has hereditary size at most $\mu$. The existence
construction here takes place in the ambient ZFC universe, not in a possibly
deficient $H_\kappa$.
\end{proof}

\subsection{The computation graph for addition and multiplication}

Take all sign prefixes of the inputs, including the inputs themselves. With
\[
 \mu=\max(\aleph_0,\card{\bd(x)},\card{\bd(y)}),
\]
this prefix family and its data have hereditary size at most $\mu$.
For addition, use finite signed words of prefix atoms $\mathsf{a}(a)$; an option moves
one atom to one of its sign-prefix options, reversing left and right for a
negative atom. The root word is $\mathsf{a}(x)+\mathsf{a}(y)$. This is the usual disjunctive-sum
construction of the numeric game for $x+y$.

For multiplication, use finite signed words of pair atoms $\mathsf{p}(a,a')$, with
$a,a'$ prefixes of the respective inputs. An option of a positive pair atom
is a signed three-term word
\begin{equation}
 \mathsf{p}(a^\circ,a')+\mathsf{p}(a,a'^\circ)-\mathsf{p}(a^\circ,a'^\circ).
 \label{hset:eq:pair-option}
\end{equation}
It is a left option when $(a^\circ,a'^\circ)$ is left--left or right--right,
and a right option for the two mixed choices. Negative atoms exchange the sides
and negate the replacement word. An option of a word replaces one atom by one
of these words. The root is $\mathsf{p}(x,y)$. This spells out Conway's multiplication
recursion without first constructing a large universe of games.

There are at most $\mu$ possible finite words. These graphs are well founded.
For addition, use the natural sum, over atoms of the word, of
$\omega^{\bd(a)}$. For multiplication put
\[
 \rho(a,a')=\bd(a)\nsum\bd(a')
\]
and use the natural sum of $\omega^{\rho(a,a')}$ over the pair atoms.
Every pair atom in \eqref{hset:eq:pair-option} has strictly smaller $\rho$ than
the replaced atom. A finite sum of lower powers of $\omega$ is smaller than the
replaced power, so each move strictly decreases the displayed measure.
The standard numeric-game arithmetic theorem identifies the values of these
graphs with $x+y$ and $xy$; their options are exactly the classical recursive
options, including the signs in the negative summands \cite{Conway}.

\begin{theorem}[Hereditary-size arithmetic certificates]
\label{hset:thm:certificates}
There are fixed first-order set-theoretic formulas defining surreal addition,
negation, and multiplication on sign sequences such that, for inputs of
hereditary size at most an infinite $\mu$, a correct computation has a
certificate of hereditary size at most $\mu$. Verification of a supplied
certificate is absolute between transitive sets containing its data.
Consequently, for every uncountable cardinal $\kappa$, these formulas define
in $H_\kappa$ exactly the external surreal operations on its sign sequences.
No assumption $H_\kappa\models\mathrm{Replacement}$ is required.
\end{theorem}
\begin{proof}
For $+$ and multiplication, existentially quantify a set of word vertices
containing the root and closed under the explicit option rule, its exact
option relations, an ordinal rank labeling decreasing on edges, and a
valuation by sign sequences passing the prefix cut test at every vertex.
The node labels specify the input prefixes and the finite signed words; all
checks of option generation quantify over those input prefix families and
over entries of the finite words. Exactness requires both that each generated
option is present and that every recorded edge is generated by the rule.
The predicates for sign sequence, prefix, numerical comparison, finite word,
and decreasing ordinal rank have bounded verifications on the supplied
sets. In particular, a supplied prefix family is not accepted merely because
its members are prefixes: for each ordinal $\xi\le\dom(s)$ it must contain
the unique restriction $s\mathbin{\upharpoonright}\xi$, and every recorded
prefix must be one of these restrictions. The same completeness check is
imposed on the prefix families of vertex values used by the cut test.
These are bounded checks against $\dom(s)+1$ and the supplied families.
A word length is a finite ordinal; this is checked by requiring every
ordinal at most that length to be zero or a successor. Thus a missing
prefix or a nonstandard-length word cannot evade an option requirement.
The verifier is consequently a fixed bounded set-theoretic predicate,
after including these domains and prefix families as certificate data.
Their cardinality and hereditary-size bounds remain $\mu$.

Any valid certificate is sound: external well-founded induction on its
rank labels identifies each vertex value with the value of its generated
numeric game. This also proves uniqueness of the root value. A certificate
exists by the preceding graph construction and \cref{hset:lem:graph-size},
with hereditary size at most $\mu$. Negation simply reverses all signs and has
the same size estimate.

For inputs in $H_\kappa$, choose the infinite bound $\mu<\kappa$ given above.
The external certificate belongs to $H_\kappa$; conversely any certificate
there is externally sound by bounded absoluteness. Hence the existential
formulas have the same intended values. The uses of cardinal regularity were
at $\mu^+$, not at $\kappa$.
\end{proof}

\begin{remark}[Scope of the certificate theorem]
This is not a proof that every class-recursive surreal operation is absolute
in every hereditary-size structure. It proves the specific finite-signature
claim needed here. Inversion is definable from multiplication by $xy=1$;
its existence in the epsilon-number cutoffs, and real closedness, are supplied
by the classical cutoff theorem. No normal-form birthday estimates or
unproved sharp product-birthday conjecture enter the argument
(\cref{hset:sub:repo}).
\end{remark}

\subsection{Second route: prefix-pair tables}
\label{hset:sub:tables}

Sources 05, 06 and 07 certify arithmetic by a single table indexed by pairs of
input prefixes. Source 05 states it at regular cardinals; source 06 isolates
the single-cardinal size bound that makes it work at every uncountable
cardinal, singular ones included, and source 07 uses the same bound. We print
this route with 06's bound.

For a sign sequence $x$ write $\operatorname{Pf}_x=\{x\mathbin{\upharpoonright}\alpha:
\alpha\le\bd(x)\}$ for its prefix set, including $x$, and let $L_x$ and $R_x$
be its proper prefixes below and above $x$; then $x=\{L_x\mid R_x\}$. Say that
$w$ is the \emph{simplest separator} of two families when every left member is
below $w$, every right member is above $w$, and no sign sequence of smaller
length has the same property. For any candidate $w\in H_\kappa$, all sign
sequences shorter than $w$ also lie in $H_\kappa$ (\cref{hset:lem:hereditary});
so this minimality test is absolute between $H_\kappa$ and $V$ whenever the
families are in $H_\kappa$. Uniqueness of a least-birthday separator follows
from signs: two distinct separators of the same least birthday would have a
common prefix at their first difference strictly between them, still
separating and of smaller birthday.

\begin{lemma}[Hereditary size of arithmetic tables]\label{hset:lem:tables}
Let $\kappa$ be an uncountable cardinal, let $x,y\in\No_{<\kappa}$ and set
\[
 \mu=\max\{\aleph_0,\card{\bd(x)},\card{\bd(y)}\}<\kappa.
\]
The functions on $\operatorname{Pf}_x\times\operatorname{Pf}_y$ given by $T_{x,y}(u,v)=u+v$ and
$U_{x,y}(u,v)=uv$ have graphs of hereditary size at most $\mu$; in particular
they belong to $H_\kappa$.
\end{lemma}
\begin{proof}
There are at most $\mu$ pairs of prefixes. Every prefix belongs to
$\No_{<\mu^+}$, which is a field by \cref{hset:lem:card-estimates} and the
classical cutoff theorem. Therefore every value in either table has birthday
below $\mu^+$ and hereditary size at most $\mu$. A set of at most $\mu$ objects,
each of hereditary size at most $\mu$, has hereditary size at most $\mu$; the
finite layers needed for ordered pairs and triples do not change this.

This argument is carried out in the ambient ZFC universe to establish the
existence of a particular witness in $H_\kappa$. It does not invoke
Replacement inside $H_\kappa$, and it uses a \emph{single} infinite cardinal
$\mu$, not a union of sizes approaching a singular $\kappa$. It also works when
$\mu^+=\kappa$.
\end{proof}

\begin{theorem}[Table definitions of the operations; sources 05, 06, 07]
\label{hset:thm:tables}
Let $\mathsf{Add}(x,y,z)$ say that there is a function $T$ on $\operatorname{Pf}_x\times\operatorname{Pf}_y$
with $T(x,y)=z$ such that every entry obeys
\[
 T(u,v)=\bigl\{T(u^L,v),T(u,v^L)\bigm| T(u^R,v),T(u,v^R)\bigr\}
 \qquad(u^L\in L_u,\ v^L\in L_v,\ u^R\in R_u,\ v^R\in R_v),
\]
the braces meaning the simplest-separator predicate. Let $\mathsf{Mult}(x,y,z)$
say that there is a function $U$ on $\operatorname{Pf}_x\times\operatorname{Pf}_y$ with $U(x,y)=z$ such that
$U(u,v)$ is the simplest separator of the left options
$U(u',v)+U(u,v')-U(u',v')$ with $(u',v')\in L_u\times L_v\cup R_u\times R_v$ and
the right options of the same form with $(u',v')\in L_u\times R_v\cup
R_u\times L_v$, where each sum and difference is certified by $\mathsf{Add}$
and sign reversal. Then for every uncountable cardinal $\kappa$ and all
$x,y,z\in\No_{<\kappa}$,
\[
 H_\kappa\models\mathsf{Add}(x,y,z)\iff z=x+y,
 \qquad
 H_\kappa\models\mathsf{Mult}(x,y,z)\iff z=xy.
\]
\end{theorem}
\begin{proof}
The actual tables are in $H_\kappa$ by \cref{hset:lem:tables} and satisfy the
canonical Conway recursions, proving existence. Conversely, any accepted table
is identified with the actual one by external well-founded induction on the
pair of prefix lengths, ordered so that decreasing either coordinate
decreases the pair: every option entry is earlier, the simplest-separator test
is absolute, and the separator is unique. For multiplication, addition has
already been identified. Each invocation of $\mathsf{Add}$ has its own
hereditarily small witness; first-order satisfaction does not require that a
single set collect the witnesses to all invocations.
\end{proof}

The two routes certify the same operations; the table route uses fewer
auxiliary objects, the option-graph route makes the verifier an explicitly
bounded predicate. Either suffices for the reverse interpretation of
\cref{hset:sec:biinterpretation}.

\begin{remark}[Regularity in source 05]\label{hset:rem:05-regular}
Source 05 states its size bookkeeping and its table lemma for uncountable
regular $\kappa$ and declines to infer a singular-cardinal theorem. Its table
argument uses only the field closure of $\No_{<\mu^+}$ and the single bound
$\mu$, exactly as in \cref{hset:lem:tables}; regularity of $\kappa$ is not used
there. The singular case is sources 06 and 07's theorem, and the base's.
\end{remark}

\section{Definable sign access and a bounded pairing map}
\label{hset:sec:bits}

\subsection{Recovering the sign tree from the two orders}

For $x,y,z$ let $\operatorname{Between}(z;x,y)$ say that $z$ lies in the closed
numerical interval with endpoints $x$ and $y$. Define
\begin{equation}
 \Pre(x,y)\quad\Longleftrightarrow\quad
 \forall z\bigl((\operatorname{Between}(z;x,y)\land z\ne x)
                 \longrightarrow\bd(x)<\bd(z)\bigr).
 \label{hset:eq:prefix}
\end{equation}

\begin{lemma}[Definable prefix relation]
\label{hset:lem:prefix}
In every $\B_\lambda$, and in $\B$, \eqref{hset:eq:prefix} holds exactly when
$x$ is a sign prefix of $y$, allowing equality.
\end{lemma}
\begin{proof}
If $x$ is a prefix of $y$, every other point between them must continue the
sign sequence of $x$; a point diverging earlier would lie outside that
interval. Its birthday is therefore larger.

If $x$ is not a prefix of $y$ and $y$ is a proper prefix of $x$, the endpoint
$z=y$ violates the formula. In the remaining case the longest common prefix
of $x,y$ lies strictly between them and has birthday smaller than that of
$x$. This prefix belongs to the same cutoff. Finally, for $x=y$ the formula
is vacuous, as required.
\end{proof}

\begin{remark}[Four prefix formulas]\label{hset:rem:prefixes}
Each source defines prefix from $<$ and $\bd$ by its own formula; all are
correct. Source 02 requires $\bd(x)\le\bd(y)$ and that every $z$ with
$\bd(z)<\bd(x)$ compares with $x$ as with $y$; source 06 adds the symmetric
comparison. Sources 05 and 07 require $\bd(x)\le\bd(y)$ and that no $z$
strictly between $x$ and $y$ has $\bd(z)\le\bd(x)$ (05) or $\bd(z)<\bd(x)$
(07). In the betweenness form the length conjunct cannot be dropped: for
$x=1$ and $y=0$ no surreal of birthday at most $1$ lies strictly between them,
but $1$ is not a prefix of $0$ (source 05, which checks this case in its finite
suite).
\end{remark}

Define the bit predicate by
\begin{align}
 \Bit(s,\alpha)\quad\Longleftrightarrow\quad &
 \Ord(\alpha)\land\alpha<\bd(s)\ \land\nonumber\\
 &\exists t\bigl(\Pre(t,s)\land\bd(t)=\alpha\land t<s\bigr).
 \label{hset:eq:bit}
\end{align}
The unique prefix of length $\alpha$ lies below $s$ precisely when the sign
of $s$ at position $\alpha$ is plus. Thus
\begin{equation}
 \{\alpha<\delta:\Bit(s,\alpha)\},\qquad\bd(s)=\delta,
 \label{hset:eq:fiber-subsets}
\end{equation}
runs through \emph{all} subsets of $\delta$, once each, for every
$\delta<\lambda$. This uses a first-order quantifier over numbers of one
birthday, not a second-order quantifier in the language. The exact length
matters: trailing minus signs are not discarded, so the code has an explicit
domain.

\subsection{Pairing without reconstructing ordinary ordinal multiplication}

For a positive embedded ordinal $\theta$, put
\begin{equation}
 p_\theta(\xi,\zeta)=\theta\xi+\zeta,
 \qquad \Lambda(\theta)=\theta^2+\theta.
 \label{hset:eq:pairing}
\end{equation}
All operations in \eqref{hset:eq:pairing} are field operations, hence natural
ordinal operations on these arguments.

\begin{lemma}[Bounded pairing]
\label{hset:lem:pairing}
The map $p_\theta$ injects $\theta\times\theta$ into $\Lambda(\theta)$.
For every epsilon number $\lambda$ and every $0<\theta<\lambda$,
$\Lambda(\theta)<\lambda$. For infinite $\theta$,
$\card{\Lambda(\theta)}=\card\theta$.
\end{lemma}
\begin{proof}
Suppose $\xi<\xi'<\theta$ and $\zeta,\zeta'<\theta$. Since
$\xi+1\le\xi'$ in ordinal order, positivity and distributivity give
\[
 \theta\xi+\zeta<\theta\xi+\theta
 =\theta(\xi+1)\le\theta\xi'\le\theta\xi'+\zeta'.
\]
For fixed $\xi$, cancellation gives injectivity in $\zeta$. Also
$\theta\xi+\zeta<\theta^2+\theta$. Closure below $\lambda$ follows from its
being a field cutoff. The cardinal estimate is
\cref{hset:lem:card-estimates}, together with $\theta\le\Lambda(\theta)$.
\end{proof}

Source 02 uses the same pairing and bound (its $B_\theta$). Consequently a sign
sequence of birthday $\Lambda(\theta)$ codes a binary relation on $\theta$. To
make this code canonical for a given relation, require every plus position to
be in the image of $p_\theta$; unused positions are minus. This last
requirement is the first-order formula
\[
 \forall\rho<\Lambda(\theta)\,
 \bigl(\Bit(r,\rho)\longrightarrow
       \exists\xi,\zeta<\theta\ [\rho=p_\theta(\xi,\zeta)]\bigr).
\]
Here and below bounded ordinal quantifiers include the predicate $\Ord$.

For a graph between two different ordinals $\delta,\epsilon<\lambda$, use
$\theta=\delta+\epsilon+1$, which is above both. Coding the graph in the
$\delta\times\epsilon$ rectangle of this square again stays below $\lambda$.
The construction requires no definable global pairing on the entire class
of ordinals, and no reconstruction of ordinary ordinal arithmetic.

\begin{remark}[Four pairings in the sources]\label{hset:rem:pairings}
The other sources pack relations differently; each variant works with the codes
below after the obvious changes. Source 05 uses the rectangle map
$p_a(i,j)=aj+i$ for $i<a$, $j<d$, with length $a(d+1)$, which is the same idea
with the roles of the coordinates exchanged. Source 06 uses the polynomial of
\cref{hset:lem:pairing06}, and source 07 the radix pairing of
\cref{hset:lem:radix}.
\end{remark}

\begin{lemma}[A one-polynomial pairing; source 06]\label{hset:lem:pairing06}
On ordinal surreals put $\varpi(\alpha,\gamma)=(\alpha+\gamma)^2+\alpha$. Then
$\varpi$ is injective on pairs of ordinals, increasing in each coordinate, and
maps $\kappa\times\kappa$ into $\kappa$ for every uncountable cardinal
$\kappa$.
\end{lemma}
\begin{proof}
Put $\sigma=\alpha+\gamma$. Since $0\le\alpha\le\sigma$,
$\sigma^2\le\varpi(\alpha,\gamma)\le\sigma^2+\sigma$. If $\sigma<\sigma'$ are
ordinals, then $\sigma+1\le\sigma'$ and $\sigma^2+\sigma<(\sigma+1)^2\le
\sigma'^2$. The intervals belonging to different sums are therefore disjoint.
Equal values force equal sums, then equal first coordinates by subtraction,
and finally equal second coordinates. Monotonicity follows from positivity and
monotonicity of addition and squaring on nonnegative elements, and the bound
is \cref{hset:lem:card-estimates}.
\end{proof}

\begin{lemma}[Admissible radices; source 07]\label{hset:lem:radix}
Let $\operatorname{Rad}(q)$ say that $q$ is an ordinal, $q>1$, and $a+a'<q$ for
all ordinals $a,a'<q$. The ordinals satisfying $\operatorname{Rad}$ are exactly
$q=\omega^\rho$ with $\rho>0$; for each of them $p_q(a,a')=qa+a'$ injects
$q\times q$ into $q^2$. Both statements are absolute between $\B$ and every
initial cutoff containing the ordinals in question. Below an epsilon number
$\lambda$ the admissible radices are cofinal and $q^2<\lambda$ for each of them.
\end{lemma}
\begin{proof}
For $q=\omega^\rho$ every ordinal below $q$ has Cantor normal form with
exponents below $\rho$, and natural addition introduces no new exponent. If
$q>1$ is not a single monomial with coefficient one, split its normal form into
two nonzero proper parts (splitting a coefficient if needed) to get $a,a'<q$
with $a\nsum a'=q$; $q=1$ is excluded. For $a<q$ the exponents of $q\nprod a$
are all at least $\rho$, those of $a'<q$ are below $\rho$, and cancellation of
natural addition recovers $a$ and $a'$; moreover $qa+a'<q(a+1)\le q^2$. For
$\xi<\lambda$ the radix $\omega^{\xi+1}$ exceeds $\xi$ and lies below
$\lambda$, and $(\omega^\rho)^2=\omega^{\rho\nsum\rho}<\lambda$ since
$\lambda=\omega^\lambda$ is closed under natural addition below itself.
All ordinals below an argument of an initial cutoff are present there, so the
quantified test is unchanged.
\end{proof}

\section{Pointed relation codes and the exact interpreted universe}
\label{hset:sec:codes}

\subsection{The definable domain of codes}

A prospective code is a triple
\[
 C=(\delta,r,t),\qquad
 0<\delta<\lambda,\quad t<\delta,\quad\bd(r)=\Lambda(\delta),
\]
with $\delta,t$ embedded ordinals and with unused bits of $r$ equal to minus.
Write
\[
 E_C(i,j)\quad\Longleftrightarrow\quad
 i,j<\delta\ \land\ \Bit(r,p_\delta(i,j)).
\]
The orientation is \emph{child belongs to parent}: $i\,E_C\,j$ will mean
that the object at node $i$ is a member of the object at node $j$.

Require extensionality,
\begin{equation}
 \forall i,j<\delta\,
 \bigl[\forall k<\delta\ (E_C(k,i)\leftrightarrow E_C(k,j))
       \longrightarrow i=j\bigr],
 \label{hset:eq:extensional}
\end{equation}
and well-foundedness in the following form:
\begin{align}
 \forall u\,[\bd(u)=\delta\ \longrightarrow\ (&
 (\exists j<\delta\ \Bit(u,j))\ \longrightarrow\nonumber\\[-0.2em]
 &\exists j<\delta\,[\Bit(u,j)\land
    \forall i<\delta\ (\Bit(u,i)\longrightarrow\neg E_C(i,j))])].
 \label{hset:eq:wf}
\end{align}
By \eqref{hset:eq:fiber-subsets}, this is \emph{actual} well-foundedness of the
set relation, not just well-foundedness against a restricted collection of
subsets.

Finally require generation by the root:
\begin{align}
 \forall u\,[\bd(u)=\delta\ \longrightarrow\ (&
 \Bit(u,t)\land
 \forall i,j<\delta\,[\Bit(u,j)\land E_C(i,j)\longrightarrow\Bit(u,i)]
 \nonumber\\[-0.2em]
 &\longrightarrow\forall i<\delta\ \Bit(u,i))].
 \label{hset:eq:rooted}
\end{align}
Thus no proper downward-closed subset of the graph contains the root.
Let $\Code(C)$ be the conjunction of these conditions. The triple is an
ordinary finite-dimensional interpretation domain in the single-sorted
structure $\B_\lambda$. Source 02 uses four-tuples with an extra domain word
that selects a subset of $\theta$ as nodes, and source 07 four-tuples with a
radix; in both, a ``clean'' or ``no junk'' condition plays the role of the
unused-bit convention.

\begin{lemma}[Rooted collapse]
\label{hset:lem:collapse}
Every valid code has a unique Mostowski collapse
$\pi_C:\delta\longrightarrow T_C$ onto a transitive set. If
$x_C=\pi_C(t)$, then
\[
 T_C=\TC(\{x_C\}),\qquad
 \card\delta=\card{\TC(\{x_C\})}.
\]
\end{lemma}
\begin{proof}
Define by well-founded recursion
\[
 \pi_C(j)=\{\pi_C(i):i\,E_C\,j\}.
\]
Extensionality and well-founded induction show that $\pi_C$ is injective:
if a collision existed, choose one minimizing the maximum of the two node
ranks; comparison of predecessor values gives equal predecessor sets, and
extensionality identifies the nodes. Its image is transitive, and the
displayed recursion gives an isomorphism. Uniqueness follows by well-founded
induction as well; this is the usual Mostowski-collapse argument.

The nodes reachable from $t$ by finitely many downward edges form a
downward-closed subset containing $t$. By \eqref{hset:eq:rooted} this is all of
$\delta$. Under the collapse, these are exactly $x_C$ and its iterated
members. Thus the image is $\TC(\{x_C\})$, giving the cardinal equality.
\end{proof}

Root generation is not needed merely to encode arbitrary sets. It is useful
here because it makes the hereditary-size cost exact rather than just an
upper bound polluted by unrelated graph components.

\subsection{Equality and membership are definable}

For codes $C=(\delta,r,t)$ and $D=(\epsilon,s,u)$, define $C\sim D$ when a
coded bijection $f:\delta\to\epsilon$ preserves $E$, sends $t$ to $u$, and
reflects $E$. Such a graph is expressed by a sign sequence in a square of
width $\delta+\epsilon+1$, as in \cref{hset:lem:pairing}. The familiar function,
bijection, and edge-preservation clauses use only bounded ordinal
quantifiers and bits. Therefore $\sim$ is a single first-order relation.

For membership, let $C\in_\I D$ mean that there is a coded injection
$f:\delta\to\epsilon$ such that
\begin{align}
 &E_D(f(t),u),\label{hset:eq:member-root}\\
 &E_C(i,j)\ \longleftrightarrow\ E_D(f(i),f(j))
       &&(i,j<\delta),\label{hset:eq:member-edge}\\
 &E_D(v,f(j))\ \longrightarrow\
       \exists i<\delta\,[E_C(i,j)\land f(i)=v]
       &&(j<\delta,0\le v<\epsilon).\label{hset:eq:member-full}
\end{align}
In \eqref{hset:eq:member-full}, the quantifier on $v$ simply ranges over the
embedded ordinals below $\epsilon$. The clause says that the image contains
\emph{all} predecessors of its nodes, not merely some of them. Sources 06 and
07 use the same downward-closed image clause; source 02 instead defines
membership through an isomorphism of the root with the cone of a predecessor,
using a definable reachability relation.

\begin{lemma}[Correctness of equality and membership]
\label{hset:lem:equal-member}
For valid codes,
\[
 C\sim D\iff x_C=x_D,
 \qquad
 C\in_\I D\iff x_C\in x_D.
\]
\end{lemma}
\begin{proof}
A root-preserving isomorphism gives equal collapsed roots by uniqueness of
collapse. Conversely, if the roots agree, the two images are the same
transitive closure by \cref{hset:lem:collapse}, and
$\pi_D^{-1}\pi_C$ gives the required isomorphism. Its graph has size at most
$\max(\card\delta,\card\epsilon)$ and is codable within the same cutoff.

For membership, \eqref{hset:eq:member-edge} and \eqref{hset:eq:member-full},
followed by induction along $E_C$, imply $\pi_D(f(j))=\pi_C(j)$ for all $j$.
The root condition then gives $x_C\in x_D$. Conversely, if $x_C\in x_D$, the
inclusion $\TC(\{x_C\})\subseteq\TC(\{x_D\})$ transports through the collapse
maps to an injection satisfying all three clauses. Its image is
predecessor-closed, so \eqref{hset:eq:member-full} really does hold. The graph
is again small enough to be coded below the cutoff.
\end{proof}

The predecessor-full clause cannot be omitted; Appendix~\ref{hset:app:regression}
gives the counterexample, which the finite suites of sources 06, 07 and 09
check.

Let $\I(\B_\lambda)$ be the quotient of valid codes by $\sim$, with the
membership relation just defined. All quantifiers in its interpreted
formulas are replaced by quantifiers over triples satisfying $\Code$.
No quotient class is assumed to be an element of the original structure.

\subsection{Cardinal rounding}

\begin{theorem}[Exact collapse range]
\label{hset:thm:collapse-range}
For every epsilon number $\lambda$, the map $[C]\mapsto x_C$ is an
isomorphism
\[
 \I(\B_\lambda)\cong(H_{\ceilcard{\lambda}},\in).
\]
More precisely, the ordinals $\delta$ supporting rooted codes for a given
set $x$ are exactly those with
\begin{equation}
 0<\delta<\lambda,
 \qquad\card\delta=\card{\TC(\{x\})}.
 \label{hset:eq:exact-code-size}
\end{equation}
\end{theorem}
\begin{proof}
Every code has $\delta<\lambda\le\kappa(\lambda)$. Hence its collapse has
hereditary size $\card\delta<\kappa(\lambda)$, by
\cref{hset:lem:collapse}. This proves one inclusion of the range.

Conversely, suppose $x\in H_{\kappa(\lambda)}$. Its nonempty transitive
closure has some cardinal $\nu<\kappa(\lambda)$. By the definition of
cardinal ceiling, $\nu<\lambda$: if $\lambda$ is a cardinal this is
immediate; otherwise $\nu\le\card\lambda<\lambda$. Choose a bijection from
$\nu$ onto $\TC(\{x\})$, transfer membership to $\nu$, and designate the
node for $x$. The resulting graph is extensional, well founded, and generated
by that node. It has a sign code of length $\Lambda(\nu)<\lambda$.
Thus $x$ occurs in the range.

For any other $\delta<\lambda$ equipotent to this transitive closure, the
same construction works with a bijection from $\delta$. Necessity in
\eqref{hset:eq:exact-code-size} was already proved by rooted collapse.
Correctness of both relations is \cref{hset:lem:equal-member}.
\end{proof}

Choice is used here only to well-order one transitive closure at a time; the
interpretation keeps all valid codes and never selects one code per set.
Source 02 proves this theorem one way, at every epsilon number, as a
four-dimensional interpretation; source 07 proves the same range at every
epsilon number (its cardinality plateaux, with $\theta(\lambda)$ in place of
$\kappa(\lambda)$), also stating it forward only; sources 05 and 06 prove it at
regular, respectively all uncountable, cardinals.

\begin{remark}[Why the bound is sharp; source 02]\label{hset:rem:sharp}
No set of hereditary size $\kappa(\lambda)$ can be represented by these codes:
its collapse would require at least $\kappa(\lambda)$ distinct nodes, while
each domain is an ordinal below $\lambda\le\kappa(\lambda)$. This obstruction
is intrinsic to the stated interpretation. It is not a claim that no
different interpretation in $\B_\lambda$ could represent a larger structure.
\end{remark}

\begin{example}[The empty set and a first nonempty code]
\label{hset:ex:first-codes}
The empty set uses one node, no edges, and root zero. With $\delta=1$, the
canonical relation string has birthday $\Lambda(1)=2$ and both signs minus.
Thus a code for the empty set is $(1,--,0)$, not the surreal zero itself.
For $\{\varnothing\}$ take $\delta=2$, root $1$, and the single edge
$0E1$. The relation string has length $6$ and its only plus is at
$p_2(0,1)=1$. Numeric values and the sets they code must not be identified.
A three-node example with two labelings is in Appendix~\ref{hset:app:example}.
\end{example}

\begin{corollary}[Countable cutoffs recover all countable sets; sources 02, 07, 09]
\label{hset:cor:countable}
For every countable epsilon number $\lambda$, in particular $\varepsilon_0$,
$\I(\B_\lambda)\cong(H_{\omega_1},\in)$.
\end{corollary}
\begin{proof}
A countable epsilon number is strictly larger than $\omega$ and strictly below
$\omega_1$, so its cardinal ceiling is $\omega_1$. Concretely, every
hereditarily countable set has a rooted code on a finite ordinal or on
$\omega$, and $\Lambda(\omega)=\omega^2+\omega<\varepsilon_0$.
\end{proof}

The recovered sets include countable ordinals far larger than $\varepsilon_0$.
Those larger ordinals are recovered as \emph{collapsed relations}, not as
all-plus scalar sign sequences in the cutoff: every countable ordinal $\alpha$
has a membership graph labeled by natural numbers, whose relation word has
birthday below $\omega^2+\omega$. Labels need not preserve ordinal order; the
collapse recovers that order from the edges. One should therefore not infer a
low set-theoretic rank of the interpreted objects from a low birthday of the
numerical codes (source 07).

\section{Bi-interpretation and recovery of the cutoff}
\label{hset:sec:biinterpretation}

\subsection{The ordinal and sign-graph bridges}

For an actual embedded ordinal $a<\lambda$, there is a distinguished code
$\OC(a)$ with domain $a+1$, root $a$, and edge relation $iEj$ iff $i<j$.
Its sign relation code is uniquely determined by the unused-bit convention.
The formula defining $\OC(a)$ is uniform, and its collapse is the
von Neumann ordinal $a$. Because $\lambda$ is a limit and a ring cutoff, both
$a+1$ and $\Lambda(a+1)$ are below $\lambda$.

\begin{definition}[Visible interpreted ordinals]
For a set code $C$ define
\[
 \Vis(C)\quad\Longleftrightarrow\quad
 \exists a\,[\Ord(a)\land C\sim\OC(a)].
\]
Under the collapse isomorphism, these are exactly the ordinals below
$\lambda$, not necessarily all ordinals of $H_{\kappa(\lambda)}$.
\end{definition}

There is also a definable bridge $\SG(s,C)$ saying that $C$ represents the
set-coded sign function of the scalar $s$: its domain is $\OC(\bd(s))$, its
values are the interpreted ordinals $0$ and $1$, and at each actual
$\alpha<\bd(s)$ it has value $1$ exactly when $\Bit(s,\alpha)$ holds.
Application of a set-coded function and the Kuratowski ordered-pair
construction are ordinary fixed formulas in the interpreted membership
language. The ordinal bridge supplies their arguments. Thus $\SG$ is a
fixed first-order formula, not an external assignment operation.

For each $s$ there is exactly one equivalence class of codes $C$ with
$\SG(s,C)$; conversely, every sign function whose length is below $\lambda$
arises this way. Existence follows from the hereditary-size estimate for
sign graphs and \cref{hset:thm:collapse-range}.

\subsection{A first-order test for a cardinal cutoff}

Define the \emph{visible-ordinal collector} formula
\begin{equation}
 \Coll(C)\quad\Longleftrightarrow\quad
 \Code(C)\land
 \forall D\,[\Code(D)\longrightarrow
       (D\in_\I C\leftrightarrow\Vis(D))].
 \label{hset:eq:collector}
\end{equation}
It says that one interpreted set collects all the scalar ordinals. This is
not the assertion that the scalar carrier itself is a set inside itself.

\begin{theorem}[The cutoff is a set exactly in the noncardinal case; sources 02, 09]
\label{hset:thm:collector}
For an epsilon number $\lambda$,
\[
 \B_\lambda\models\exists C\,\Coll(C)
 \quad\Longleftrightarrow\quad
 \lambda\text{ is not a cardinal}.
\]
When it exists, the collector is unique modulo $\sim$ and its collapse is
exactly the ordinal $\lambda$.
\end{theorem}
\begin{proof}
The elements specified in \eqref{hset:eq:collector} are exactly all ordinals
$a<\lambda$. In the ambient set universe the unique set with precisely
those elements is $\lambda$. Hence a collector exists precisely when
$\lambda\in H_{\kappa(\lambda)}$, equivalently
$\card\lambda<\kappa(\lambda)$, since $\TC(\{\lambda\})=\lambda+1$. For a
cardinal $\lambda$ the right side fails; otherwise
$\kappa(\lambda)=(\card\lambda)^+$ and it holds. Uniqueness follows from
extensionality of the interpreted universe.
\end{proof}

Source 02's sentence $\mathsf{CollectOrd}$ is this sentence written with its
own codes. The theorem recovers the \emph{ordinal value} of a noncardinal
cutoff as an interpreted object. It is not a construction of a surreal number
with birthday $\lambda$ inside $\No_{<\lambda}$. This collector is the
additional information needed to recover a noncardinal cutoff after
interpreting its hereditary-size universe. Merely saying that $\B_\lambda$
interprets $H_{\kappa(\lambda)}$ would lose it.

\subsection{The reverse interpretation}

If $\lambda=\kappa$ is a cardinal, interpret surreals in $H_\kappa$ as all
set-coded functions from an ordinal into $\{0,1\}$. By
\cref{hset:lem:hereditary}, these are exactly the sign sequences of length
below $\kappa$. Define the order by first disagreement, birthday by the
all-one function on the domain ordinal, and arithmetic by the certificate
formulas of \cref{hset:thm:certificates} (or the tables of
\cref{hset:thm:tables}). Call this interpretation $\J$.

If $\lambda$ is not a cardinal, work instead in the pointed structure
$(H_{\kappa(\lambda)},\in,\lambda)$. Use the same definitions, with the
extra domain restriction that the sign length is below the named ordinal
$\lambda$. The certificate verifier may use auxiliary sign sequences
anywhere in $H_{\kappa(\lambda)}$; it is not required to carry out a large
recursion inside the smaller interpreted sign domain. The root outputs are
below $\lambda$ by classical cutoff closure. This gives a uniform
interpretation $\J_{\rm cut}$ in the language with one named cutoff.
Without the named cutoff, interpreting sign sequences in $H_{\omega_1}$ gives
$\No_{<\omega_1}$, not $\No_{<\varepsilon_0}$; this is why source 07 stated its
plateau theorem forward only.

\begin{theorem}[Uniform local bi-interpretation]
\label{hset:thm:biinterpretation}
The interpretations just constructed prove all assertions of
\cref{hset:thm:main-rounding}. The maps witnessing the two bi-interpretations
are definable, not merely external bijections.
\end{theorem}
\begin{proof}
Start on the set side. A code in the interpreted surreal structure gives
an extensional well-founded relation on an ordinal. Its Mostowski-collapse
root defines the map back to the original set. The graph of the collapse
belongs to the same $H_\kappa$: its domain, range, and all values have
hereditary size bounded by a fixed infinite cardinal below $\kappa$.
Consequently this correspondence is defined by a first-order formula in
$H_\kappa$, even when Replacement is not an axiom of that structure.
The collapse equation is a bounded verification of the supplied function.
Every set has a code, and equality and membership are respected by
\cref{hset:lem:equal-member}. In the pointed case, the collector collapses to
the named cutoff, so the correspondence preserves that constant as well.

Start on the surreal side. The formula $\SG(s,C)$ identifies $s$ with the
sign function representing it in the copy of the set universe. In the
cardinal case all sign functions in the interpreted $H_\lambda$ occur.
In the noncardinal case, the collector defines $\lambda$, and the restricted
reverse interpretation uses exactly lengths below that ordinal. Thus $\SG$
is a definable bijection to the twice-interpreted surreal structure.
It preserves numerical order and birthday directly from the sign rules,
and it preserves arithmetic by the certificate theorem. This proves the
required definable isomorphisms in both directions.
\end{proof}

At cardinal cutoffs sources 06 and 07 give the same two round trips with
their own codes (06 by a formula $\mathsf{Eval}$ for the collapse and a
sign-function bridge; 07 by an ordinal-code comparison and the interpreted sign
name $\langle\bd(x),A_x\rangle$), and source 05 gives them with graphs modulo
bisimulation (\cref{hset:sec:bisim}).

\begin{corollary}[The first epsilon cutoff and hereditary countability]
\label{hset:cor:epsilon0}
The structure $\B_{\varepsilon_0}$ is parameter-free bi-interpretable with
$(H_{\omega_1},\in)$.
\end{corollary}
\begin{proof}
The theorem first gives the pointed structure
$(H_{\omega_1},\in,\varepsilon_0)$. The ordinal $\varepsilon_0$ is
parameter-free definable in $H_{\omega_1}$ as the least nonzero fixed point
of ordinary $\alpha\mapsto\omega^\alpha$. For any countable
arguments, the functions witnessing ordinal arithmetic and the relevant
well-orders are countable sets, hence belong to $H_{\omega_1}$. Their
well-founded recursive verification is absolute. The least-fixed-point
specification therefore picks out the actual $\varepsilon_0$ there.
Naming it is a definitional expansion, which may be removed.
\end{proof}

The same argument applies to any noncardinal epsilon number that is
parameter-free definable, with its correct ordinal value, in the recovered
hereditary-size universe. It is not an assertion that every ordinal can be
named parameter-free.

\begin{remark}[Questions of the sources answered here]\label{hset:rem:answered-bi}
Source 02 claimed only a one-way interpretation and warned that a two-way
theorem ``should specify whether the target carries a predicate for that
segment or a constant for $\lambda$''. \Cref{hset:thm:biinterpretation}, source
09's, is exactly such a theorem, with a constant; at $\varepsilon_0$ it is even
unpointed. Source 06's question on noncardinal cutoffs, source 07's question on
recovering a noncardinal cutoff, and the noncardinal half of source 05's
question on optimal cutoffs are answered for epsilon-number cutoffs by the same
theorem; the singular-cardinal half of 05's question is answered by sources 06
and 07. The questions stay open for cutoffs that are not epsilon numbers
(\cref{hset:sec:questions}).
\end{remark}

\subsection{Coherence of the interpretation}

\begin{proposition}[Coherence under cutoff inclusions]
\label{hset:prop:coherence}
If $\lambda\le\eta$ are epsilon numbers, the inclusion of scalar cutoffs
induces, under $\I$, the actual inclusion
\[
 H_{\kappa(\lambda)}\longrightarrow H_{\kappa(\eta)}.
\]
A code keeps the same collapsed value, and all sign-graph bridges agree.
For cardinal cutoffs the reverse interpretations also commute with the
actual inclusions.
\end{proposition}
\begin{proof}
For a code already below $\lambda$, the larger structure has no additional
sign sequences of birthday $\delta$ for any $\delta<\lambda$: both contain
\emph{all} sequences of each such length in the same ambient universe.
Thus the subset tests in well-foundedness and root generation are unchanged.
Equality and membership agree by their collapse characterizations; any
needed witness graph has a code below $\lambda$ as proved above. The
collapse is defined from the same relation, so its value is unchanged.
The sign-graph bridge refers to the same prefixes and bits. Finally the
reverse arithmetic formulas compute the same external operations.
\end{proof}

Sources 02 (at epsilon numbers, forward), 06 and 07 (at cardinal cutoffs,
both directions) prove the same compatibility. If two noncardinal cutoffs have
the same cardinal ceiling, the induced inclusion on recovered
\emph{unpointed} set universes is the identity. Their recovered distinguished
cutoffs nevertheless differ. This is precisely why elementary inclusion will
usually fail.

\section{Second route: graphs modulo bisimulation}
\label{hset:sec:bisim}

Source 05 codes sets by well-founded graphs that need not be extensional or
generated by their root, and takes \emph{bisimilarity} as equality. This avoids
any extensionality or canonical-enumeration requirement: two different
terminal nodes both represent $\varnothing$, and a disjoint component below no
root changes nothing. We write its codes with the pairing of
\cref{hset:lem:pairing}; source 05 uses its rectangle map
(\cref{hset:rem:pairings}), which changes nothing below. Source 05 proved the
resulting bi-interpretation at
uncountable regular cardinals; the argument below uses regularity nowhere (the
certificates of \cref{hset:sec:certificates} are the only arithmetic input), and
we state it at the cardinal cutoffs of \cref{hset:thm:main-rounding}, where the
base already proves the conclusion by the first route. The value of the second
route is its different equality test.

A \emph{graph code} is a triple $G=(\nu,e,t)$ with $\nu>0$ and $t<\nu$ embedded
ordinals and $\bd(e)=\Lambda(\nu)$; its relation is
$E_G(i,j)\iff i,j<\nu\land\Bit(e,p_\nu(i,j))$, again oriented from a member to
a set. Let $\operatorname{WF}(G)$ be the well-foundedness clause
\eqref{hset:eq:wf}. No extensionality or root generation is required. The
\emph{decoration} of $G$ is the unique function with
\[
 v_G(j)=\{v_G(i):i\,E_G\,j\},\qquad\operatorname{dec}(G)=v_G(t);
\]
ordinary well-founded recursion suffices, and the Mostowski collapse is its
extensional special case.

\begin{proposition}[Exactly the hereditary sets; source 05]
\label{hset:prop:bisim-range}
For an uncountable cardinal $\kappa$, the decorations of graph codes in
$\B_\kappa$ are exactly the members of $H_\kappa$.
\end{proposition}
\begin{proof}
Every member of $\TC(\{v_G(j)\})$ is a value at a node reachable by a finite
predecessor chain, so it lies in the range of $v_G$, of cardinality at most
$\card\nu<\kappa$. Conversely, pull membership on a transitive set containing
$a\in H_\kappa$ back to an ordinal $\nu<\kappa$ along a bijection and root it at
the label of $a$; induction gives decoration $a$.
\end{proof}

For graph codes $G=(\nu,e,t)$ and $G'=(\nu',e',t')$, a relation
$Q\subseteq\nu\times\nu'$ is a \emph{bisimulation} when every $(i,j)\in Q$
satisfies
\begin{align*}
 i'E_Gi&\Longrightarrow\exists j'\,(j'E_{G'}j\land(i',j')\in Q),\\
 j'E_{G'}j&\Longrightarrow\exists i'\,(i'E_Gi\land(i',j')\in Q).
\end{align*}
Code $Q$ by one sign sequence $q$ in a square of width $\nu+\nu'+1$, and let
$\operatorname{Bis}(G,G',q)$ be the resulting first-order formula. Define
\[
 G\approx G'\iff\exists q\,[\operatorname{Bis}(G,G',q)\land (t,t')\in Q_q],
 \qquad
 G\in_{\approx}G'\iff\exists j<\nu'\,[E_{G'}(j,t')\land G\approx G'[j]],
\]
where $G'[j]=(\nu',e',j)$ is the same graph rerooted at $j$.

\begin{theorem}[Equality is bisimilarity; source 05]\label{hset:thm:bisim}
For well-founded graph codes,
$G\approx G'\iff\operatorname{dec}(G)=\operatorname{dec}(G')$ and
$G\in_\approx G'\iff\operatorname{dec}(G)\in\operatorname{dec}(G')$.
Consequently the graph codes modulo $\approx$, with $\in_\approx$, form a
parameter-free interpretation of $(H_\kappa,\in)$ in $\B_\kappa$.
\end{theorem}
\begin{proof}
Suppose $Q$ is a bisimulation relating the roots. By well-founded induction on
$i$ along $E_G$, every $j$ with $(i,j)\in Q$ satisfies $v_G(i)=v_{G'}(j)$: a
member $v_G(i')$ of $v_G(i)$ is matched by the forth clause with a
predecessor $j'$ of $j$ and equals $v_{G'}(j')$ by induction, and the back
clause gives the other inclusion. Conversely, if the root values agree, the set
relation $Q=\{(i,j):v_G(i)=v_{G'}(j)\}$ is a bisimulation relating the roots,
because equal decorated sets supply the forth and back witnesses without
selecting unique matches; it has a code in the same cutoff. Membership
follows by applying this to the rerooted graph. Hence $\approx$ is an
equivalence relation, a conclusion established semantically rather than
presumed from the shape of the definition, and \cref{hset:prop:bisim-range}
gives surjectivity.
\end{proof}

\begin{example}[Duplicates and irrelevant components]
One terminal node represents $\varnothing$. A graph with two terminal nodes and
a root above both represents $\{\varnothing\}$, as does a graph with only one
terminal child. Their rooted graphs are not isomorphic, but their roots are
bisimilar. Graph isomorphism is therefore \emph{not} the right equality for
these codes; source 05's finite suite exercises exactly this case.
\end{example}

The round trips of source 05 are the analogues of those in
\cref{hset:thm:biinterpretation}. The set-side map is the decoration relation
``there is a function $v$ on $\nu$ satisfying the recursion with $v(t)=a$'',
whose witness has hereditary size at most $\max(\aleph_0,\card\nu)$. On the
surreal side, for an ordinal $\alpha$ let $\OC(\alpha)$ be the graph on
$\alpha+1$ with $iEj\iff i<j$ and root $\alpha$, so that
$G\approx\OC(\alpha)$ iff $\operatorname{dec}(G)=\alpha$; and for a surreal
$x$ let $A_x$ be the interpreted set whose members are the classes
$\OC(\gamma)$ for the ordinals $\gamma<\bd(x)$ with $\Bit(x,\gamma)$. The
number $x$ is then matched with the interpreted sign name
$\langle\bd(x),A_x\rangle$, which externally decodes to
$\langle\bd(x),\{\gamma:x(\gamma)=+\}\rangle$. Source 07 uses the same
sign-name round trip with extensional codes.

\begin{remark}[Internal versions and models of ZFC; sources 05, 07]
\label{hset:rem:internal}
All constructions above are theorems of ZFC and can be carried out inside any
model $M\models\ZFC$, with $M$'s own ordinals, subsets and well-founded
relations; for the full class $\No^M$ they recover $M$ itself. For a
nonstandard $M$ this is an internal statement: graph well-foundedness is $M$'s
notion, not external well-foundedness. For set models $M,N$ of ZFC, writing
$\B^M$ and $\B^N$ for the internal birthday structures regarded externally,
source 07 derives
\[
 M\equiv N\iff\B^M\equiv\B^N,\qquad M\cong N\iff\B^M\cong\B^N,
\]
by transporting sentences and isomorphisms through the uniform
interpretations and their definable composites. These full-class and internal
statements are instances of the announced Chen--Hamkins--Yang theorem
(\cref{hset:thm:chy}), not results of sources 05 or 07.
\end{remark}

\section{Ring cutoffs and the finite boundary}
\label{hset:sec:rings}

The classical ring statement is weaker than the field statement: the infinite
ring cutoffs $\No_{<\lambda}$ are exactly those with $\lambda=\omega^{\omega^\alpha}$
\cite[Theorems 1.1--1.2]{BG}. Source 02 observes that the forward
interpretation needs only ring operations.

\begin{proposition}[Ring version; source 02]\label{hset:prop:ring}
The interpretation $\I$ and its exact image remain valid for every infinite
ordinal $\lambda$ for which $\No_{<\lambda}$ is a subring of $\No$, when
$\B_\lambda$ is regarded as a birthday-enriched ordered ring. Its image is
$H_{\kappa(\lambda)}$, where now $\kappa(\lambda)$ is the least \emph{infinite}
cardinal $\ge\lambda$. In particular
\[
 \I(\No_{<\omega};+,-,\cdot,<,\bd)\cong(H_\omega,\in),
\]
the hereditarily finite sets, and $\I(\B_{\omega^\omega})\cong(H_{\omega_1},\in)$.
\end{proposition}
\begin{proof}
Every formula of the interpretation uses only the ordered-ring operations and
birthday. The bounds $\Lambda(\theta)$ and $\theta+\epsilon+1$ remain in the
cutoff by ring closure, and an infinite ring cutoff is a limit ordinal, so the
ordinal codes also fit. All later coding and collapse arguments use only these
facts and the cardinal-ceiling calculation. For $\lambda=\omega$ every domain is
finite, all codes for its subsets and finite isomorphisms have finite length,
and the represented transitive closures are exactly the finite ones.
\end{proof}

Real closedness matters for the comparison with decidable field theories and
for the surcomplex transfer, but it is not what enables the set coding. The
decisive ingredients are complete sign layers, definable reading of those
signs, and enough ordinal arithmetic to fit relation codes below the cutoff.
Only the forward interpretation is claimed at ring cutoffs that are not epsilon
numbers. At $\omega$ the cutoff is the dyadic ring, not a field, and it would
be incorrect to reuse any real-closed-field argument there. Arbitrary ordinal
cutoffs that are not ring cutoffs are not classified.

\part{Consequences inside one universe}\label{hset:part:inside}

\section{The exact elementary-inclusion spectrum}
\label{hset:sec:elementary}

The elementary-inclusion result is stronger than the observation that a
birthday expansion is not governed by real-closed-field quantifier
elimination. It identifies exactly which initial inclusions are possible.

\begin{lemma}[Noncardinal lower cutoffs never include elementarily; sources 02, 09]
\label{hset:lem:lower-noncardinal}
If $\lambda<\eta$ are epsilon numbers and $\lambda$ is not a cardinal, then
$\B_\lambda\not\prec\B_\eta$.
\end{lemma}
\begin{proof}
Choose a code $C$ in $\B_\lambda$ for its visible-ordinal collector. Its
collapse is $\lambda$, and it satisfies $\Coll(C)$ there. By coherence it
still collapses to $\lambda$ in $\B_\eta$. But $\lambda$ itself is now a
scalar ordinal, since $\lambda<\eta$. Its interpreted ordinal code is
visible, and is not a member of the set coded by $C$, because
$\lambda\notin\lambda$. Thus $\Coll(C)$ is false in the larger structure.
This single formula, with the three entries of $C$ as parameters, witnesses
failure of elementarity.
\end{proof}

Source 02's witness is the weaker formula ``the coded set contains every
literal ordinal'', which is also true in $\B_\lambda$ and false in $\B_\eta$
for the same code. As source 02 observes, $\B_{\varepsilon_0}$ and
$\B_{\varepsilon_1}$ interpret the same $H_{\omega_1}$, yet their inclusion is
not elementary: interpretation does not discard the additional first-order
information of the original structure.

\begin{lemma}[A cardinal lower cutoff cannot include elementarily into a
noncardinal cutoff]
\label{hset:lem:upper-noncardinal}
If $\lambda$ is a cardinal and $\eta$ is not, then
$\B_\lambda\not\equiv\B_\eta$.
\end{lemma}
\begin{proof}
The parameter-free sentence $\exists C\,\Coll(C)$ is false in the first
structure and true in the second, by \cref{hset:thm:collector}.
\end{proof}

\begin{theorem}[Classification of elementary initial inclusions]
\label{hset:thm:elementary}
For epsilon numbers $\lambda<\eta$,
\[
 \B_\lambda\prec\B_\eta
 \quad\Longleftrightarrow\quad
 \lambda,\eta\text{ are cardinals and }H_\lambda\prec H_\eta.
\]
\end{theorem}
\begin{proof}
The preceding two lemmas force both cutoffs to be cardinals; they are
uncountable, since no countable cardinal is an epsilon number. For cardinal
cutoffs, an elementary inclusion of the surreal structures induces an
elementary inclusion of their interpreted set structures. By
\cref{hset:prop:coherence} this is the actual inclusion
$H_\lambda\subseteq H_\eta$. Conversely, an elementary inclusion of those set
structures induces an elementary inclusion of the uniformly interpreted sign
structures. The reverse interpretation and its definable sign-graph
isomorphism identify this with the original inclusion
$\B_\lambda\subseteq\B_\eta$.
\end{proof}

\begin{remark}[Source 02's open question, answered]\label{hset:rem:answered-02}
Source 02 proved the forward implication (its necessary conditions) and asked
whether they could be sharpened to a sufficient criterion, noting that any
proof would have to control the reverse interpretation locally. Source 09's
\cref{hset:thm:elementary} is that criterion: the reverse interpretation of
\cref{hset:thm:biinterpretation} and the coherence of
\cref{hset:prop:coherence} give the converse. Between cardinal cutoffs the
equivalence is also proved by source 05 (regular cardinals) and sources 06 and
07 (all uncountable cardinals).
\end{remark}

\begin{corollary}[Two kinds of elementary extension]
The ordered field inclusion $\No_{<\lambda}\subseteq\No_{<\eta}$ is elementary
for all epsilon numbers $\lambda<\eta$. Adding birthday changes its elementary
status according to \cref{hset:thm:elementary}. In particular,
\[
 \No_{<\varepsilon_0}\prec\No_{<\omega_1},
 \qquad
 \B_{\varepsilon_0}\not\prec\B_{\omega_1}.
\]
\end{corollary}
\begin{proof}
Both fields are real closed, so ordered-field quantifier elimination gives
the first assertion. The birthday assertions are instances of the theorem.
\end{proof}

Despite the last failure, both structures in that example are separately
bi-interpretable with $H_{\omega_1}$. There is no contradiction: a
bi-interpretation does not make \emph{every particular inclusion} elementary.
The smaller structure defines an interpreted cutoff set; that formula changes
its meaning under the displayed inclusion.

\begin{example}[Two field-equivalent cutoffs separated by one sentence; sources 05, 07]
\label{hset:ex:omega12}
Let $\chi$ say that every set admits an injection into $\omega$ (defined as the
least infinite ordinal). Then $H_{\omega_1}\models\chi$ and
$H_{\omega_2}\models\neg\chi$, since $\omega_1\in H_{\omega_2}$. Hence the
parameter-free sentence $\chi^\I$ gives $\B_{\omega_1}\not\equiv\B_{\omega_2}$,
although their ordered-field reducts are elementarily equivalent and the
field inclusion is elementary.
\end{example}

\begin{corollary}[Transfer of elementary equivalence; source 06]\label{hset:cor:equiv}
For uncountable cardinals $\kappa,\lambda$,
$\B_\kappa\equiv\B_\lambda\iff H_\kappa\equiv H_\lambda$. The complete theories
$\operatorname{Th}(H_\kappa,\in)$ and $\operatorname{Th}(\B_\kappa)$ are
computably many-one reducible to each other by translations independent of
$\kappa$ (source 05, at regular cardinals).
\end{corollary}
\begin{proof}
Translate sentences through the uniform interpretations in both directions;
the translations are effective syntactic operations.
\end{proof}

\begin{remark}[Bi-interpretation versus elementary equivalence; source 05]
Distinct membership structures can be bi-interpretable without being
isomorphic or elementarily equivalent in their own language; Freire and Hamkins
give relative-consistency instances involving $H_{\omega_1}$ and $H_{\omega_2}$
\cite{FreireHamkins}. \Cref{hset:ex:omega12} does not contradict this: the
theorem here supplies one uniform interpretation between each cutoff and its
own hereditary-size structure, and does not classify all interpretations
between two different hereditary-size structures.
\end{remark}

\begin{corollary}[Comparison with the full class]
\label{hset:cor:fullclass}
For an epsilon number $\lambda$, interpreted formula by formula,
\[
 \B_\lambda\prec\B
 \quad\Longleftrightarrow\quad
 \lambda\text{ is a cardinal and }H_\lambda\prec V.
\]
\end{corollary}
\begin{proof}
The same collector obstruction excludes a noncardinal lower cutoff. For a
cardinal lower cutoff, the uniform interpretation and reverse sign-graph
interpretation extend to definable classes (\cref{hset:thm:chy}). The induced
set map is $H_\lambda\subseteq V$, so the elementary equivalence follows for each
formula. This uses no internal satisfaction class for $V$.
\end{proof}

The condition $H_\lambda\prec H_\eta$ in \cref{hset:thm:elementary} is not
replaced here by a conjectural large-cardinal characterization. It is an exact
reduction of one elementary-inclusion problem to another familiar one,
together with the complete exclusion of every noncardinal initial cutoff. If
both cutoffs are strongly inaccessible, $H_\rho=V_\rho$ and the condition reads
$V_\lambda\prec V_\eta$ (source 05).

\subsection{Elementary embeddings}
\label{hset:sub:embeddings}

Sources 05 (regular cardinals) and 06 (all uncountable cardinals) prove that the
uniform bi-interpretation also transfers elementary embeddings that are not
inclusions. We state the theorem at uncountable cardinals and give source 05's
explicit formula. Write $\name(x)=\langle\bd(x),A_x\rangle$ with
$A_x=\{\alpha<\bd(x):x(\alpha)=+\}$ for the set-theoretic sign name of $x$.

\begin{theorem}[Elementary-embedding correspondence; sources 05, 06]
\label{hset:thm:embeddings}
For uncountable cardinals $\kappa,\lambda$ there is a canonical bijection
\[
 \Emb\bigl((H_\kappa,\in),(H_\lambda,\in)\bigr)
 \longleftrightarrow
 \Emb(\B_\kappa,\B_\lambda),
\]
respecting identities and composition. For $j:H_\kappa\to H_\lambda$ the
corresponding map $e_j$ is given by
\[
 \name(e_j(x))=j(\name(x))=\langle j(\bd(x)),j(A_x)\rangle,
\]
and for $e:\B_\kappa\to\B_\lambda$ the corresponding $j_e$ is
$j_e(x_C)=x_{e(C)}$, with $e$ applied coordinatewise to codes. In particular
$e_j(\alpha)=j(\alpha)$ on ordinals and $e_j(x)(j(\alpha))=x(\alpha)$ for
$\alpha<\bd(x)$.
\end{theorem}
\begin{proof}
Sign names, their order, birthday and arithmetic are uniformly defined in the
membership language (\cref{hset:sec:biinterpretation}), so an elementary $j$
maps sign names to sign names and preserves every translated formula; this
proves elementarity of $e_j$. Since the code domain and $\sim$ are
parameter-free, an elementary $e$ maps codes to codes and equivalent codes to
equivalent codes, so $j_e$ is well defined, and applying $\I$ to each
membership formula proves it elementary. The definable round-trip isomorphisms
commute with elementary maps, because an elementary map preserves their
defining formulas and their unique outputs modulo the equivalence relations;
hence $e_{j_e}=e$ and $j_{e_j}=j$, and composition and identities are
respected. The name of an ordinal $\alpha$ is $\langle\alpha,\alpha\rangle$,
which gives $e_j(\alpha)=j(\alpha)$, and elementarity of $\Bit$ gives the last
identity.
\end{proof}

\begin{remark}[Do not replace $j(A)$ by the pointwise image]
For a nontrivial elementary set embedding, $j(A)$ need not equal
$j``A$. The formula uses the actual image \emph{set} $j(A)$: the sign sequence
$e_j(x)$ has domain $j(\bd(x))$ and may have coordinates outside $j``\bd(x)$,
which a pointwise transport of signs would leave undetermined.
\end{remark}

\begin{corollary}[Critical ordinals and the identity case; source 05]
\label{hset:cor:critical}
Corresponding embeddings move exactly the same ordinals; when a least moved
ordinal exists it is the same for both. An elementary embedding
$\B_\kappa\to\B_\lambda$ fixing every ordinal below $\kappa$ is the canonical
inclusion.
\end{corollary}
\begin{proof}
If $e$ fixes all ordinals, then $\bd(e(x))=e(\bd(x))=\bd(x)$, and elementarity
of the bit formula gives the same signs below that birthday.
\end{proof}

The correspondence is a classification by reduction, not an existence theorem
for such embeddings: it proves no new large-cardinal consistency statement and
does not refute rank-into-rank hypotheses. A proper canonical inclusion moves
no ordinal, so ``proper embedding'' must not be confused with ``embedding
having a critical point''. Every automorphism of $\B_\kappa$ is the identity
(\cref{hset:lem:rigidity}); that statement overlaps the collection's existing
simplicity rigidity and is not counted as new.

\section{Birthday growth and set-existence axioms}
\label{hset:sec:growth}

\subsection{Birthday eternity}

\begin{definition}[The eternity scheme $\Et$; source 09]
For every formula $\varphi(a,x,\vec p)$ in the birthday-field language,
include the assertion that if $\gamma$ is an embedded ordinal and
\[
 \forall a<\gamma\ \exists!x\ \varphi(a,x,\vec p),
\]
then some embedded ordinal $\varrho$ satisfies
\[
 \forall a<\gamma\ \forall x\,
       (\varphi(a,x,\vec p)\longrightarrow\bd(x)<\varrho).
\]
The tuple $\vec p$ consists of arbitrary parameters from the structure.
\end{definition}

This is a precisely specified version of birthday boundedness for definable
ordinal-indexed functions. The analogy with the eternity principle in the
Chen--Hamkins--Yang announcements is intentional; the cutoff classification
below is proved directly, rather than imported from an unexpanded axiom
scheme in the slides.

\begin{theorem}[Eternity forces cardinality of an epsilon cutoff]
\label{hset:thm:eternity-noncardinal}
If $\lambda$ is a noncardinal epsilon number, then
$\B_\lambda\not\models\Et$.
\end{theorem}
\begin{proof}
Put $\mu=\card\lambda$, so $\mu<\lambda$ and
$\kappa(\lambda)=\mu^+$. Choose a bijection $f:\mu\to\lambda$.
Its graph has hereditary size at most $\mu$, so it belongs to
$H_{\mu^+}$ and has a code in $\B_\lambda$. Fix one such code as a
parameter. Using the interpreted graph of $f$ and the ordinal bridge, a
single formula now defines a scalar map on the actual ordinal $\mu$:
at $a<\mu$ its value is the embedded ordinal $f(a)<\lambda$.
Its birthdays exhaust $\lambda$, hence have no bound below $\lambda$.
This violates the corresponding instance of $\Et$.
\end{proof}

At a singular \emph{cardinal} $\kappa$, an external cofinal map into
$\kappa$ generally has hereditary size $\kappa$ and need not belong to
$H_\kappa$. Thus the same argument cannot simply be repeated with
$\mu=\cf(\kappa)$. Whether an internally definable cofinal map exists is the
point of the next theorem.

\subsection{The axioms that all hereditary-size universes already satisfy}

\begin{lemma}[The fixed background axioms]
\label{hset:lem:H-background}
For every uncountable cardinal $\kappa$, $H_\kappa$ satisfies
Extensionality, Empty Set, Pairing, Union, Infinity, Foundation, every
Separation instance, and Choice. Moreover the actual transitive closure of
each of its elements belongs to $H_\kappa$, and its cardinality is correctly
computed there.
\end{lemma}
\begin{proof}
Transitivity gives Extensionality and Foundation. Empty Set, Pairing, Union,
and Infinity preserve hereditary size below $\kappa$. For any fixed formula
and parameters, the externally defined subset of a set $a\in H_\kappa$
satisfying that formula in the set structure $H_\kappa$ has hereditary size
bounded by that of $a$ (up to finitely many additional objects). It therefore
belongs to $H_\kappa$, proving Separation, including formulas with unbounded
quantifiers in that structure.

An ambient choice function for a family $a\in H_\kappa$ has its domain and
chosen values inside the transitive closure of $a$, and its graph has the
same infinite hereditary-cardinality bound. It belongs to $H_\kappa$.
This proves Choice for set families, not a class-choice principle.

Writing $T=\TC(\{x\})$, the transitive closure of $T$ adds at most the set
$T$ to its already transitive data, so $T\in H_\kappa$. Its description as
the least transitive set containing $x$ as an element is correct internally:
the actual least such set is available there. A bijection between $T$ and
its least equipotent ordinal has hereditary size below $\kappa$, and any
bijection asserted internally is externally a bijection. Minimality among
ordinals is therefore also computed correctly.
\end{proof}

\subsection{Replacement is exactly eternity at cardinal cutoffs}

\begin{theorem}[Eternity--Replacement equivalence]
\label{hset:thm:eternity-replacement}
For every uncountable cardinal $\kappa$,
\[
 \B_\kappa\models\Et
 \quad\Longleftrightarrow\quad
 H_\kappa\models\text{every instance of Replacement}.
\]
\end{theorem}
\begin{proof}
Assume first Replacement in $H_\kappa$. A definable scalar function on an
ordinal $\gamma<\kappa$ corresponds under the bi-interpretation to a
set-theoretically definable function assigning to $a<\gamma$ its unique
sign graph. Replacement collects those sign graphs in one set of
$H_\kappa$. The union of their domains, with one added successor for each
length, is bounded below $\kappa$: otherwise the transitive closure of the
collected set would contain at least $\kappa$ ordinals. More explicitly,
\[
 \varrho=\sup\{\dom(s)+1:s\text{ is an image sign graph}\}<\kappa
\]
bounds their birthdays strictly. This proves the corresponding eternity
instance; the empty-domain case is vacuous.

Conversely, assume eternity and consider a formula defining, in
$H_\kappa$, a function $F:a\to H_\kappa$, where $a\in H_\kappa$ and the
formula has fixed parameters. Choose a bijection $e:\delta\to a$ with
$\delta<\kappa$. Its graph belongs to $H_\kappa$. For each $i<\delta$, put
externally
\[
 \nu_i=\card{\TC(\{F(e(i))\})},
\]
regarded as an initial ordinal. By \cref{hset:lem:H-background}, this is a
uniquely defined ordinal of $H_\kappa$, with the correct external value.
The definition of $F$, the parameter $e$, and the hereditary-cardinality
formula translate into a definable scalar map $i\mapsto\nu_i$ in
$\B_\kappa$. Here the ordinal bridge is available for every ordinal below
$\kappa$. No definable choice of a code for $F(e(i))$ is required: its
\emph{hereditary cardinal} is unique.

Eternity gives a $\varrho<\kappa$ bounding all these ordinals. Ambient
Replacement forms the actual image set $A=F[a]$, and its transitive closure
has cardinality at most
\[
 \max(\aleph_0,\card\delta,\card\varrho)<\kappa.
\]
Indeed it is contained, up to the image set itself, in the union of
$\card\delta$ transitive closures, each of cardinality at most
$\max(\aleph_0,\card\varrho)$. Thus $A\in H_\kappa$, proving the desired
Replacement instance in that structure. This proof uses a fixed cardinal
bound below $\kappa$; it never assumes that an arbitrary union of fewer than
$\kappa$ small sets is small.
\end{proof}

\begin{corollary}\label{hset:cor:regular-eternity}
If $\kappa$ is uncountable and regular, then $\B_\kappa\models\Et$.
\end{corollary}
\begin{proof}
Any externally given function on an ordinal below $\kappa$ has fewer than
$\kappa$ birthdays in its image. Regularity bounds them below $\kappa$.
This argument applies in particular to definable functions.
\end{proof}

This corollary genuinely uses regularity; the equivalence of the theorem does
not.

\subsection{Power Set and packing a birthday fiber}

\begin{lemma}[Power Set in a hereditary-size universe; sources 02, 06, 09]
\label{hset:lem:Hpower}\label{hset:prop:powerset}
For every infinite cardinal $\kappa$,
\[
 H_\kappa\models\text{Power Set}
 \quad\Longleftrightarrow\quad
 \kappa\text{ is strong limit}\quad(2^\nu<\kappa\text{ for all cardinals }\nu<\kappa).
\]
For an individual $\alpha<\kappa$, the actual power set $\Pow(\alpha)$ is
an element of $H_\kappa$ exactly when $2^{\card\alpha}<\kappa$. Regularity is
not required.
\end{lemma}
\begin{proof}
Every subset of a set $a\in H_\kappa$ is already in $H_\kappa$, and the
internal subset relation is the actual one by transitivity. So an internal power
set of $a$, if it exists, is the full actual power set $\Pow(a)$. Its
hereditary size is $2^{\card a}$ up to harmless finite or countable bounds.
This proves necessity (apply it to the cardinals $\nu<\kappa$) and the
individual assertion. Conversely, for $x\in H_\kappa$ with $\kappa$ uncountable,
the transitive closure of $\Pow(x)$ has size at most
$\max(\aleph_0,\card{\TC(\{x\})},2^{\card x})$, which is below a strong-limit
$\kappa$. If $\kappa=\omega$, power sets of hereditarily finite sets are
hereditarily finite.
\end{proof}

\begin{theorem}[The translated Power Set sentence; sources 02, 06]\label{hset:thm:power}
Let $\PowI$ be the $\I$-translation of the Power Set axiom
$\forall a\,\exists p\,\forall x\,(x\in p\leftrightarrow\forall y(y\in x\to y\in a))$;
it is one fixed parameter-free sentence of the birthday-field language, not a
scheme. For every epsilon number $\lambda$,
\[
 \B_\lambda\models\PowI
 \quad\Longleftrightarrow\quad
 \lambda\text{ is a strong-limit cardinal}.
\]
Among uncountable regular cutoffs it holds exactly at the strongly inaccessible
ones.
\end{theorem}
\begin{proof}
By \cref{hset:thm:collapse-range} and \cref{hset:lem:Hpower}, $\PowI$ holds
exactly when $\kappa(\lambda)$ is strong limit. If $\lambda$ is not a cardinal,
write $\nu=\card\lambda$; then $\kappa(\lambda)=\nu^+$ and Cantor's theorem gives
$2^\nu\ge\nu^+$, so the condition fails. If $\lambda$ is a cardinal, its ceiling
is itself. Adding regularity yields strong inaccessibility.
\end{proof}

Source 02 proves this at every epsilon number, source 06 at uncountable
cardinals (its $\mathsf{Pow}^\sharp$). It would be incorrect to remove the words
``among uncountable regular cutoffs'': a singular strong-limit cardinal such as
$\beth_\omega$ also satisfies $\PowI$, and \cref{hset:sec:prikry} shows that the
same enriched field can have a regular cutoff in one universe and a singular
one in another. The sentence detects the strong-limit property, not
regularity, Replacement, or full ZFC.

\begin{definition}[The fiber-packing principle $\Pack$; source 09]
For every embedded ordinal $\alpha$, there exist a positive embedded ordinal
$\theta$ and a scalar $r$, with
\[
 \chi=\theta+\alpha+1,
 \qquad\bd(r)=\Lambda(\chi),
\]
such that
\begin{equation}
 \forall u\,[\bd(u)=\alpha\longrightarrow
   \exists\xi<\theta\ \forall\zeta<\alpha\,
    (\Bit(u,\zeta)\leftrightarrow\Bit(r,p_\chi(\xi,\zeta)))].
 \label{hset:eq:packing}
\end{equation}
\end{definition}

This is one first-order sentence, using the previously defined bit and
pairing formulas; it is \emph{not} the Power Set axiom. Its rows list every sign
string of birthday $\alpha$; repetitions are allowed. Unused coordinates can
be ignored. It does not use ordinary ordinal concatenation or a primitive
power-set operation.

\begin{theorem}[Exact packing threshold]
\label{hset:thm:packing}
In $\B_\lambda$, the instance of \eqref{hset:eq:packing} for a particular
$\alpha<\lambda$ holds exactly when
\[
 2^{\card\alpha}<\kappa(\lambda),
 \quad\text{equivalently, }2^{\card\alpha}<\lambda,
\]
where the cardinal on the left is compared as an ordinal in the second
inequality. Consequently $\Pack$ holds exactly when $\lambda$ is an
uncountable strong-limit cardinal.
\end{theorem}
\begin{proof}
A packing table with $\theta<\lambda$ rows gives a surjection from $\theta$
onto $\Pow(\alpha)$, so $2^{\card\alpha}\le\card\theta<\kappa(\lambda)$.
Conversely, if that power-set cardinal is $\nu<\kappa(\lambda)$, then
$\nu<\lambda$ by cardinal rounding. Choose an enumeration of the fiber by
$\theta=\nu$ and encode its table in the square of width
$\chi=\theta+\alpha+1$. The square code stays below $\lambda$ by closure,
so the required $r$ exists. The equivalence of the two inequalities follows
because no cardinal lies in $[\lambda,\kappa(\lambda))$.

For a cardinal $\lambda$, validity for all $\alpha<\lambda$ is exactly the
strong-limit property. If $\lambda$ is not a cardinal, let
$\mu=\card\lambda<\lambda$ and use $\alpha=\mu$. Cantor's theorem gives
$2^\mu\ge\mu^+=\kappa(\lambda)$, so this instance fails.
\end{proof}

Notice that for a noncardinal cutoff $\lambda$, the recovered
$H_{\kappa(\lambda)}$ and the scalar domain do not have the same ordinal
height. Assertions about \emph{all} interpreted ordinals and assertions
about only scalar ordinals must be kept separate. The failure at
$\alpha=\card\lambda$ is enough to settle the full packing principle.

\begin{center}\small
\begin{tabular}{@{}P{0.3\textwidth}P{0.2\textwidth}P{0.14\textwidth}P{0.14\textwidth}@{}}
\toprule
Cutoff $\lambda$ & Recovered universe & $\exists C\,\Coll(C)$ & $\PowI$, $\Pack$\\
\midrule
Countable epsilon number & $H_{\omega_1}$ & true & false\\
Noncardinal, $\card\lambda=\nu$ & $H_{\nu^+}$ & true & false\\
Uncountable successor cardinal $\kappa$ & $H_\kappa$ & false & false\\
Singular strong-limit cardinal $\kappa$ & $H_\kappa$ & false & true\\
Strongly inaccessible $\kappa$ & $H_\kappa=V_\kappa$ & false & true\\
\bottomrule
\end{tabular}
\end{center}
The table is source 02's, with the packing column added from source 09.

\section{Worldly cutoffs and separation examples}
\label{hset:sec:worldly}

\subsection{The exact conjunction}

Call an uncountable cardinal $\kappa$ \emph{worldly} here when
$V_\kappa\models\mathrm{ZFC}$. This definition does not impose regularity.
We need the following familiar set-theoretic fact, for which a proof is
included to fix the precise meaning.

\begin{lemma}[When hereditary size becomes cumulative rank]
\label{hset:lem:worldly-HV}
For an uncountable cardinal $\kappa$,
\[
 H_\kappa\models\mathrm{ZFC}
 \quad\Longleftrightarrow\quad
 H_\kappa=V_\kappa\text{ and }V_\kappa\models\mathrm{ZFC}.
\]
Moreover, $V_\kappa\models\mathrm{ZFC}$ by itself implies
$H_\kappa=V_\kappa$.
\end{lemma}
\begin{proof}
Suppose $H_\kappa\models\mathrm{ZFC}$. Its internal cumulative hierarchy
exists to every ordinal $\alpha<\kappa$. Every subset of each of its sets
belongs to $H_\kappa$, so its internal power sets are actual power sets.
Induction on $\alpha$ shows that its internal $V_\alpha$ is the actual
$V_\alpha$; the limit step is ordinary union over the actual smaller
ordinals. Thus every $V_\alpha$ for $\alpha<\kappa$ belongs to $H_\kappa$,
and transitivity gives $V_\kappa\subseteq H_\kappa$.
The reverse inclusion is \cref{hset:lem:hereditary}.

Conversely, assume $V_\kappa\models\mathrm{ZFC}$. For every
$\alpha<\kappa$, that model can well-order $V_\alpha$ and supply a bijection
with some ordinal $\delta<\kappa$. These are actual sets and actual
bijections, because the model is transitive. Hence
$\card{V_\alpha}<\kappa$. The transitive closure of any
$x\in V_\kappa$ is contained in some such $V_\alpha$, and therefore has
size below $\kappa$. Again $H_\kappa=V_\kappa$ follows. The remaining
implication is immediate from equality.
\end{proof}

\begin{theorem}[Worldly-cutoff characterization]
\label{hset:thm:worldly}
For an epsilon number $\lambda$,
\begin{align*}
 \B_\lambda\models\Et+\Pack
 &\iff \lambda\text{ is a cardinal and }H_\lambda\models\mathrm{ZFC}\\
 &\iff \lambda\text{ is a worldly cardinal}.
\end{align*}
For regular uncountable $\lambda$, this holds exactly when $\lambda$ is
strongly inaccessible.
\end{theorem}
\begin{proof}
Either growth principle excludes noncardinal epsilon cutoffs. At cardinal
cutoffs,
\cref{hset:thm:eternity-replacement,hset:lem:Hpower,hset:lem:H-background}
identify the two missing ZFC requirements as Replacement and Power Set.
The equivalence with the rank formulation is \cref{hset:lem:worldly-HV}.
For regular uncountable cardinals, eternity is automatic and packing is
exactly strong limit, giving the usual inaccessible condition.
\end{proof}

This proves \cref{hset:thm:main-growth}. It does not identify the two formulas
with every formulation of a theory called SA. It is an exact theorem about the
explicit scheme and sentence defined here.

\subsection{The two growth principles do different work}

\begin{example}[Eternity without fiber packing]
At $\kappa=\omega_1$, regularity gives eternity. The birthday-$\omega$ fiber
has $2^{\aleph_0}\ge\aleph_1$ members, so no ordinal below $\omega_1$ can
index a packing of it. Hence
\[
 \B_{\omega_1}\models\Et,
 \qquad\B_{\omega_1}\not\models\Pack.
\]
In the decoded structure, every individual real is present, but the set of
all reals is not an element of $H_{\omega_1}$.
\end{example}

\begin{example}[Fiber packing without eternity]
\label{hset:ex:bethomega}
Let $\kappa=\beth_\omega=\sup_{n<\omega}\beth_n$, with
$\beth_0=\aleph_0$ and $\beth_{n+1}=2^{\beth_n}$. This is a strong-limit
cardinal: for every $\mu<\kappa$ some $\beth_n$ is at least $\mu$, and then
$2^\mu\le\beth_{n+1}<\kappa$. Thus packing holds.

In $H_\kappa$, each finite initial segment of the beth recursion is present
and correctly computed. The necessary finite sequence of power sets,
cardinals, and bijections has hereditary size below $\kappa$, and Power Set
is correct there by \cref{hset:lem:Hpower}. A single formula therefore defines
$n\mapsto\beth_n$ on the actual $\omega$. Transporting this formula through
the ordinal bridge gives a definable scalar function whose birthdays are
cofinal in $\kappa$. Consequently
\[
 \B_{\beth_\omega}\models\Pack,
 \qquad\B_{\beth_\omega}\not\models\Et.
\]
The infinite sequence of all these ordinals is precisely what cannot be
collected as an element of $H_\kappa$.
\end{example}

Source 06 isolates the same phenomenon as one explicit sentence, and uses it to
separate $\beth_\omega$ from its successor. Let $\vartheta(n,\gamma)$ be the
membership formula saying that a finite sequence of length $n+1$ begins with
$\omega$, ends with $\gamma$, and at each step takes the least ordinal
equinumerous with the power set of the previous entry; and let
\[
 \RepB:\quad
 [\forall n\in\omega\ \exists!\gamma\ \vartheta(n,\gamma)]
 \longrightarrow
 \exists y\ \forall\gamma\,[\gamma\in y\leftrightarrow
 \exists n\in\omega\ \vartheta(n,\gamma)],
\]
a single Replacement instance.

\begin{theorem}[A singular Replacement boundary; source 06]\label{hset:thm:beth}
With $\kappa=\beth_\omega$, $H_\kappa\not\models\RepB$ and
$H_{\kappa^+}\models\RepB$, the antecedent being true in both. Consequently
\[
 \B_{\beth_\omega}\not\prec\B_{\beth_\omega^+},
\]
although $\B_{\beth_\omega}\models\PowI$ and the ordered-field inclusion is
elementary.
\end{theorem}
\begin{proof}
Every finite beth computation, with the power sets and bijections at its
finitely many stages, has hereditary size below $\kappa$ and belongs to both
structures, where it is correctly computed; so $\vartheta(n,\gamma)$ picks out
the actual $\beth_n$ in both. A set collecting all $\beth_n$ has every ordinal
below $\kappa$ in its transitive closure, so it is not in $H_\kappa$; the actual
set $\{\beth_n:n<\omega\}$ has hereditary size $\kappa$ and lies in
$H_{\kappa^+}$. Transfer by \cref{hset:thm:elementary} (or the translation
$\RepB^\I$) gives the failure of elementarity, and \cref{hset:thm:power} the
Power Set assertion.
\end{proof}

The theorem gives one singular example with failed Replacement. It does
\emph{not} show that every singular $\kappa$ fails Replacement in $H_\kappa$: an
externally cofinal sequence can fail to be definable over $H_\kappa$ or fail to
be an element of it. Nor does internal ZFC imply external inaccessibility;
\cref{hset:cor:prikry} gives a sharp example.

\begin{example}[A noncardinal field cutoff]
Both principles fail at $\varepsilon_0$. Packing fails at the ordinal
$\omega$. Eternity fails by \cref{hset:thm:eternity-noncardinal}; one may also
use the definable finite power-tower sequence cofinal in $\varepsilon_0$,
with ordinary ordinal exponentiation computed in the decoded
$H_{\omega_1}$. Nevertheless the field is real closed and its birthday
expansion is bi-interpretable with hereditary countability.
\end{example}

These are actual structures exhibiting independence of the two displayed
growth requirements within this family. They are not an assertion of
independence from an unspecified list of all the other axioms of SA.

\subsection{Singular cutoffs satisfying both principles}

\begin{proposition}[A conditional source of singular worldly cutoffs]
\label{hset:prop:singular-worldly}
If $\Omega$ is strongly inaccessible, there is a cardinal $\kappa<\Omega$
of cofinality $\omega$ such that
\[
 V_\kappa\prec V_\Omega.
\]
In particular, $\kappa$ is worldly and $\B_\kappa$ satisfies both growth
principles.
\end{proposition}
\begin{proof}
Enumerate all existential formulas in the language of set theory. Construct
strictly increasing infinite cardinals $\alpha_n<\Omega$. Given
$\alpha_n$, consider the first $n$ formulas and all their finite parameter
tuples from $V_{\alpha_n}$. Whenever $V_\Omega$ has a witness, choose one.
There are fewer than $\Omega$ such requirements, since
$\card{V_{\alpha_n}}<\Omega$. Regularity bounds their witness ranks below
$\Omega$. Choose a cardinal $\alpha_{n+1}<\Omega$ above that bound and above
$\alpha_n$.

Set $\kappa=\sup_n\alpha_n$. Regularity of $\Omega$ gives $\kappa<\Omega$;
it is a cardinal of cofinality $\omega$. Every finite tuple in $V_\kappa$
appears in one stage, and every existential formula is treated at all
sufficiently late stages. Thus every existential statement true in
$V_\Omega$ with parameters in $V_\kappa$ has a witness in $V_\kappa$.
The Tarski--Vaught criterion gives elementarity. Since $V_\Omega$ satisfies
ZFC, so does $V_\kappa$. Apply \cref{hset:thm:worldly}.
\end{proof}

The reflection construction is standard set-theoretic reasoning, included
as an illustration rather than as a novelty claim, and it needs the stated
inaccessible. It demonstrates why replacing ``worldly'' by ``inaccessible'' in
\cref{hset:thm:worldly} would be wrong. No inaccessible cardinal is assumed in
the unconditional cutoff theorems.

\section{Model-theoretic consequences}
\label{hset:sec:modeltheory}

\subsection{Full second-order arithmetic and undecidability}

\begin{proposition}[The standard natural numbers and all their subsets; sources 06, 07, 09]
\label{hset:prop:arithmetic}
In $\B$ and in every $\B_\lambda$ ($\lambda$ an epsilon number) there is a
parameter-free interpretation of the two-sorted structure
$(\mathbb N,\Pow(\mathbb N);0,1,+,\cdot,<,\in)$ with its \emph{full} actual
subset sort. The number sort consists of the ordinals below the least nonzero
limit ordinal $w$ (externally $w=\omega$), the subset sort of the $s$ with
$\bd(s)=w$, and membership is $\Bit(s,n)$.
\end{proposition}
\begin{proof}
``$q$ is a nonzero limit ordinal'' is
$\Ord(q)\land q>0\land\forall a<q\,\exists d\,(a<d<q)$ with ordinal guards, and
$w$ is its least solution. The ordinals below $w$ are the standard finite
ordinals, and natural arithmetic on them is ordinary arithmetic. The fiber
$\bd(s)=\omega$ is in bijection with all subsets of $\omega$ by
\eqref{hset:eq:fiber-subsets}; no restriction to definable subsets is imposed.
\end{proof}

When working inside a model $M$, the subset sort is $\Pow(\omega)^M$, as it
must be. Sources 06 and 07 state the proposition, and the next two results, at
uncountable cardinals and source 09 at $\varepsilon_0$; the proofs use only
$\omega<\lambda$, which holds at every epsilon number. This is an ordinary
first-order interpretation of a two-sorted
structure, not a replacement of first-order by second-order syntax. It does not
make set theory definable in the pure ordered-field reduct, where the birthday
structure is absent.

\begin{corollary}[Undecidability; sources 05, 06, 07, 09]\label{hset:cor:undecidable}
True first-order arithmetic reduces effectively to the complete theory of
every $\B_\lambda$, and formula by formula to $\B$. Consequently these complete
theories are neither decidable nor computably enumerable.
\end{corollary}
\begin{proof}
Relativize arithmetic quantifiers to the number sort. If the target theory were
computably enumerable, the inverse image would enumerate true arithmetic; a
complete computably enumerable theory is decidable, and true arithmetic is not.
\end{proof}

These are logical-complexity consequences, not new algorithms. A code can
contain an arbitrary noncomputable subset of an ordinal, and first-order
definability in the birthday expansion must not be confused with decidability
from a finite symbolic representation.

\subsection{Independence and an omitted type}

\begin{proposition}[Independence property and a countable omitted type; source 07]
\label{hset:prop:ip}
The formula $\Bit(s,n)$ has the independence property in every $\B_\lambda$.
The parameter-free partial type
$\{x<w\land\Ord(x)\}\cup\{x>n:n\in\mathbb N\}$, with each standard $n$ named by
its definition, is finitely satisfiable but omitted.
\end{proposition}
\begin{proof}
For distinct finite ordinals $n_1,\dots,n_k$ and any
$S\subseteq\{1,\dots,k\}$, a length-$\omega$ sign code positive exactly at the
$n_j$ with $j\in S$ realizes the pattern. Finite parts of the type are realized
by large finite ordinals; a realization of the whole type would be a standard
natural number above all standard natural numbers.
\end{proof}

\begin{remark}[Saturation of the pure field, and the sibling report]
\label{hset:rem:saturation}
For regular uncountable $\kappa$, the pure real closed field
$\No_{<\kappa}$ realizes every one-variable type over fewer than $\kappa$
parameters: by quantifier elimination a nonalgebraic type over the real closed
subfield generated by the parameters is a cut, regularity bounds the birthdays
of its fewer than $\kappa$ endpoints, and \eqref{hset:eq:cut-bound} supplies a
separator. Source 07 records this as the usual saturation property of the pure
field, not a new result. Naming birthday nevertheless produces the omitted type
of \cref{hset:prop:ip}. External saturation of the pure field $\No^M$ inside an
outer universe is the subject of the sibling report
\repo{foundations-and-computation/surreal-fields-across-universes}; this report
does not study it.
\end{remark}

\subsection{Birthday is not visible in the field reduct}

\begin{proposition}[Nondefinability of birthday; sources 05, 07]\label{hset:prop:nondef}
On $\No_{<\lambda}$, for $\lambda$ an epsilon number, neither birthday nor prefix
simplicity is definable, even with parameters, in the ordered-field reduct, nor
in any o-minimal expansion of it. The same holds formula by formula on $\No$.
\end{proposition}
\begin{proof}
If birthday were definable, its fixed-point set, the embedded ordinals, would
be definable. It is infinite and contains no interval: between ordinals
$\alpha<\gamma$ lies $\alpha+\tfrac12$, which is not an ordinal. An o-minimal
structure (in particular a real closed field, by quantifier elimination) has
only finite unions of points and intervals as definable subsets of the line. If
prefix simplicity were definable, the ordinals would be defined by
``every proper prefix of $x$ lies below $x$'', with the same contradiction.
\end{proof}

Source 05 states the proposition at regular cardinals and source 07 at
uncountable cardinals; the proof uses only real closedness (o-minimality). The
collection's foundations report notes that the field language has no
birthday symbol; this proposition is the corresponding standard
nondefinability proof, new to the collection but not a novelty claim.

\subsection{How little arithmetic is needed initially}

Let $G_\varpi(\alpha,\gamma,\rho)$ be the ternary relation on $\No_{<\kappa}$ saying
that $\alpha,\gamma,\rho$ are ordinals and $\rho=\varpi(\alpha,\gamma)$, with
$\varpi$ the polynomial of \cref{hset:lem:pairing06}.

\begin{corollary}[Reconstruction from ordinal pairing; source 06]\label{hset:cor:reduct}
For every uncountable cardinal $\kappa$, all operations and constants of
$\B_\kappa$ are uniformly definable in $(\No_{<\kappa};<,\bd,G_\varpi)$. Hence the two
structures are definitionally equivalent.
\end{corollary}
\begin{proof}
Run the coding of \cref{hset:sec:codes} with relation words of length
$\varpi(\delta,\delta)+1$ whose plus positions are the $\varpi(u,v)$, as source
06 does; ordinal successor is definable from the ordinal order. The code
formulas then use field arithmetic only to evaluate $\varpi$ on ordinals, which
$G_\varpi$ supplies, so this reduct still interprets $H_\kappa$. Translate the
arithmetic definitions of \cref{hset:sec:certificates} through that
interpretation and use the surreal-side round trip to put their graphs on the
original carrier. The converse is the polynomial formula for $G_\varpi$.
\end{proof}

No minimal-language claim is made: the corollary does not remove birthday data
or show that arithmetic is definable from order and prefix simplicity alone.
It is partial progress on the minimal-signature questions of sources 02, 05 and
07, which stay open (\cref{hset:sec:questions}).

\section{The exact Choice boundary of ordinal graph coding}
\label{hset:sec:choice}

This section, source 06's, works in \emph{ZF}. It considers the full class of
sign sequences and the graph interpretation of \cref{hset:sec:codes} without a
cutoff. All its defining formulas remain first-order formulas about individual
sets or surreals, and class assertions are read formula by formula.

\begin{definition}
Let $\HWO=\{x:\TC(\{x\})\text{ admits a well-ordering}\}$. One well-ordering of
the whole transitive closure is required, with no definability condition such as
the one in $\mathrm{HOD}$.
\end{definition}

\begin{theorem}[Choice-sensitive range; source 06]\label{hset:thm:choice}
Over ZF, the pointed ordinal-graph interpretation has quotient structure
isomorphic to $(\HWO,\in)$. Its external collapse map is surjective onto $V$ if
and only if the Axiom of Choice holds.
\end{theorem}
\begin{proof}
The Mostowski collapse of a valid code $(\delta,r,t)$ is injective with domain
the ordinal $\delta$, and its range is $\TC(\{x_C\})$ by root generation;
transporting the ordinal order well-orders that closure. Conversely, a
well-order of $\TC(\{x\})$ gives a bijection from an ordinal, along which
membership pulls back to a valid code for $x$; no simultaneous choice for a
family is used. Equality and membership remain correct in ZF: all subsets of a
fixed ordinal exist and have characteristic words, so the well-foundedness test
still quantifies over every subset, and Mostowski collapse is a theorem of ZF.
If Choice holds every transitive closure is well-orderable. Conversely,
surjectivity well-orders every $\TC(\{x\})$ and hence every $x\subseteq
\TC(\{x\})$, and the well-ordering principle implies Choice.
\end{proof}

\begin{proposition}[Internal Choice is a different assertion; source 06]
\label{hset:prop:internal-choice}
In ZF, $\HWO$ is transitive, contains every ordinal and every subset of an
ordinal, and internally satisfies the statement that every set can be
well-ordered. If $\Pow(\omega)$ is not well-orderable in the ambient universe,
$(\HWO,\in)$ fails the Power Set axiom at $\omega$.
\end{proposition}
\begin{proof}
If $y\in x\in\HWO$ then $\TC(\{y\})\subseteq\TC(\{x\})$ is well-orderable. For
$a\subseteq\alpha$, $\TC(\{a\})\subseteq\{a\}\cup\alpha$. For $x\in\HWO$ the
restriction to $x$ of a well-order of $\TC(\{x\})$, and every subset of $x$, lie
in $\HWO$, since finite products and finite subsets of a well-orderable set are
well-orderable in ZF; so the internal well-order has a witness. Every subset of
$\omega$ is in $\HWO$, so an internal power set of $\omega$ would be the actual
$\Pow(\omega)$, which is then well-orderable, contrary to the hypothesis.
\end{proof}

The surjectivity statement of \cref{hset:thm:choice} is an external comparison
of ranges, not the value of the internal Choice sentence. The theorem is about
the specified ordinal-graph construction; no claim is made that every possible
interpretation of sets in surreals has this boundary. Global Choice is
unnecessary here and for the local ZFC bi-interpretations: existence of a code
is proved object by object and the quotient identifies all codes.

\begin{remark}[Source 09's fourth target, answered]\label{hset:rem:answered-choice}
Source 09 listed ``the choice-free analogue'' among its targets and observed
that ordinal-indexed relations naturally recover the sets with well-orderable
transitive closures. \Cref{hset:thm:choice}, source 06's, is that analogue for
the full class; source 07 records the same range in a remark without a separate
proof. Whether a richer coding signature can escape the boundary, and a
choice-free version of the \emph{cutoff} theorems, remain open.
\end{remark}

\part{Outer models}\label{hset:part:outer}

Throughout this part, $M\subseteq N$ are transitive models of ZFC with the same
ordinals. They may be set models viewed from an ambient metatheory, or class
models with formulas interpreted one at a time. Superscripts indicate the model
in which a construction is performed, and all references to $M$-cardinality are
internal to $M$. The intended examples are set-forcing extensions, but no
result below assumes forcing unless it says so.

\section{Old arithmetic and cross-model coherence}
\label{hset:sec:absolute}\label{hset:sec:extensions}

\begin{lemma}[An old simplest separator stays simplest; sources 02, 05, 07]
\label{hset:lem:oldcuts}
If $L,R\in M$ are separated sets of old sign sequences and $w\in M$ is their
simplest separator computed in $M$, then $w$ is their simplest separator in
$N$.
\end{lemma}
\begin{proof}
Comparisons of old sign sequences are absolute, so $w$ still separates. Suppose
$z\in N$ separates and is shorter. If $z$ is a prefix of $w$ it is a restriction
of $w$ to an ordinal, hence old, contradicting minimality in $M$. Otherwise $z$
and $w$ first disagree before $z$ ends; their common prefix $p$ is a
restriction of $w$, hence old, lies strictly between $z$ and $w$, and is a
shorter separator by convexity of the set of separators, with all its
comparisons with $L,R$ holding in $M$. This again contradicts minimality.
\end{proof}

Source 02 gives a constructive variant: build the separator sign by sign,
appending $+$ if some $\ell\in L$ is at least the current word and $-$
otherwise, and taking unions at limits; every test involves only the given sets
and the previous prefix, so the stopping place and every sign agree in $M$ and
$N$.

\begin{proposition}[Absolute old field structure; sources 02, 05, 07]
\label{hset:prop:absolute}
The inclusion of old sign sequences $\No^M\subseteq\No^N$ preserves and
reflects order and preserves birthday, constants, negation, addition,
multiplication and inverses. Its ordered-field reduct is an elementary
inclusion, and so is the pure-field inclusion $\No^M[i]\subseteq\No^N[i]$. The
same holds for $\No_{<\lambda}^M\subseteq\No_{<\lambda}^N$ at every common epsilon
number $\lambda$.
\end{proposition}
\begin{proof}
Order, length, constants and sign reversal are absolute. For addition induct on
pairs of old prefixes: all prefixes of an old number are old, the options are
identified by induction, and \cref{hset:lem:oldcuts} identifies their simplest
separator. Multiplication follows by the same induction using addition, and
inverses by uniqueness. Both real fields are real closed and both complex
fields algebraically closed of characteristic zero, so quantifier elimination
makes the inclusions elementary \cite{Marker,Deloro}. For proper classes the
equivalences are applied one formula at a time.
\end{proof}

This is the usual absoluteness of the canonical construction, not an assertion
about arbitrary chosen presentations of isomorphic fields. It does not say that
the outer universe cannot distinguish the inner one once extra structure is
named: the rest of this part shows that birthday does.

\begin{proposition}[Cross-model coherence; source 09]
\label{hset:prop:cross-model}
Let $\kappa$ be an uncountable cardinal of both models. There are actual
inclusions
\[
 \B_\kappa^M\subseteq\B_\kappa^N,
 \qquad H_\kappa^M\subseteq H_\kappa^N,
\]
compatible with the interpretations and their definable isomorphisms.
Consequently
\[
 \B_\kappa^M\prec\B_\kappa^N
 \quad\Longleftrightarrow\quad H_\kappa^M\prec H_\kappa^N.
\]
\end{proposition}
\begin{proof}
A sign sequence of $M$ remains the same sign sequence in $N$; order and
birthday are absolute. A small arithmetic certificate existing in $M$
remains valid in $N$, so the field operations agree. An old set of hereditary
size below $\kappa$ comes with an old bijection witnessing that bound; since
$\kappa$ is still a cardinal, it belongs to $H_\kappa^N$.

For code validity, one must not claim that the two models have the same
subsets of the coding ordinal. Instead, a relation well founded in $M$ has
an ordinal rank function in $M$, which remains decreasing in $N$. Its
well-foundedness therefore persists despite the extra subsets in $N$.
Root generation says that every node is connected to the root by a finite
downward path, as established in \cref{hset:lem:collapse}; these paths are
absolute. Extensionality is bounded and absolute. Failure of any of the
code conditions in $M$ also persists, using its old witness. Thus validity
on old codes is unchanged.

The old Mostowski collapse still satisfies its bounded defining equations
in $N$ and hence is the same collapse. Equality and membership of decoded
old sets are unchanged. Witnesses to their equality and membership can be
constructed in $M$ from those collapses, so their required graph codes
already exist there. The reverse sign-graph interpretation likewise sends
old objects to their actual images. Applying each interpretation to an
elementary inclusion now proves the equivalence.
\end{proof}

Unlike two cutoffs in a single universe, $\B^M_\kappa$ and $\B^N_\kappa$ may
have different birthday fibers even at the same small ordinal. Their ordered
field reducts nevertheless always include elementarily
(\cref{hset:prop:absolute}).

\section{The preservation criterion and the first new birthday}
\label{hset:sec:preservation}

\begin{theorem}[Exact preservation criterion; sources 02, 06, 07, 09]
\label{hset:thm:preservation}
For every common ordinal $\gamma$,
\begin{equation}
 (\No_{<\gamma})^M=(\No_{<\gamma})^N
 \quad\Longleftrightarrow\quad
 \forall\alpha<\gamma\ \ \Pow(\alpha)^M=\Pow(\alpha)^N.
 \label{hset:eq:stage-agreement}
\end{equation}
If $\gamma=\lambda$ is an epsilon number, these conditions are equivalent to the
literal equality $\B_\lambda^M=\B_\lambda^N$ of the canonically presented
structures, including all operations and the birthday function. If
$\gamma=\kappa$ is an uncountable cardinal in both models, they are also
equivalent to $H_\kappa^M=H_\kappa^N$.
\end{theorem}
\begin{proof}
Taking plus positions is an absolute bijection between sign sequences of
length $\alpha$ and subsets of $\alpha$; this gives \eqref{hset:eq:stage-agreement}.
Once the carriers agree, \cref{hset:prop:absolute} gives equality of all
operations; conversely, equal structures have equal carriers.

For $H_\kappa$: a subset of $\alpha<\kappa$ has hereditary size below $\kappa$ in
either model, so $H^M_\kappa=H^N_\kappa$ implies that no bounded subset is new.
Conversely assume no new bounded subsets. Old sets of hereditary size below
$\kappa$ stay so in $N$ (\cref{hset:prop:cross-model}). For $x\in H_\kappa^N$ take in
$N$ a valid code $(\delta,r,t)$ with $\delta<\kappa$. Its relation word has
length $\Lambda(\delta)<\kappa$, so it is old, as are $\delta$ and $t$. The coded
relation is well founded in $N$, hence in $M$ (an old nonempty subset without a
minimal element would also be one in $N$); extensionality is absolute; the
collapse computed in $M$ satisfies its recursion in $N$, so it is the collapse
there. Hence $x\in M$, with hereditary size below $\kappa$ witnessed in $M$.
\end{proof}

Source 07 proves \eqref{hset:eq:stage-agreement} for every common ordinal,
source 02 the epsilon-number clause, source 06 the cardinal clause with
$H_\kappa$, and source 09 the structure equality inside
\cref{hset:thm:extension-rigidity}. The right side of
\eqref{hset:eq:stage-agreement} must \emph{not} be replaced by ``no new short
sequences of ordinals'': a short sequence with unbounded range can be new while
every bounded subset below a large cutoff stays old, which is exactly what the
Prikry example of \cref{hset:sec:prikry} exploits. Nor is ``adds no real''
identified with ``adds no surreal'' (\cref{hset:ex:closed}).

\begin{corollary}[Countable cutoffs detect exactly new reals; source 02]
\label{hset:cor:noreals}
If $\lambda$ is a countable epsilon number, then
$\B_\lambda^M=\B_\lambda^N\iff\Pow(\omega)^M=\Pow(\omega)^N$.
\end{corollary}
\begin{proof}
For infinite $\alpha<\lambda$ a bijection between $\omega$ and $\alpha$ lies in
$M$ and remains one in $N$, so a new subset of $\alpha$ gives a new subset of
$\omega$; finite ordinals have no new subsets; and $\omega<\lambda$.
\end{proof}

\begin{lemma}[New sets imply a new ordinal subset; sources 05, 07]
\label{hset:lem:newsubset}
If $M\subsetneq N$, then $N$ has a subset of some ordinal that is not in $M$.
\end{lemma}
\begin{proof}
Suppose not, and let $a\in N$. In $N$ enumerate $\TC(\{a\})$ by an ordinal and
code the pulled-back membership relation by a subset of an ordinal through an
old injective pairing (for example the natural-arithmetic rectangle map). By
hypothesis the code, hence the relation, is in $M$. Well-founded in $N$, it is
well founded in $M$; its decoration computed in $M$ is the decoration in $N$ by
uniqueness of recursion, and its root value is $a$. So $a\in M$, and $M=N$.
\end{proof}

\begin{theorem}[First new birthday; sources 02, 05, 07]
\label{hset:thm:firstbirthday}
Let $M\subsetneq N$ and
\[
 \beta=\beta(M,N)=\min\{\alpha:\Pow(\alpha)^N\ne\Pow(\alpha)^M\}.
\]
Then $\beta$ exists and:
\begin{enumerate}[label=(\roman*)]
\item every surreal of $N$ of birthday below $\beta$ belongs to $M$;
\item $N$ has a surreal of birthday exactly $\beta$ not in $M$, all of whose
proper prefixes are in $M$; thus $\beta$ is the least birthday of an element of
$\No^N\setminus\No^M$;
\item $\beta$ is an infinite cardinal of $M$.
\end{enumerate}
\end{theorem}
\begin{proof}
$\beta$ exists by \cref{hset:lem:newsubset}. A sign sequence of birthday
$\alpha$ is determined by $\alpha$ and its plus set, which is old when
$\alpha<\beta$; this proves (i), and a new subset of $\beta$ gives the new sign
of (ii), whose proper prefixes are old by (i). No finite ordinal has a new
subset. If $\beta$ were not an $M$-cardinal, $M$ would have a bijection
$f:\nu\to\beta$ with $\nu<\beta$; for a new $A\subseteq\beta$ the pullback
$f^{-1}[A]$ would be old by minimality, and then $A=f[f^{-1}[A]]$ would be old.
\end{proof}

\begin{remark}[Source 02's version]\label{hset:rem:firstbirthday-02}
Source 02 proved the theorem for a set-forcing extension $N=M[G]\ne M$,
obtaining existence of $\beta$ from the generic filter: a ground-model
well-order of the forcing set turns $G$ into a new subset of an ordinal. Sources
05 and 07 remove the forcing hypothesis by \cref{hset:lem:newsubset}. No claim
that all cardinals are preserved is needed. The theorem does not assert that
$\beta$ is regular, and it does not identify the signs of birthday at most
$\omega$ with real numbers.
\end{remark}

\begin{example}[Adding a real]\label{hset:ex:cohen}
If $N=M[g]$ adds a Cohen subset $g$ of $\omega$, no finite subset is new and $g$
is, so $\beta=\omega$. The sign with plus set $g$ is a new surreal of birthday
$\omega$, and each of its finite prefixes is old. Such an extension changes
$\No_{<\kappa}$ for every common uncountable cardinal $\kappa$.
\end{example}

\begin{example}[No new reals, but new surreals; sources 05, 07]\label{hset:ex:closed}
Force over $M$ with countable partial functions from $\omega_1^M$ to $2$,
ordered by extension. The union of a descending $\omega$-chain is a condition,
so the forcing is countably closed; the successive-decision argument shows that
it adds no new countable sequence of ordinals, hence no subset of any old
countable ordinal, and preserves $\omega_1$. The generic union differs from
every old subset of $\omega_1^M$ by a density argument \cite{Jech}. Hence
$\beta=\omega_1^M$: all countable-birthday surreals are unchanged, there are
new signs of birthday $\omega_1^M$, and the full birthday inclusion still fails
elementarity (\cref{hset:thm:outer}). Agreement on the reals, or even on every
countable sign sequence, does not exhaust what the birthday expansion remembers.
\end{example}

\begin{remark}[The sibling report]\label{hset:rem:sibling}
The sibling report
\repo{foundations-and-computation/surreal-fields-across-universes} works in the
same setting of transitive $M\subseteq N$ with the same ordinals and uses the
same letter $\beta$ for the first new birthday; it writes $\delta$ for the first
new ordinal-sequence length, which is not named here. It studies the external
saturation of the \emph{pure} ordered field $\No^M$ inside $N$, which is
elementary in $\No^N$ by \cref{hset:prop:absolute}; the present report studies
the birthday expansion, whose inclusion is non-elementary as soon as $M\ne N$
(\cref{hset:thm:outer}). The first-new-birthday theorem and the absoluteness of
old arithmetic are shared ground; the countably closed and Prikry examples
appear in both.
\end{remark}

\section{Detecting new subsets}
\label{hset:sec:detectors}

\subsection{The old-power-set detector at a cutoff}

\begin{theorem}[Uniform detector for an added subset; source 09]
\label{hset:thm:detector}
Let $\kappa$ be an uncountable cardinal of both models. Suppose $\alpha<\kappa$,
$N$ has a subset of $\alpha$ not in $M$, and
\[
 \nu=(2^{\card\alpha})^M<\kappa.
\]
Then $\B_\kappa^M\not\prec\B_\kappa^N$. One fixed interpreted formula
witnesses the failure using parameters coding $\alpha$ and
$A=\Pow^M(\alpha)$.
\end{theorem}
\begin{proof}
The cardinal bound gives $A\in H_\kappa^M$. In the membership language
consider
\begin{equation}
 \exists y\,[\forall v\in y\ (v\in\alpha)\ \land\ y\notin A].
 \label{hset:eq:new-subset-detector}
\end{equation}
This is false in $H_\kappa^M$, because all its subsets of $\alpha$ lie in
$A$. It is true in $H_\kappa^N$: any newly added subset of $\alpha$ has
hereditary size below the still-cardinal $\kappa$ and is not an element of
the old set $A$. Translate \eqref{hset:eq:new-subset-detector} through $\I$ and
choose old codes for the two parameters. Cross-model coherence gives the
claimed failure in the birthday structures.
\end{proof}

The membership-language formula \eqref{hset:eq:new-subset-detector} is
$\Sigma_1$: its matrix has only bounded quantifiers. Its fully expanded
birthday-field translation need \emph{not} be existential or $\Sigma_1$ in
that different language, because code validity includes the bit-quantified
well-foundedness test. No quantifier-complexity claim is made across that
translation.

\begin{proposition}[Sharpness of the detector's parameter cost; source 09]
\label{hset:prop:detector-cost}
The old parameter $\Pow^M(\alpha)$ belongs to $H_\kappa^M$ exactly when
$(2^{\card\alpha})^M<\kappa$. In the rooted coding scheme, every code for
that parameter has node-domain cardinality
$\card{\TC(\{\Pow^M(\alpha)\})}^M$, which equals the power-set cardinal for
infinite $\alpha$. In that infinite case, a code of birthday cardinality no
larger than that cardinal can be chosen.
\end{proposition}
\begin{proof}
The first statement is \cref{hset:lem:Hpower} inside $M$. The exact
node-domain cardinality is \eqref{hset:eq:exact-code-size}. For infinite
$\alpha$, the transitive closure consists of the power set, its subsets,
and the ordinal data beneath $\alpha$, giving cardinality
$2^{\card\alpha}$. Choosing its initial ordinal as the node domain gives
relation birthday $\Lambda(\nu)$, of cardinality $\nu$ by
\cref{hset:lem:pairing}.
\end{proof}

This is not a theorem that elementarity must hold when the bound fails.
Another formula can detect a change at a smaller cutoff. It says exactly
when this particular, universally applicable old-power-set witness is
available.

\begin{example}[A conditional continuum-sized illustration]
Suppose $M$ satisfies CH, $N$ preserves $\omega_2^M$, and $N$ adds a real.
Then the old set of reals has hereditary size $\omega_1^M$, so the detector
works at the cutoff $\omega_2^M$. That old set is not an element of
$H_{\omega_1}^M$, so the same parameter cannot be used at the cutoff
$\omega_1^M$. No claim of elementarity at the smaller cutoff follows.
\end{example}

\subsection{Rigidity under extensions at a strong limit}

\begin{theorem}[Strong-limit extension rigidity; source 09]
\label{hset:thm:extension-rigidity}
Let $\kappa$ be an uncountable cardinal of both models, strong limit in $M$.
Then
\begin{align*}
 \B_\kappa^M\prec\B_\kappa^N
 &\iff \B_\kappa^M=\B_\kappa^N\\
 &\iff \Pow^M(\alpha)=\Pow^N(\alpha)
                 \quad\text{for every }\alpha<\kappa.
\end{align*}
\end{theorem}
\begin{proof}
If some subset of an $\alpha<\kappa$ is new, the strong-limit property
makes its old power-set parameter small enough, and \cref{hset:thm:detector}
excludes elementarity. Conversely, if all these power sets agree, the
structures are equal by \cref{hset:thm:preservation}, hence the inclusion is
elementary. Equality of the scalar domains forces equality of the
corresponding plus-sets, proving the remaining direction.
\end{proof}

This proves \cref{hset:thm:main-extension}. It applies, for example, to any
forcing extension satisfying the stated cardinal-preservation hypotheses,
but uses no particular forcing construction. It is a theorem about
\emph{standard cutoff structures at a fixed common cardinal}, not a
prohibition on abstract elementary extensions supplied by model theory.
Source 02 explicitly declined to claim that a proper enlargement at a fixed
cutoff is non-elementary merely because it is proper; the theorem above and
\cref{hset:thm:detector} give sufficient conditions, and nothing more is
claimed. Their ordered field reducts still include elementarily even when the
birthday expansion fails the detector; a set-theoretic change can be invisible
to ordered-field formulas yet visible to one birthday-field formula with old
parameters.

\subsection{The full class, and an explicit row-search formula}

For the full classes write $\B^M=(\No^M;0,1,+,-,\cdot,<,\bd)$ and likewise
$\B^N$; by \cref{hset:prop:absolute} the old signs form a substructure.

\begin{theorem}[An outer universe is detectable; sources 05, 07]
\label{hset:thm:outer}
$\B^M\prec\B^N\iff M=N$. When $M\ne N$ the failure is witnessed with old
parameters: by the $\I$-translation of \eqref{hset:eq:new-subset-detector} with
codes for $\alpha$ and $\Pow(\alpha)^M$, for any $\alpha$ with a new subset
(source 07), and by the formula $\New$ of \cref{hset:thm:new} (source 05).
\end{theorem}
\begin{proof}
If $M\ne N$, \cref{hset:lem:newsubset} gives an ordinal $\alpha$ with a new
subset. Old valid codes remain valid in $N$ and decode to the same sets (as in
\cref{hset:prop:cross-model}, without a cardinal bound), and
\eqref{hset:eq:new-subset-detector} is false in $M$ and true in $N$ with those
parameters.
\end{proof}

Source 05's detector reads the birthday language directly. Fix a nonzero
ordinal $\theta$ and in $M$ an enumeration $(A_j)_{j<\nu}$, $\nu>0$, of
$\Pow(\theta)^M$ (a surjection suffices), and code it by one old sign $m$ of
birthday $\Lambda(\theta+\nu+1)$ whose bit at $p_{\theta+\nu+1}(j,\xi)$ is plus
exactly when $\xi\in A_j$. Let
\[
 \New(m,\theta,\nu):\quad
 \exists x\Bigl[\bd(x)=\theta\ \land\
 \forall j<\nu\ \exists\xi<\theta\
 \neg\bigl(\Bit(x,\xi)\leftrightarrow
 \Bit(m,p_{\theta+\nu+1}(j,\xi))\bigr)\Bigr],
\]
together with the ordinal and length conditions on $\theta,\nu,m$; it asks for a
row of length $\theta$ different from every old row.

\begin{theorem}[An explicit new-subset detector; source 05]\label{hset:thm:new}
With such parameters, $\B^M\models\neg\New(m,\theta,\nu)$, and
$\B^N\models\New(m,\theta,\nu)$ if and only if $\Pow(\theta)^N\ne\Pow(\theta)^M$.
\end{theorem}
\begin{proof}
Every sign of birthday $\theta$ in $M$ has an old plus set, which is some row
$A_j$. The old sign $m$ has the same bits in $N$, so it still has exactly the
old rows. A new subset of $\theta$ gives a witness in $N$; conversely a witness
has a plus set differing from every old subset, hence new.
\end{proof}

Source 05 writes the matrix with its rectangle map and length $\theta(\nu+1)$;
the formula is the same up to that pairing. The enumeration is a set in $M$ and
the matrix one sign of set length; no class of ground-model signs is named. The
formula is not claimed to be existential after the abbreviations are expanded,
since the prefix and bit predicates have quantifiers of their own. The
detector branches off after the pairing lemma and does not need the arithmetic
certificates.

\begin{remark}[The bounded version]\label{hset:rem:bounded-new}
At a cutoff the matrix must fit. Source 05 states: if a common ordinal $\kappa$
is an uncountable regular cardinal in both models, the old matrix has
$\theta,\nu<\kappa$ and length below $\kappa$, and $N$ has a new subset of
$\theta$, then $\B^M_\kappa\not\prec\B^N_\kappa$, witnessed by $\New$. Its proof
uses regularity only through closure of natural arithmetic below $\kappa$, which
holds at every uncountable cardinal (\cref{hset:lem:card-estimates}); so the
conclusion holds whenever $\kappa$ is an uncountable cardinal of both models and
$(2^{\card\theta})^M<\kappa$, the hypothesis of \cref{hset:thm:detector}.
Without a bound on the old enumeration the corollary cannot be inferred from
$\theta<\kappa$ alone, because the set of old subsets of a small ordinal may
have cardinality at least $\kappa$.
\end{remark}

\begin{remark}[The conclusions use parameters]
These theorems say that the \emph{canonical inclusion} is not elementary. They
do not by themselves show that the two structures have different
parameter-free complete theories, nor do they exclude unrelated embeddings or
isomorphisms. If the extension changes a parameter-free set-theoretic sentence,
its translation does distinguish the theories (sources 05, 07).
\end{remark}

\section{Prikry forcing: an unchanged field with a singularized cutoff}
\label{hset:sec:prikry}

\subsection{The imported forcing theorem}

We use the classical ordinary Prikry forcing associated with a normal measure
$U$ on a measurable cardinal $\kappa$. Conditions consist of a finite increasing
stem $s\subset\kappa$ and a measure-one set $A\in U$ above that stem; an
extension may lengthen the stem using points of $A$ and shrink $A$, a direct
extension only shrinks $A$. Its standard properties are: the generic stem union
is cofinal of order type $\omega$ in $\kappa$; the forcing adds no bounded
subsets of $\kappa$; and it preserves cardinals. These are imported from the
exposition of Poveda--Rinot--Sinapova, especially Lemma 2.10 and Section 3.1
\cite{PRS}, and are not reproved in full. For orientation, source 02 sketches
the bounded-subset consequence of the Prikry property: decide a name for a
subset of $\alpha<\kappa$ by successive direct extensions, using that
intersections of fewer than $\kappa$ measure-one sets have measure one; bounded
subset preservation keeps the cardinals below $\kappa$ and $\kappa$ itself, and the
$\kappa^+$-chain condition the cardinals above.

\begin{theorem}[Forcing invisibility, conditional; sources 02, 06]
\label{hset:thm:prikry}
Suppose a transitive model $M$ of ZFC has a measurable cardinal $\kappa$, and let
$N=M[G]$ be an ordinary Prikry extension at $\kappa$. Then
\[
 \B_\kappa^M=\B_\kappa^N,\qquad H_\kappa^M=H_\kappa^N=V_\kappa^M,
\]
with the same carrier, functions and relations, although
\[
 M\models\cf(\kappa)=\kappa,
 \qquad N\models\cf(\kappa)=\omega.
\]
The common membership structure satisfies ZFC in both universes.
\end{theorem}
\begin{proof}
In $M$ the measurable $\kappa$ is strongly inaccessible, hence an epsilon number
by \cref{hset:lem:card-estimates}; it remains a cardinal, hence an epsilon
number, in $N$, although singular there. Prikry forcing adds no subset of any
$\alpha<\kappa$, so \cref{hset:thm:preservation} gives both equalities of
structures. Since $\kappa$ is inaccessible in $M$, $H_\kappa^M=V_\kappa^M$ (sets in
$V_\kappa$ have small transitive closures by strong limit and regularity, and a
well-founded relation on fewer than $\kappa$ nodes has rank below $\kappa$), and
$V_\kappa^M\models\ZFC$. The generic sequence witnesses $\cf^N(\kappa)=\omega$.
Truth of each fixed first-order formula in one unchanged set structure is
absolute between transitive models, by induction on formulas.
\end{proof}

\begin{corollary}[No uniform regularity detector; sources 02, 06]
\label{hset:cor:prikry}
Relative to the existence of such a pair of models, no first-order sentence in
the language of $\B_\kappa$ holds, in every universe and at every uncountable
cardinal cutoff, exactly when the external cutoff is regular, or exactly when it
is strongly inaccessible. More strongly, no invariant determined by the
isomorphism type of the enriched structure alone determines the external
cofinality of its cutoff in a universe-independent way.
\end{corollary}
\begin{proof}
Such a sentence or invariant would take different values on the two identical
structures of \cref{hset:thm:prikry}.
\end{proof}

This is a relative-consistency separation using a measurable cardinal: no
measurable cardinal is constructed or proved to exist, and ZFC is not claimed to
prove the existence of such a pair of universes. The statement concerns fixed
first-order sentences; unrestricted second-order satisfaction is \emph{not}
asserted to be unchanged, since the family of external subsets of the carrier
changes. The equality established in $N$ is with the \emph{old} $V_\kappa^M$;
it is not claimed that $H_\kappa^N=V_\kappa^N$, and the rank hierarchy of the
extension can contain new objects. The corollary also shows why diagnosing
singularity by failure of some unspecified Replacement instance is unreliable:
the small universe keeps every axiom of ZFC while the ambient universe acquires
an unbounded sequence invisible to it.

\begin{proposition}[Where the missing sequence lives; sources 02, 06]
\label{hset:prop:onset}
In the situation of \cref{hset:thm:prikry}, the first new birthday is exactly
$\beta(M,N)=\kappa$. The Prikry sequence $h:\omega\to\kappa$, and its range, have
hereditary size $\kappa$, so they lie in $H_{\kappa^+}^N\setminus H_\kappa^N$; any
faithful pointed code for them needs at least $\kappa$ nodes. For every epsilon
number $\lambda$ with $\kappa<\lambda<\kappa^+$, of which there are some, the
carrier $\No_{<\lambda}$ changes and its interpreted universe $H_{\kappa^+}$
contains the new sequence.
\end{proposition}
\begin{proof}
No subset below $\kappa$ is new, while the range of $h$ is a new subset of
$\kappa$ (already by regularity of $\kappa$ in $M$); \cref{hset:thm:firstbirthday}
gives $\beta=\kappa$, and every $\lambda>\kappa$ contains that sign layer. The
values of $h$ are cofinal ordinals below $\kappa$, so the transitive closure of
$h$ contains all of $\kappa$ and has size $\kappa$. For $\kappa<\lambda<\kappa^+$ the
cardinal ceiling is $\kappa^+$ in both models. Iterating $\alpha\mapsto
\omega^\alpha$ $\omega$ times from above $\kappa$ gives a fixed point of
cardinality $\kappa$, an epsilon number in the interval.
\end{proof}

The proposition distinguishes the \emph{appearance of new elements} from a
\emph{change of the complete first-order theory}; it proves the former and does
not assert the latter for every enlarged cutoff. By contrast, an extension
adding a real changes $\No_{<\kappa}$ at every common uncountable cardinal
(\cref{hset:ex:cohen}); the preservation criterion distinguishes the two kinds
of extension without referring to the names of their forcing notions. The
surcomplex enrichment is likewise unchanged in the Prikry situation
(\cref{hset:thm:complex-forcing}).

\part{The full class and the surcomplex numbers}\label{hset:part:class}

\section{The full class: the announced theorem and Kunen rigidity}
\label{hset:sec:fullclass}

\subsection{The announced global bi-interpretation}

\begin{theorem}[Chen--Hamkins--Yang, announced \cite{CHYKobe,CHYNotreDame,CHYCUNY}]
\label{hset:thm:chy}
In ZFC there are parameter-free first-order interpretations $\I$ and $\J$ such
that, formula by formula, $\I(\B)$ is definably identified with $(V,\in)$,
$\J(V,\in)$ with $\B$, and both composites have definable isomorphisms to the
original structures.
\end{theorem}

\paragraph{Status and credit.} This is the global theorem announced by Chen,
Hamkins and Yang, stated there for the surreal field with birthday order; by
\cref{hset:lem:birthday-language} that language and the birthday-function
language are definitionally equivalent. Their slides describe the coding of sets
by pointed well-founded extensional relations and then by sign sequences, and
mark the work as still being checked; the collection's foundations report
records the announcement in \lbl{found:sub:announcement}. The written argument
below is source 07's; source 05 gives the same statement with its bisimulation
codes, and internally in every model of ZFC (\cref{hset:rem:internal}). It is
printed as \emph{a proof of the announced theorem}, not as a new result, and it
is not an audit of an unpublished manuscript of Chen--Hamkins--Yang or of their
proposed axiomatization.

\begin{proof}[Proof (source 07)]
Use the codes of \cref{hset:sec:codes} without a cutoff. Every set has a valid
code: choose an injection of $\TC(\{a\})$ into an ordinal (one application of
Choice) and code the transported membership relation. The correctness of
$\sim$ and $\in_\I$ (\cref{hset:lem:equal-member}) and of the collapse is proved
for set-sized graphs and applies to each code. The reverse interpretation takes
sign names $\langle\alpha,A\rangle$ with $A\subseteq\alpha$ and defines the
arithmetic by the certificates of \cref{hset:sec:certificates}, which now exist
without a size restriction. The two composites are the collapse relation and the
sign-name relation $x\mapsto\langle\bd(x),A_x\rangle$ of
\cref{hset:sec:biinterpretation,hset:sec:bisim}; both are single formulas, and
their correctness is proved one formula at a time, with no truth predicate for
$V$ and no class of chosen representatives.
\end{proof}

The local refinements proved in \cref{hset:part:spine} were stated for $V$ by
the announcement; the cutoff theorems, the noncardinal pointed version, and the
local certificates at singular cutoffs are what sources 02, 05, 06, 07 and 09
propose to add, with the priority caveats of \cref{hset:sec:scope}. Source 09
already noted that its refinements may overlap unpublished details of the
Chen--Hamkins--Yang project.

\subsection{Transfer of Kunen's inconsistency}
\label{hset:sub:kunen}

An automorphism and a non-surjective elementary self-embedding must not be
conflated. For automorphisms of the full class the argument of
\cref{hset:lem:rigidity} applies (the ordinals form a definable well-ordered
class). It fails for a non-surjective elementary self-embedding $j$, whose
restriction to the ordinals need not be an automorphism, and which satisfies
$\bd(j(x))=j(\bd(x))$ rather than the externally birthday-fixed equation
$\bd(j(x))=\bd(x)$. Literature on birthday-fixed surreal morphisms, such as the
structures of Rangel--Mariano \cite{RM}, must be compared with this distinction
in mind (source 07).

\begin{lemma}[Transferring an elementary embedding; source 07]
\label{hset:lem:embedding-transfer}
Let $j:\B\to\B$ be a class elementary embedding. For each set $a$ choose any
code $C$ with $x_C=a$ and put $\widehat j(a)=x_{j(C)}$, with $j$ applied
coordinatewise. This is a well-defined class elementary embedding
$(V,\in)\to(V,\in)$, and $\widehat j$ is the identity only if $j$ is.
\end{lemma}
\begin{proof}
Elementarity preserves validity and $\sim$, so the value does not depend on the
code, and $\widehat j$ is a class defined from $j$ without choosing
representatives. For each set-theoretic formula and codes for its parameters,
\cref{hset:thm:chy} and elementarity of $j$ give
$\varphi(\vec a)\iff\varphi(\widehat j(\vec a))$. The definable sign-name
composite is natural under elementary maps; if $\widehat j$ is the identity, it
identifies $j(x)$ and $x$ with the same interpreted sign name, so $j(x)=x$.
\end{proof}

\paragraph{The class-theoretic hypothesis.} The classical Kunen inconsistency
rules out a nonidentity elementary embedding $V\to V$ under Choice, in the form
that permits the embedding as a class parameter in Separation and Replacement
(Hamkins--Kirmayer--Perlmutter \cite{HKP}). We work, as source 07 does, in GBC
for convenience; global choice is stronger than needed. The class $\widehat j$
is definable from $j$ and inherits availability as a class parameter.
Elementarity may be read formula by formula; for a single class formulation one
may use $\Delta_0$-elementarity with membership-cofinality, which holds because
$\widehat j(\alpha)\ge\alpha$ and $\widehat j(V_\alpha)=V_{\widehat j(\alpha)}$.

\begin{theorem}[No nontrivial class elementary self-embedding; source 07]
\label{hset:thm:kunen}
In GBC, or under the class-parameter Separation and Replacement hypotheses just
described, every class elementary self-embedding of $\B$ is the identity.
\end{theorem}
\begin{proof}
A nonidentity $j$ would give a nonidentity elementary $\widehat j:V\to V$ by
\cref{hset:lem:embedding-transfer}, contradicting the cited Kunen theorem.
\end{proof}

The set-theoretic input is the classical Kunen inconsistency; this is not a new
proof of it. The theorem concerns the full birthday-expanded class. It does not
assert that the pure surreal ordered field has no nontrivial embeddings or
automorphisms; it does not classify the elementary self-embeddings of the bounded
$\B_\kappa$ (which \cref{hset:thm:embeddings} reduces to those of $H_\kappa$); and
it does not prohibit external self-maps of a countable nonstandard model's
internal surreal structure, which need not be classes of that model with the
required Separation and Replacement properties.

\section{Surcomplex numbers and the orientation obstruction}
\label{hset:sec:surcomplex}

\subsection{The enriched surcomplex structures}

For an epsilon number $\lambda$, $\No_{<\lambda}[i]$ is algebraically closed
because $\No_{<\lambda}$ is real closed. Let $c(x+iy)=x-iy$ be conjugation and
\[
 \cb(x+iy)=\max(\bd(x),\bd(y)),
\]
the \emph{coordinate birthday}, whose value is an embedded real ordinal. Put
\[
 \C_\lambda=(\No_{<\lambda}[i];0,1,+,-,\cdot,c,\cb),
 \qquad
 \C^i_\lambda=(\C_\lambda,i),
\]
and $\C$, $\C^i$ for the full classes on $\No[i]$. The coordinate birthday is a
stated convention, not the birthday of a combinatorial-game representative of a
complex object, and not asserted to be a pre-existing universal definition of a
``surcomplex birthday''. Conjugation, $\cb$ and $i$ are explicit extra
structure; none is asserted to be definable in the pure algebraically closed
field, whose inclusions are elementary in the pure language \cite{Marker}.

\begin{lemma}[Independence of orientation; source 06]
\label{hset:lem:cb-orientation}
$\cb$ is unchanged if $i$ is replaced by $-i$, and its restriction to
$\No_{<\lambda}$ is $\bd$.
\end{lemma}
\begin{proof}
The replacement changes coordinates $(x,y)$ to $(x,-y)$, and negation reverses
every sign without changing length, so $\bd(-y)=\bd(y)$; also $\bd(0)=0$.
\end{proof}

\begin{proposition}[Recovery of the real birthday structure; all five sources]
\label{hset:prop:real-recovery}
Both $\C_\lambda$ and $\C^i_\lambda$ interpret $\B_\lambda$ without parameters on
the fixed field of $c$, with order
\[
 x<y\quad\Longleftrightarrow\quad
 \exists t\,[c(t)=t\land t\ne0\land t^2=y-x]
\]
for real $x,y$. The structure $\C^i_\lambda$ is uniformly bi-interpretable with
$\B_\lambda$, hence with $(H_\lambda,\in)$ at cardinal cutoffs and with
$(H_{\kappa(\lambda)},\in,\lambda)$ otherwise; the same holds for $\C^i$, $\B$ and
$(V,\in)$ formula by formula.
\end{proposition}
\begin{proof}
The fixed field is $\No_{<\lambda}$, real closed, where positivity is being a
nonzero square; field operations restrict, and $\cb$ restricts to $\bd$.
Conversely interpret complex numbers as pairs of reals with
$(x,y)(u,v)=(xu-yv,xv+yu)$, $c(x,y)=(x,-y)$, $i=(0,1)$ and
$\cb(x,y)=(\max(\bd x,\bd y),0)$. With $i$ named, $(x,y)\mapsto x+iy$ is
definable with inverse
\[
 z\longmapsto\Bigl(\frac{z+c(z)}2,\frac{z-c(z)}{2i}\Bigr),
\]
and the real composite is $x\mapsto(x,0)$. These are the definable isomorphisms
of a bi-interpretation; compose with \cref{hset:thm:biinterpretation} and
\cref{hset:thm:chy}.
\end{proof}

Thus cardinal rounding, recovery of a noncardinal cutoff, and all the
set-theoretic distinctions above transfer to the oriented surcomplex
structures, and adjoining $i$ creates no new set-theoretic size requirement.
The real-axis recovery is standard quadratic-extension algebra; its role here is
to identify which enrichment carries the foundational information. In the
Prikry situation of \cref{hset:thm:prikry}, $\C_\kappa$ and $\C^i_\kappa$ are also
unchanged, because both coordinates and all displayed operations are; so
adjoining $i$ and retaining conjugation does not repair the inability to recover
the ambient cofinality (sources 02, 06).

\subsection{Rigidity and the two-element automorphism group}

\begin{lemma}[Rigidity of a standard birthday cutoff; sources 05, 06, 07, 09]
\label{hset:lem:rigidity}
Every automorphism of $\B_\lambda$ is the identity, and every automorphism of a
transitive structure $(H_\kappa,\in)$, or $(H_\kappa,\in,\lambda)$, is the
identity.
\end{lemma}
\begin{proof}
The embedded ordinals are definable and numerically well ordered, so any order
automorphism of this segment fixes every ordinal by transfinite induction; the
automorphism therefore fixes all birthdays, preserves the definable bit
predicate at each fixed coordinate, and fixes every sign sequence. For
$H_\kappa$, external well-founded induction gives $h(x)=\{h(y):y\in x\}=x$;
naming an ordinal does not change this.
\end{proof}

The first assertion is a case of the collection's existing simplicity
rigidity: the report \repo{surcomplex/surcomplex-field-automorphisms} shows, in
its section \lbl{saut:sec:exponential} (subsection ``Simplicity restores the
canonical surreal numbers''), that a bijection preserving numerical order and
prefix simplicity is the identity, and prefix is definable from order and
birthday (\cref{hset:lem:prefix}). It is printed here once, not as a new result.
Rigidity is not pointwise definability: a rigid structure can have elements
that are not individually definable without parameters.

\begin{theorem}[The orientation obstruction; sources 05, 06, 09, after Je\v{r}\'abek]
\label{hset:thm:orientation}
For every epsilon number $\lambda$, $\Aut(\C_\lambda)=\{\id,c\}$. Consequently
$\C_\lambda$ is parameter-free mutually interpretable with, but not
parameter-free bi-interpretable with, $(H_\lambda,\in)$ at cardinal cutoffs, or
$(H_{\kappa(\lambda)},\in,\lambda)$ otherwise; after naming $i$ the
bi-interpretation holds (\cref{hset:prop:real-recovery}).
\end{theorem}
\begin{proof}
An automorphism preserves the fixed field of $c$, its definable order and
birthday, so its restriction is the identity by \cref{hset:lem:rigidity}; it
sends $i$ to a root of $X^2+1$, so it is the identity or conjugation, as in
\lbl{saut:thm:axis} ($\Aut(F(i)/F)=\{\id,c\}$ for real closed $F$). Both
preserve $c$ and $\cb$ (\cref{hset:lem:cb-orientation}).

A parameter-free bi-interpretation induces mutually inverse homomorphisms
between the automorphism groups: an automorphism preserves the definable
interpretation domain and equivalence relation, induces an automorphism of the
interpreted structure, and the definable round-trip isomorphism recovers it. The
set structure is rigid (\cref{hset:lem:rigidity}), whereas conjugation is a
nonidentity automorphism of $\C_\lambda$. This excludes \emph{every}
parameter-free bi-interpretation, not only the coordinate one. Mutual
interpretation holds: recover the real birthday structure inside $\C_\lambda$ and
then its set universe; in the other direction use the reverse surreal
interpretation and the pair field, omitting the constant $i$.
\end{proof}

The general real/complex obstruction, that real and complex fields with
conjugation are mutually interpretable but need a named imaginary unit for a
parameter-free bi-interpretation, was explained by Je\v{r}\'abek in 2021
\cite{Jerabek}; only source 06 credited it, and the credit applies to the
versions in sources 05 and 09 as well. The group computation is a consequence of
classical rigidity and of \lbl{saut:thm:axis}, not a new classification of the
unrestricted field automorphisms of $\No[i]$ studied in the automorphisms report.
The point here is the distinction between mutual interpretability and
bi-interpretability, and why an unnamed imaginary unit cannot be treated as
harmless notation in the latter claim. The two roots $i$ and $-i$ have the same
parameter-free type in $\C_\lambda$; interpreting a rigid universe does not force
the interpreting structure to be rigid, and it is the definability of both round
trips that transmits rigidity (source 06).

\subsection{Two lifts of each real embedding}

\begin{theorem}[Two lifts, and only two; source 06]\label{hset:thm:two-lifts}
Let $\kappa,\lambda$ be uncountable cardinals. Every elementary embedding
$f:\B_\kappa\to\B_\lambda$ has exactly two elementary extensions
$\C_\kappa\to\C_\lambda$,
\[
 \widehat f_\epsilon(x+iy)=f(x)+\epsilon\,i\,f(y),\qquad\epsilon\in\{1,-1\},
\]
relative to reference roots of $-1$, and every elementary embedding
$\C_\kappa\to\C_\lambda$ is one of these. After naming the reference imaginary
units only the positive lift remains. Hence
\[
 \Emb(\C_\kappa,\C_\lambda)\simeq\Emb(H_\kappa,H_\lambda)\times\{1,-1\},
 \qquad
 \Emb(\C^i_\kappa,\C^i_\lambda)\simeq\Emb(H_\kappa,H_\lambda),
\]
the signs multiplying under composition, and $\C^i_\kappa\prec\C^i_\lambda\iff
H_\kappa\prec H_\lambda$ for $\kappa<\lambda$ (source 05 proved the last
equivalence, and the positive-lift correspondence, at regular cardinals).
\end{theorem}
\begin{proof}
The displayed maps are field embeddings commuting with $c$ and preserving
$\cb$, because $f$ preserves order and birthdays, and they differ on $i$. For
elementarity, name $i$ and translate each complex formula into $\B_\kappa$ by
coordinates, as in \cref{hset:prop:real-recovery}; elementarity of $f$ gives
elementarity of $\widehat f_1$, and composing with $c$ gives the other lift in
the language without $i$. Conversely an elementary embedding of the complex
structures restricts to an elementary embedding of the definable real
structures and sends $i$ to $\pm i$, and the field laws force the displayed
form. Combine with \cref{hset:thm:embeddings}.
\end{proof}

At $\kappa=\lambda$ with $f=\id$ this gives $\Aut(\C_\kappa)=\{\id,c\}$ again. The
statement concerns the specified birthday-and-conjugation enrichment; it is
emphatically not a classification of automorphisms of the pure algebraically
closed field, which has far more symmetry (see the automorphisms report). It is
a specialization combining the local theorem with coordinate complexification,
not a new general Galois principle.

\subsection{The exact parameter threshold for orientation}

\begin{theorem}[Orientation threshold; source 06]\label{hset:thm:threshold}
Let $\kappa$ be an uncountable cardinal and $A\subseteq\No_{<\kappa}[i]$ any set of
parameters. The following are equivalent:
\begin{enumerate}[label=(\roman*)]
\item some root of $X^2+1$ is definable in $\C_\kappa$ over $A$;
\item a linear order of the whole carrier is definable in $\C_\kappa$ over $A$;
\item $A$ contains a nonreal element.
\end{enumerate}
In particular neither an imaginary unit nor a linear order is definable over any
collection of real parameters, and one nonreal parameter suffices for both.
\end{theorem}
\begin{proof}
If all parameters are real, conjugation fixes them and exchanges the two roots,
so neither root is definable; and a definable linear order $\triangleleft$
would be preserved by $c$, while for nonreal $z$ exactly one of
$z\triangleleft c(z)$ and $c(z)\triangleleft z$ holds and $c$ reverses it. If
$z\in A$ is nonreal, $v=z-c(z)$ is nonzero and purely imaginary, $-v^2$ is real
and positive, and
\[
 j=\frac{z-c(z)}{\sqrt{-(z-c(z))^2}},
\]
with the positive square root of the definable real field, is definable from
$z$ and satisfies $j^2=-1$. The definable coordinates
$\bigl(\tfrac{w+c(w)}2,\tfrac{w-c(w)}{2j}\bigr)$, ordered lexicographically,
give a definable linear order of the carrier; it is not a field order.
\end{proof}

Each defining formula uses finitely many parameters; allowing arbitrarily many
real parameters does not help because conjugation fixes them all. The proof uses
only that $\No_{<\kappa}$ is real closed. A nonreal parameter removes the twofold
obstruction to the coordinate round trip, but if it is kept as a named constant,
its two real coordinates must be kept in the real interpretation and the
corresponding sets in the $H_\kappa$ interpretation. It is \emph{not} inferred
that an arbitrary expansion $(\C_\kappa,z)$ is bi-interpretable with the
\emph{unexpanded} $H_\kappa$; naming a root of $-1$ avoids the issue because its
coordinates $(0,\pm1)$ are already definable.

\subsection{The full class, under GBC}

\begin{theorem}[Surcomplex elementary self-embeddings; source 07]
\label{hset:thm:complex-rigid}
In the class setting of \cref{hset:thm:kunen}, the class elementary
self-embeddings of $\C$ are exactly the identity and conjugation, and those of
$\C^i$ only the identity. In particular every such self-embedding is surjective.
\end{theorem}
\begin{proof}
An elementary self-embedding of $\C$ restricts to one of the definable real
structure $\B$, which is the identity by \cref{hset:thm:kunen}; it sends $i$ to
$\pm i$, and the coordinates are fixed, so it is the identity or conjugation.
Both are automorphisms. Naming $i$ excludes conjugation.
\end{proof}

This is the one surcomplex result of this report that goes beyond the
automorphism statements of the automorphisms report: it concerns elementary
self-embeddings that are not assumed surjective, of the full structure, under
GBC. With \cref{hset:lem:rigidity} in place of Kunen's theorem the same argument
gives only the two-element \emph{automorphism} group of the bounded structures.
Nothing is asserted about unrestricted field embeddings, and the specified
conjugation and coordinate birthday are not claimed to be recoverable from the
pure surcomplex field.

\subsection{Outer models on the surcomplex side}

\begin{theorem}[The surcomplex outer-model contrast; sources 05, 07]
\label{hset:thm:complex-forcing}
For transitive $M\subseteq N$ with the same ordinals, the inclusion of the full
surcomplex fields is elementary in the pure field language, and also after naming
$i$ and conjugation. After additionally naming $\cb$, with or without $i$, it is
elementary exactly when $M=N$. If $M\ne N$, the least coordinate birthday of a
new surcomplex number is $\beta(M,N)$.
\end{theorem}
\begin{proof}
The pure-field statement is \cref{hset:prop:absolute}. With $i$ and $c$ named,
translate every formula into real coordinates; it becomes an ordered-field
formula, so the inclusion is elementary. With $\cb$ named the real axis with its
birthday is definable, so an elementary inclusion would restrict to
$\B^M\prec\B^N$, which forces $M=N$ by \cref{hset:thm:outer}. If
$\cb(x+iy)<\beta$, both coordinates are old; a new real sign of birthday $\beta$
is a new surcomplex number of coordinate birthday $\beta$.
\end{proof}

The result concerns field and coordinate structure, not complex analysis: it
involves no contour integral, topology or total surcomplex exponential, and is
not a consequence of an analytic continuation or a valuation argument.

\part{Scope, questions and provenance}\label{hset:part:scope}

\section{What the theorems do and do not imply}
\label{hset:sec:scope}

\subsection{Logical strength versus field strength}

Already $\B_{\varepsilon_0}$ recovers $H_{\omega_1}$. Within that membership
structure, the natural numbers and all subsets of the natural numbers form the
usual full two-sorted structure of second-order arithmetic
(\cref{hset:prop:arithmetic}). The birthday-field interpretation thus includes
quantification over all actual subsets of $\omega$, not merely over computable
subsets, because the numbers coding those subsets are first-order elements.
This does not make set theory definable in the pure ordered-field reduct, where
birthday is not definable (\cref{hset:prop:nondef}), and it does not give a
decision procedure (\cref{hset:cor:undecidable}).

\subsection{Four statements that must not be conflated}

For a cutoff $\lambda$ the following can have very different answers:
\begin{center}
\begin{tabular}{@{}P{0.31\textwidth}P{0.61\textwidth}@{}}
\toprule
Question & Exact issue in this report\\
\midrule
Is $\No_{<\lambda}$ a real closed field? & Classical epsilon-number cutoff theorem.\\[0.35em]
Which sets can be decoded? & Exactly $H_{\kappa(\lambda)}$, by rooted relation codes.\\[0.35em]
Does the decoded universe satisfy ZFC? & At cardinal cutoffs, exactly the worldly condition.\\[0.35em]
Is an initial inclusion elementary? & Both cutoffs must be cardinals, followed by elementarity of their hereditary-size universes.\\
\bottomrule
\end{tabular}
\end{center}
An algebraic closure theorem alone does not answer any of the last three
questions. Conversely, a transitive set model satisfying ZFC need not arise
from a regular inaccessible cutoff (\cref{hset:prop:singular-worldly}), and the
same enriched field can sit below a regular and a singular cutoff in two
universes (\cref{hset:thm:prikry}). Hereditary size, not rank, is what the
birthday structure measures: the recovered universe is $H_\kappa$, which in
general differs from $V_\kappa$; the two agree, for example, at every worldly
$\kappa$ (\cref{hset:lem:worldly-HV}).

\subsection{Choice and quotients}

Ordinary Choice is used to enumerate an individual transitive closure and, in
some proofs, to enumerate a particular fiber or choose a particular bijection.
The interpretation does not choose one representative simultaneously for every
set. It keeps all valid codes and quotients by a definable equivalence relation.
This is why Global Choice is unnecessary outside \cref{hset:sub:kunen}.

The exact cardinal-rounding conclusion is stated in ZFC. Without Choice, not
every transitive closure is well-orderable, and a code on an ordinal cannot
represent every set; \cref{hset:thm:choice} identifies the exact range of the
full-class coding in ZF. No choice-free version of the cutoff theorems is
claimed, and nothing in the ZFC proofs justifies silently deleting Choice.

\subsection{Set models, nonstandard models, and class truth}

At a fixed cutoff, the structures in the main theorems are actual set
structures with ordinary first-order satisfaction; their interpretations are not
hypothetical syntactic models satisfying only a fragment of surreal arithmetic.
A nonstandard model of ZFC can internally perform the same constructions, but
its coded well-foundedness and its internal cardinal ceiling must be read in that
model; an externally ill-founded model is not covered by an external
Mostowski-collapse claim merely because it internally believes its code is well
founded. Arbitrary nonstandard models of a proposed surreal-arithmetic axiom
system need a separate analysis.

For the full class $\No$, one applies the translation to a fixed formula and
fixed set parameters. Statements about ``all formulas'' are metatheoretic
schemes. We do not form a set of all class functions, a proper-class membership
quotient, or a truth predicate for $V$ inside ZFC. The cutoff results do not
require a class theory, a Grothendieck universe, or a large-cardinal existence
axiom; GBC is used only in \cref{hset:sub:kunen} and
\cref{hset:thm:complex-rigid}, a measurable cardinal only in
\cref{hset:sec:prikry}, and an inaccessible only in
\cref{hset:prop:singular-worldly}.

\subsection{Why uncountability and the epsilon hypothesis matter}

For the hereditary-size bi-interpretation at cardinal cutoffs, uncountability
ensures that countably sized computation certificates and the ordinary infinite
set $\omega$ remain below the cutoff. At $\omega$ the finite-birthday surreals
form the dyadic ring, not a field, and only the forward interpretation onto the
hereditarily finite sets is claimed (\cref{hset:prop:ring}). For noncardinal
cutoffs, the epsilon hypothesis supplies closure of the bounded pairing terms
and field operations. Arbitrary ordinal sign cutoffs that fail those closure
conditions are not classified; the full birthday-field statement would then need
a different language or partial operations. The reconstruction also applies to
the canonical birthday function only, and uses the presence of every sign of
each length below the cutoff: a small subfield with an inherited birthday may
omit most subset codes, and no such generalization is asserted (source 05). For
outer models all ordinals are assumed common; models of different heights are a
different problem.

\subsection{Non-claims of the sources}
\label{hset:sub:nonclaims}

Every limitation, disclaimer and priority caveat stated by a source is kept.
Most appear at their point of use above; this subsection lists them by source.

\paragraph{Across the report.} The proofs are written mathematical arguments.
They have not been independently refereed and have not been formalized in Lean;
no Lean source is supplied. The finite programs check finite analogues only. No
named longstanding conjecture is resolved. Priority is not certified by any of
the five targeted literature and repository searches; ``not found'' is a
bounded search result, not proof of absence.

\paragraph{Source 02.}
The full-class interpretation is Chen--Hamkins--Yang's prior announced work; its
contributions are the bounded refinement and its consequences, with priority not
certified. Interpreting $H_{\kappa(\lambda)}$ is not interpreting the ambient
universe and does not assert $H_{\kappa(\lambda)}\models\ZFC$. The translated
Power Set sentence is one axiom; it does not characterize Replacement,
regularity or full ZFC. Its own elementary-inclusion conditions were necessary
only (the converse is now \cref{hset:thm:elementary}). The results concern
concrete externally well-founded structures. Sharpness is intrinsic to the
interpretation. The collector recovers an ordinal value, not a surreal of
birthday $\lambda$. Prikry forcing is imported; no measurable cardinal is
constructed; the statement is relative consistency and concerns fixed
first-order sentences, not second-order satisfaction. The preservation criterion
must not be replaced by ``no new short sequences'', and a proper enlargement at
a fixed cutoff is not claimed non-elementary merely because it is proper. The
onset proposition concerns new elements, not a change of complete theory. In
the surcomplex part $i$, $c$ and the coordinate birthday are explicit structure,
not definable in the pure algebraically closed field, and no new automorphism
classification is claimed. Only one-way interpretation was claimed at ring
cutoffs and by source 02 generally. The Rangel--Mariano program is not settled.
The finite tests do not prove the transfinite results.

\paragraph{Source 05.}
Its package was called candidate-original: not independently refereed or
formally verified; no named problem settled; the audit targeted; priority
unverified. It omitted the Chen--Hamkins--Yang credit, which is restored here.
Classical inputs (sign arithmetic, cardinal closure, well-founded recursion,
quantifier elimination, and the simplicity rigidity of the automorphisms report)
are credited. Regularity was its uniform hypothesis and it inferred no singular
theorem; $H_\kappa\ne V_\kappa$ in general and $\Pow(\alpha)$ need not be a set in
$H_\kappa$. The coordinate birthday is a convention, not a game birthday. No
sharp product-birthday bound is used. The outer-model theorem needs transitive
same-ordinal models; nonstandard models are handled internally. Its
non-elementarity is witnessed with parameters, not by distinct parameter-free
theories. The first new birthday is not asserted regular, and signs of birthday
at most $\omega$ are not identified with reals. The bounded detector needs an
enumeration bound. No small-subfield generalization and no different ordinal
heights are covered. The embedding correspondence is a reduction, not an
existence theorem, and claims no large-cardinal consistency result; exponential
automorphisms are untouched, and the collection's exponential faithfulness
result is outside its proof dependencies. The detector formula is not claimed to
be existential after expansion. Its Python checks are diagnostic only.

\paragraph{Source 06.}
Chen--Hamkins--Yang is prior work and the central coding idea is theirs;
Je\v{r}\'abek (2021) is prior for the conjugation obstruction; field closure and
Prikry are imported. Novelty is provisional; not peer-reviewed or Lean-certified;
no named open problem settled. The coordinate birthday is a convention, and the
automorphism conclusions do not apply to the pure surcomplex field or classify
its automorphisms. The $\beth_\omega$ theorem does not show that every singular
$\kappa$ fails Replacement, and internal ZFC does not imply external
inaccessibility. The ZF boundary concerns this construction only, and internal
Choice is not the range statement. Prikry gives relative consistency, not a ZFC
theorem, and not $H_\kappa^N=V_\kappa^N$. A nonreal parameter does not give
bi-interpretability with an unexpanded $H_\kappa$. Rigidity is not pointwise
definability. Its reduct corollary is not a minimal-language claim. It left open
noncardinal cutoffs (answered here for epsilon numbers), ill-founded models of a
surreal-arithmetic theory, quantifier elimination for an order-and-simplicity
language, and a universal embedding theorem without Global Choice. The finite
tests do not verify the theorems.

\paragraph{Source 07.}
It omitted the Chen--Hamkins--Yang credit and presented the full-class
bi-interpretation as a proposed contribution; that is corrected in
\cref{hset:thm:chy}, and it had no Je\v{r}\'abek credit, now added. Its results
were proposed contributions without priority certification, refereeing or Lean.
It claims no new proof of Kunen's inconsistency, no rigidity of the pure field,
no classification of bounded elementary self-embeddings, and no distinct
parameter-free theories from non-elementarity; nothing about unrestricted field
embeddings. Its plateau theorem is forward only, and its Choice remark says the
construction is not a reconstruction of $V$ without Choice. The Kunen transfer
needs GBC or class-parameter hypotheses and does not forbid external self-maps of
countable nonstandard models. Nonstandard models are handled internally. The
reconstruction is a logical representation, not literal membership between
chosen representatives; there is no canonical code, no decidability result, and
no simple closed sign rules for arithmetic. Rangel--Mariano and Bertram are
distinct programs whose categorical obstructions are not refuted. The finite
checks are limited.

\paragraph{Source 09.}
The global result of Chen--Hamkins--Yang, the van den Dries--Ehrlich cutoff
theorem with its erratum, and ordinary well-founded coding are prior work; the
refinements may overlap unpublished details of the Chen--Hamkins--Yang project;
no named conjecture is resolved; no Lean, no peer review. It did not credit
Je\v{r}\'abek; that credit is added. Cardinal rounding is not a rank cutoff. The
principles $\Et$ and $\Pack$ are explicit schemes, not identified with any
formulation of SA, and the independence examples are within this family only.
The detector bound is sharp only for the old-power-set parameter; no
quantifier-complexity claim is made across the translation; failure of the bound
does not imply elementarity. Extension rigidity concerns standard cutoff
structures and does not forbid abstract elementary extensions. The orientation
result is a consequence of classical rigidity, not a new automorphism
classification. Choice is used object by object and no ZF version of its
theorems was claimed (\cref{hset:thm:choice} is source 06's). Nonstandard models
are internal; no class theory, Grothendieck universe or large cardinal is needed
for the cutoff theorems. At $\omega$ the cutoff is a ring; non-epsilon cutoffs
are not classified. The singular worldly construction is a standard illustration
needing an inaccessible.

\section{Open questions, re-scoped}
\label{hset:sec:questions}

\paragraph{Answered inside this report.}
\begin{itemize}
\item Source 02's first question, a sufficient criterion for elementary
birthday inclusions: answered by source 09 (\cref{hset:thm:elementary},
\cref{hset:rem:answered-02}).
\item Source 02's caution that a two-way theorem must name the cutoff: answered
by source 09's pointed bi-interpretation (\cref{hset:rem:answered-bi}).
\item Source 05's optimal cutoffs at singular cardinals: answered by sources 06
and 07; at noncardinal epsilon numbers: answered by source 09. Source 06's and
source 07's noncardinal-cutoff questions: answered by source 09 for epsilon
numbers.
\item Source 09's fourth target, the choice-free analogue: answered for the
full-class graph coding by source 06 (\cref{hset:rem:answered-choice}).
\end{itemize}

\paragraph{Still open.}
\begin{question}[Cutoffs that are not epsilon numbers]
For which ordinals $\eta$ that are not epsilon numbers, beyond the ring cutoffs
of \cref{hset:prop:ring} (forward only), is there a uniform reconstruction from
$\No_{<\eta}$ with its inherited operations, and what exactly is recovered?
(Re-scoped from sources 05, 06 and 07.)
\end{question}

\begin{question}[Minimal signature]
How much of $\{0,1,+,-,\cdot,<,\bd\}$ can be deleted while retaining a uniform
interpretation of hereditary sets? \Cref{hset:cor:reduct} shows that $<$,
$\bd$ and the graph of one ordinal pairing polynomial suffice; whether $<$ and
$\bd$ alone, or addition with birthday, suffice is open, as is whether
multiplication is necessary. (Sources 02, 05, 07; partial progress by 06.)
\end{question}

\begin{question}[Elementary embeddings and explicit criteria]
Replace the condition $H_\kappa\prec H_\tau$ in \cref{hset:thm:elementary} by
useful explicit hypotheses in interesting families of cardinals (source 09).
Classify the elementary self-embeddings of particular bounded $\B_\kappa$; by
\cref{hset:thm:embeddings,hset:thm:two-lifts} this is the corresponding problem
for $H_\kappa$ (source 07).
\end{question}

\begin{question}[Forcing and the bounded theory]
Which changes of the complete bounded theory can occur under forcing that adds
new bounded subsets, and how much of the cutoff predicate is first-order
recoverable from an appropriately expanded $H_\kappa$ (source 02)? The detectors
of \cref{hset:sec:detectors} show failure of elementarity with parameters; they
do not decide changes of parameter-free theories.
\end{question}

\begin{question}[Quantifier complexity]
What is the least quantifier complexity of a uniform formula detecting all
proper same-ordinal outer-model inclusions in the birthday language (source 05)?
\end{question}

\begin{question}[Surcomplex reducts]
What set-theoretic information is recoverable from
$(\No_{<\kappa}[i];0,1,+,-,\cdot,\cb)$ when neither conjugation nor $i$ is
named (source 05)? \Cref{hset:thm:threshold} concerns the language with
conjugation.
\end{question}

\begin{question}[Axiomatics and formalization]
Establish definitional equivalence, or its failure, between $\Et$ and $\Pack$
and particular published axiom formulations of surreal arithmetic, once
finalized (source 09); analyse ill-founded models of such a theory, quantifier
elimination for an order-and-simplicity language, and a universal embedding
theorem without Global Choice (source 06); and formalize the size-certified
interpretation, beginning with the prefix, bit and pairing lemmas, then the
certificate verifier and its hereditary-size bound (sources 05, 06, 07, 09).
\end{question}

The collection's Lean library already contains relevant sign-sequence
infrastructure (prefix truncation, the failure of birthday precedence to be
prefix simplicity, Hessenberg agreement of ordinal arithmetic), but none of the
interpretation theorems of this report is formalized.

\section{Provenance and verification}
\label{hset:sec:audit}

\subsection{The sources' repository and literature audits}

All five sources inspected the repository read-only at the pinned commit
\begin{center}\repoSHA\end{center}
and modified nothing. Source 02 inspected the tree, the root and documentation indexes, the
catalogue source and the foundations report, and records its scope in
\path{02-birthday-cutoffs-SOURCE_AUDIT.md}, kept byte-identical. Source 05
records its requested line ranges, some truncated, in
\path{data/05-birthdays-recover-sets-repository_audit.json}. Source 06
inspected listings, the first 200 lines of the documentation catalogue, the
first 150 lines of the foundations README, and the README and first 240 lines
of the automorphisms article. Source 07 inspected the tree, READMEs, catalogue,
formalization ledger, the foundations introduction and the automorphisms README.
Source 09 inspected the catalogue, the relevant parts of the foundations report
and the scope of the automorphisms report. The corrections these audits need
are in \cref{hset:sub:repo}.

The sources' literature checks were: the Chen--Hamkins--Yang announcements and
slides (sources 02, 06, 09; source 06 cites slide pages 103, 107 and 122, and
source 02 pages 96--122, which are compatible); van den Dries--Ehrlich and its
erratum, Bournez--Guilmant, Poveda--Rinot--Sinapova, and Rangel--Mariano
(sources 02, 06, 09); Freire--Hamkins, Ehrlich's paper on the simplicity
hierarchy \cite{EhrlichTree} (record and abstract only) and Ehrlich--Kaplan
\cite{EhrlichKaplan} (background, not an input) (source 05, whose session could
not load the publisher PDF of the erratum); Bertram and Hamkins--Kirmayer--Perlmutter
(source 07); Je\v{r}\'abek (source 06). These external checks were not repeated
for this merge.

\subsection{Dependency and novelty ledger}

\begin{longtable}{@{}P{0.26\textwidth}P{0.34\textwidth}P{0.32\textwidth}@{}}
\toprule
Result & Main inputs & Status here\\
\midrule\endhead
Global birthday-field bi-interpretation (\cref{hset:thm:chy}) & Pointed well-founded relation coding & Prior Chen--Hamkins--Yang announcement; written proof from source 07; not claimed new.\\[0.4em]
Real-closed epsilon cutoffs & Classical surreal field theory & Prior van den Dries--Ehrlich, with erratum.\\[0.4em]
Small arithmetic certificates, two routes & Option graphs (09); prefix-pair tables (05, 06, 07) & Explicit local proof ingredients; no priority claim for the general locality principle.\\[0.4em]
Cardinal rounding and cutoff recovery & Bounded pairing, rooted collapse, collector, sign-graph bridge & Proposed local refinement (09), with cases in 02, 05, 06, 07.\\[0.4em]
Elementary-inclusion spectrum and embeddings & Collector and coherent bi-interpretations & Proposed classification (09); embeddings (05, 06).\\[0.4em]
Power Set, eternity, packing, worldly, $\beth_\omega$ & Hereditary cardinal bounds & Proposed calibrations (02, 06, 09); the underlying $H_\kappa$ facts are standard.\\[0.4em]
ZF range & Collapse in ZF & Direct refinement of the coding (06); no claim about other interpretations.\\[0.4em]
Preservation, first new birthday & Sign/subset bijection & Direct consequences (02, 05, 06, 07), tools rather than isolated discoveries.\\[0.4em]
Detectors, extension rigidity, $M=N$ & Old-power-set parameter; row search & Proposed explicit refinements (05, 07, 09); outer-model reasoning is standard.\\[0.4em]
Prikry invisibility & Imported Prikry forcing & Proposed application under a measurable cardinal (02, 06).\\[0.4em]
Kunen transfer & Imported Kunen inconsistency (GBC) & Proposed application (07).\\[0.4em]
Surcomplex orientation, automorphisms & Real-axis recovery, rigidity, Je\v{r}\'abek & Consequence of classical rigidity and the automorphisms report.\\[0.4em]
Two lifts, orientation threshold, full-class self-embeddings & Coordinates, rigidity, Kunen & Proved specializations (06, 07); the only surcomplex items new to the collection.\\
\bottomrule
\end{longtable}

The places where the argument could not be replaced by a general slogan are the
arithmetic certificates (saying that the global construction is definable does
not prove that the same formula computes arithmetic in an $H_\kappa$ without
Replacement), and the collector (a noncardinal cutoff is not recovered merely by
interpreting the unpointed hereditary-size universe).

\subsection{Finite checks}\label{hset:app:verification}

Each source supplied a standard-library Python program checking finite analogues
of its coding conventions. All five were rerun for this merge with Python 3.14.4
on copies of the files; each passed, and each recorded output was reproduced
exactly, apart from line endings. The recorded runs are:
\begin{center}\small
\begin{tabular}{@{}P{0.07\textwidth}P{0.85\textwidth}@{}}
\toprule
Source & Recorded finite checks\\
\midrule
02 & 65,025 prefix comparisons and 1,538 bit positions through length 7; 173,880 natural-ordinal matrix addresses on 81 Cantor-normal-form samples; 687 relabellings of all 16 sets of $V_4$; 1,024 code pairs for equality and membership; 4 invalid codes rejected.\\[0.3em]
05 & 65,025 prefix comparisons through birthday 7, including the counterexample of \cref{hset:rem:prefixes}; 64 acyclic graphs on four nodes, 4,096 graph pairs, 65,536 root-pair checks each for bisimulation equality and membership; 64 matrix round trips; 49 rectangle packings; 30 missing-row detector tests.\\[0.3em]
06 & 16,129 prefix and 642 bit instances on 127 words of length at most 6; pairing on 10,000 pairs; all 1,570 rooted relations on at most three vertices, 15 valid codes, 225 pairs each for equality and membership.\\[0.3em]
07 & 16,129 prefix comparisons, 642 sign checks, 819 radix addresses (3 radices), 42,849 graph pairs each for isomorphism and membership over 207 labeled graphs decoding 12 sets; the missing-range counterexample confirmed.\\[0.3em]
09 & 16,129 prefix pairs, 642 bit positions, 11,440 pairing checks; 530 relations and 1,570 rooted relations on at most three nodes, 15 valid codes, 225 pairs each for equality and membership; the counterexample of Appendix~\ref{hset:app:regression} confirmed.\\
\bottomrule
\end{tabular}
\end{center}
These tests can expose a reversed edge, a missing root or range condition, a bad
pairing index, or an incorrect finite sign formula. They do not prove
transfinite well-foundedness, cardinal bounds, bi-interpretability, Replacement
equivalences, forcing statements, or any large-cardinal assertion; those rest on
the written proofs and the explicitly imported classical theorems.

\section{Conclusion}

Birthday truncation has a precise set-theoretic effect. At an epsilon-number
cutoff $\lambda$, rooted sign codes recover the hereditary-size universe at the
cardinal ceiling of $\lambda$; when $\lambda$ is not itself a cardinal, its value
survives as a definable interpreted set, which explains why noncardinal cutoff
inclusions never become elementary even though all their ordered field reducts
are elementary subfields. Packing a birthday fiber requires its full power-set
cardinal to fit below the cutoff; bounding definable ordinal-indexed birthday
growth is Replacement at cardinal cutoffs; together they detect worldly, not
merely inaccessible, cardinals. Across universes, the birthday structure below a
cardinal is unchanged exactly when no bounded subset is added, while any new set
is detected by an explicit formula; yet a measurable cardinal can be singularized
without changing the structure at all. The same information is available in
surcomplex arithmetic once conjugation, coordinate birthday and an imaginary
unit are specified; omitting the unit preserves mutual interpretation but leaves
a two-element obstruction to bi-interpretation. The decisive structure is not the
field alone: it is the field with a carefully specified memory of how its
elements are born.

\appendix
\section{Formula expansion and size checks}
\label{hset:app:formulas}

\subsection{What ``first-order'' means in the coding proof}

Every abbreviation expands to a finite formula in
$\{0,1,+,-,\cdot,<,\bd\}$. Bounded ordinal quantifiers are shorthand for
ordinary scalar quantifiers with guards $\Ord(x)$ and $x<\delta$. A set variable
in an interpreted membership formula becomes a triple of scalar variables.
Equality becomes $\sim$, membership becomes $\in_\I$, and quantifiers are
relativized to $\Code$:
\[
 (u=v)^\I=(C_u\sim C_v),\quad (u\in v)^\I=(C_u\in_\I C_v),\quad
 (\exists u\,\varphi)^\I=\exists\delta_u\exists r_u\exists t_u\,
 (\Code(C_u)\land\varphi^\I),
\]
commuting with the connectives. Induction on formulas proves
$H_{\kappa(\lambda)}\models\varphi(x_{C_1},\dots)\iff\B_\lambda\models
\varphi^\I(C_1,\dots)$ for valid codes, independently of the representatives.

For example, a graph code $f$ on $\delta\times\epsilon$, embedded in the square
of width $\theta=\delta+\epsilon+1$, uses
\[
 F(i,j)\iff i<\delta\land j<\epsilon\land\Bit(f,p_\theta(i,j)).
\]
It is a total function exactly when
\[
 \forall i<\delta\ \exists j<\epsilon\,
 \bigl(F(i,j)\land\forall k<\epsilon\ (F(i,k)\longrightarrow k=j)\bigr).
\]
Injectivity is
$\forall i,i'<\delta\ \forall j<\epsilon\,(F(i,j)\land F(i',j)\to i=i')$,
and surjectivity is $\forall j<\epsilon\,\exists i<\delta\,F(i,j)$.
Occurrences of $f(i)$ in the main text can therefore be eliminated by these
unique-value clauses. Root preservation and edge correspondence add only
finitely many quantifiers. The graph code's birthday is $\Lambda(\theta)$; all
plus positions outside the relevant rectangle can be required absent, which
makes canonical relation coding a definable operation when the intended relation
is specified by a formula and its bit string exists.

The dependency order of the abbreviations is
\begin{center}\small
\begin{tabular}{@{}P{0.2\textwidth}P{0.36\textwidth}P{0.34\textwidth}@{}}
\toprule
Predicate & Earlier definitions needed & Quantification simulated\\
\midrule
$\Ord(a)$ & $\bd(a)=a$ & Embedded ordinals\\
$\Pre(x,y)$ & $<,\bd$ & Literal prefix test\\
$\Bit(s,a)$ & $\Ord,\Pre,<,\bd$ & One sign position\\
$E_C$ & $\Bit,p_\delta$ & Edges of a code\\
$\Code(C)$ & shape, clean bits, \eqref{hset:eq:extensional}--\eqref{hset:eq:rooted} & Pointed membership graphs; all subsets of one domain\\
$\sim$, $\in_\I$ & $\Code$, $E$, coded maps & Interpreted equality and membership\\
$\OC(a)$, $\Vis$ & $\Ord$, $\Bit$, $p_\theta$, $\sim$ & Literal ordinal as a set code\\
$\Coll$ & $\Code$, $\in_\I$, $\Vis$ & Collection of literal ordinals\\
\bottomrule
\end{tabular}
\end{center}
The table follows source 02's ledger, with this report's names. A tuple variable
is eliminated by replacing it with ordinary variables; ``exactly one'' has its
usual two-quantifier expansion; variables are renamed before substitution to
prevent capture. The specification is an ordinary finite formula, albeit a long
one.

\subsection{Every potentially large construction has a stated bound}

\begin{longtable}{@{}P{0.35\textwidth}P{0.56\textwidth}@{}}
\toprule
Object & Bound used in the proof\\
\midrule\endhead
Subsets tested in well-foundedness & Individual sign strings of birthday $\delta$; no collection of that whole fiber is required.\\[0.4em]
A relation on a code domain $\delta$ & Sign birthday $\Lambda(\delta)=\delta^2+\delta<\lambda$.\\[0.4em]
An isomorphism or member embedding & Square width $\delta+\epsilon+1<\lambda$, with the corresponding $\Lambda$ bound.\\[0.4em]
Collapse value of a rooted code & Transitive closure has exactly $\card\delta$ members.\\[0.4em]
An arithmetic computation & At most $\mu$ word vertices or table cells, hereditary size $\mu$, where $\mu$ bounds the inputs and is infinite.\\[0.4em]
Rank labels in that computation & Below $\mu^+$, hence each has cardinality at most $\mu$.\\[0.4em]
Replacement image after eternity & One bound $\max(\aleph_0,\card\delta,\card\varrho)<\kappa$, not an arbitrary singular-cardinal union.\\[0.4em]
A packed birthday-$\alpha$ fiber & Exists exactly if $2^{\card\alpha}<\kappa(\lambda)$.\\[0.4em]
The old-power-set parameter & Hereditary size $(2^{\card\alpha})^M$ for infinite $\alpha$.\\[0.4em]
Actual power set of a set of hereditary size $\mu$ & At most $2^\mu$; this need not be below $\kappa$, and no earlier row needs it to be.\\
\bottomrule
\end{longtable}

\section{An essential negative regression}
\label{hset:app:regression}

In the membership definition, edge preservation alone is insufficient.
Take the code for $x=\varnothing$ and the three-node chain coding
$y=\{\{\varnothing\}\}$. Map the sole node of the first code to the middle
node of the chain, which codes $\{\varnothing\}$. The image root is a child
of the target root, and the one-node induced subgraph has no edges, just
like the source. But $\varnothing\notin\{\{\varnothing\}\}$.
The predecessor-full clause \eqref{hset:eq:member-full} rejects this false
witness, because the middle node has a predecessor omitted from the image.
The finite suites of sources 06, 07 and 09 confirm this counterexample.

\section{A small concrete membership code}
\label{hset:app:example}

Let $x=\{\varnothing,\{\varnothing\}\}$, the von Neumann ordinal $2$ (source 02's
example). Its transitive closure with the root included has the three nodes
$0\leftrightarrow\varnothing$, $1\leftrightarrow\{\varnothing\}$,
$2\leftrightarrow x$. Take $\delta=3$, root $t=2$, and a relation word of length
$\Lambda(3)=12$. The edges are $(0,1)$, $(0,2)$ and $(1,2)$; under
$p_3(u,v)=3u+v$ the plus positions are $1,2,5$, and all other positions of the
length-12 word are minus. Its collapse is the identity on the ordinal $3$, with
root value $2$.

A different labeling can put $x$ at label $0$, the empty set at label $2$, and
$\{\varnothing\}$ at label $1$. The relation word then changes, and the root is no
longer the largest label. Nevertheless the isomorphism code identifies the two
roots, so interpreted equality cannot be literal equality of the triples. The
example also shows why label order must not be confused with membership order in
the countable-ordinal overflow of \cref{hset:cor:countable}.

\begin{thebibliography}{99}

\bibitem{Conway}
J. H. Conway, \emph{On Numbers and Games}, second edition, A K Peters, 2001;
first edition, Academic Press, 1976. Classical number forms and recursive
arithmetic.

\bibitem{Gonshor}
H. Gonshor, \emph{An Introduction to the Theory of Surreal Numbers},
London Mathematical Society Lecture Note Series 110, Cambridge University
Press, 1986. \url{https://doi.org/10.1017/CBO9780511629143}.

\bibitem{vdDE}
L. van den Dries and P. Ehrlich, ``Fields of surreal numbers and
exponentiation,'' \emph{Fundamenta Mathematicae} \textbf{167} (2001),
173--188. \url{https://doi.org/10.4064/fm167-2-3}. Epsilon-cutoff closure:
Proposition 4.7 and Corollary 4.9. Read with the following erratum.

\bibitem{vdDEerr}
L. van den Dries and P. Ehrlich, ``Erratum to `Fields of surreal numbers and
exponentiation','' \emph{Fundamenta Mathematicae} \textbf{168} (2001),
295--297. Bibliographic record: \url{https://eudml.org/doc/281975}.

\bibitem{BG}
O. Bournez and Q. Guilmant, ``Surreal fields stable under exponential and
logarithmic functions,'' arXiv:2201.08199, 2022.
\url{https://arxiv.org/abs/2201.08199}. Theorems 1.1--1.2 restate the corrected
ring, field and real-closedness cutoff results.

\bibitem{CHYKobe}
J. D. Hamkins, reporting joint work with J. Chen and R. Yang,
``The elementary theory of surreal arithmetic is bi-interpretable with set
theory,'' Kobe, September 2025; announcement dated August 20, 2025, with linked
slides.
\url{https://jdh.hamkins.org/theory-of-surreal-arithmetic-bi-interpretable-with-set-theory-japan-september-2025/}.
The slides announce the global result and discuss the coding and growth
principles; they are not a complete refereed paper.

\bibitem{CHYNotreDame}
J. D. Hamkins, reporting joint work with J. Chen and R. Yang,
``Surreal arithmetic,'' Notre Dame Logic Seminar; posted 12 November 2025, talk
18 November 2025, with linked slides.
\url{https://jdh.hamkins.org/surreal-arithmetic-notre-dame-logic-seminar-nov-2025/}.

\bibitem{CHYCUNY}
J. D. Hamkins, reporting joint work with J. Chen and R. Yang, ``Surreal
arithmetic is bi-interpretable with set theory,'' CUNY Logic Workshop,
13 March 2026; announcement posted 4 February 2026.
\url{https://jdh.hamkins.org/surreal-arithmetic-cuny-logic-workshop-march-2026/};
workshop page
\url{https://nylogic.github.io/logic-workshop/2026/03/13/surreal-arithmetic-is-bi-interpretable-with-set-theory.html};
slides
\url{https://jdh.hamkins.org/wp-content/uploads/Slides-Elementary-theory-of-surreal-arithmetic-Hamkins-CUNY-March-2026.pdf}.
The slides describe joint work in progress, with details still being checked.

\bibitem{Jerabek}
E. Je\v{r}\'abek, answer to O. Cunningham's question ``Are these theories of
real and complex number biinterpretable?'', MathOverflow, 13 September 2021.
\url{https://mathoverflow.net/questions/403410/are-these-theories-of-real-and-complex-number-biinterpretable}.
Original explanation of the named-imaginary-unit and automorphism-group
obstruction; not a journal publication.

\bibitem{FreireHamkins}
A. R. Freire and J. D. Hamkins, ``Bi-interpretation in weak set theories,''
\emph{The Journal of Symbolic Logic} \textbf{86} (2021), 609--634.
\url{https://doi.org/10.1017/jsl.2020.72}; arXiv:2001.05262.

\bibitem{PRS}
A. Poveda, A. Rinot and D. Sinapova, Sigma-Prikry forcing I: the axioms,
arXiv:1912.03335v2, 2020. \url{https://arxiv.org/abs/1912.03335}.
Lemma 2.10 and Section 3.1: bounded-subset preservation and ordinary
normal-measure Prikry forcing. Published version:
\url{https://doi.org/10.4153/S0008414X20000425}.

\bibitem{HKP}
J. D. Hamkins, G. Kirmayer and N. L. Perlmutter, ``Generalizations of the Kunen
inconsistency,'' arXiv:1106.1951v2, 2012. \url{https://arxiv.org/abs/1106.1951v2}.
Theorem 1 and the metamathematical discussion give the class-embedding form
used here.

\bibitem{Jech}
T. Jech, \emph{Set Theory}, third millennium edition, revised and expanded,
Springer Monographs in Mathematics, Springer, 2003.
\url{https://doi.org/10.1007/3-540-44761-X}.

\bibitem{Marker}
D. Marker, \emph{Model Theory: An Introduction}, Graduate Texts in Mathematics
217, Springer, 2002. \url{https://doi.org/10.1007/b98860}.

\bibitem{Deloro}
A. Deloro, \emph{Basic Model Theory of Algebraically Closed Fields}, mini-course,
Moscow Higher School of Economics, June 2013, Sections 3.4--3.5.
\url{https://webusers.imj-prg.fr/~adrien.deloro/teaching-archive/Moskva-ACF.pdf}.

\bibitem{RM}
D. R. Rangel and H. L. Mariano, ``An algebraic (set) theory of surreal numbers,
I,'' arXiv:1911.12726, 2019. \url{https://arxiv.org/abs/1911.12726}.
A distinct categorical foundational program, not a dependency of the proofs.

\bibitem{Bertram}
W. Bertram, ``On Conway's Numbers and Games, the Von Neumann Universe, and Pure
Set Theory,'' arXiv:2501.04412v2, 2025. \url{https://arxiv.org/abs/2501.04412v2}.

\bibitem{EhrlichTree}
P. Ehrlich, ``Conway names, the simplicity hierarchy and the surreal number
tree,'' \emph{Journal of Logic and Analysis} \textbf{3} (2011), article 1,
1--26. \url{https://doi.org/10.4115/jla.2011.3.1}.

\bibitem{EhrlichKaplan}
P. Ehrlich and E. Kaplan, ``Surreal ordered exponential fields,''
arXiv:2002.07739, 2020. \url{https://arxiv.org/abs/2002.07739}.

\end{thebibliography}
\end{document}
